%-----------------------------------------------------------------------------
\documentclass[12pt,a4paper]{report}
%---------------------------------------------------------------------------
\usepackage{amsmath}
\usepackage{amssymb}
\usepackage{amsthm}
% enumitem is loaded for the web build, not for the page layout. Without it
% tex4ht throws away the optional label on \item[(a).] and numbers the list
% 1, 2, 3 itself - so a question whose parts are (i), (ii), (iii) came out
% lettered a, b, c, and prose saying "the test in (a)" pointed at an item
% labelled "1.". With enumitem loaded the label survives verbatim, at every
% depth, and print and web finally agree.
\usepackage{enumitem}
\usepackage{amsfonts}
\usepackage{array}
\usepackage{longtable}
\usepackage{appendix}
\usepackage{blkarray}
\usepackage{bibentry}
\usepackage[none]{hyphenat}
\usepackage[noadjust]{cite}
\usepackage{filecontents}
\usepackage[all,knot,arc,poly]{xy}
\usepackage[pdftex]{graphicx}
\graphicspath{{images/}{../images/}}
\DeclareGraphicsExtensions{.pdf,.jpeg,.png, .jpg}
\usepackage{lscape}
\usepackage{epsfig}
\usepackage{mathrsfs}
\usepackage[left=40mm,right=25mm,top=25mm,bottom=25mm]{geometry}
\usepackage{url}
\usepackage{multicol}
\usepackage{multirow}
\usepackage{color}
\usepackage[table]{xcolor}
\definecolor{pink}{rgb}{1,0.5,0.5}
\usepackage{graphicx}
\graphicspath{ {images/} }
\usepackage{float}
\floatstyle{boxed} 
\restylefloat{figure}

\usepackage{tikz}
\usetikzlibrary{patterns,graphs,arrows.meta,shapes,trees,bending,fadings,quotes,angles}





\usepackage{pgfplots}
\usepgfplotslibrary{fillbetween}
\pgfplotsset{compat=1.18}

\usepackage{rotating}
%-----------------------------------------------------------------------------
\usepackage{setspace}
\usepackage{titletoc}
\titlecontents{chapter}
[0pt]
{\bfseries}
{\MakeUppercase\chaptername\ \thecontentslabel:\quad}
{\MakeUppercase}
{\hfill\contentspage}[]
%-----------------------------------------------------------------------------
\usepackage{float}
\textwidth 6.5in \textheight 9.175in\topmargin 0in \headheight 0in
\oddsidemargin 0in \evensidemargin 0in
\parskip 0.5\baselineskip
\parindent 0pt
\renewcommand{\baselinestretch}{1.25}
%-----------------------------------------------------------------------------
\theoremstyle{definition}
% One shared counter running chapter.section.n, as in the other courses here.
\newtheorem{thm}{Theorem}[section]
\newtheorem{defn}[thm]{Definition}
\newtheorem{exa}[thm]{Example}
\newtheorem{prop}[thm]{Proposition}
\newtheorem{coro}[thm]{Corollary}
\newtheorem{lem}[thm]{Lemma}
\newtheorem{result}[thm]{Result}
\newtheorem{exe}[thm]{Exercise}
\newtheorem*{note}{Note}
\newtheorem*{remark}{Remark}
\newtheorem{tech}{Technique}[section]
\newtheorem{fact}[thm]{Fact}
\newtheorem{problemx}{Problem}[section]
\newenvironment{problem}{\begin{problemx}}{\end{problemx}}
\newenvironment{solution}
  {\renewcommand{\qedsymbol}{}\begin{proof}[\textit{Solution}.]}
  {\end{proof}}
% A heading that is legal inside a theorem body, a list or a float, where
% \subsubsection is not.
\newcommand{\subhead}[1]{%
  \par\medskip\noindent\textbf{#1}\par\nobreak\smallskip}
\usepackage[nottoc]{tocbibind}
\renewcommand{\contentsname}{TABLE OF CONTENTS}
\renewcommand\bibname{REFERENCES}
%\renewcommand{\qedsymbol}{$\blacksquare$}

\parindent 0ex

%-----------------------------------------------------------------------------
\usepackage{titlesec}
\titleformat{\chapter}[hang]
          {\normalfont\filright\ttfamily\scshape\sffamily\bfseries
          	\centering\large}
          {\MakeUppercase{\chaptertitlename}\ \thechapter}{10pt\ :}{\large\MakeUppercase}        
          \titlespacing*{\chapter}{0pt}{-90pt}{16pt}
                       
\titleformat{\section}[hang]
           {\normalfont\sffamily\bfseries%
            \large}
            {\thesection.}{1ex}{}%[\titlerule]
%-----------------------------------------------------------------------------


%-----------------------------------------------------------------------------
\usepackage{fancyhdr}
\fancyhead{}

\pagestyle{fancy}

\fancyhf{}
%\fancyfoot[C]{\small Prepared by N.~Jeshurun \quad \thepage}

\fancyfoot[C]{\thepage}


%\fancyfoot[C]{\small Prepared by N.~Jeshurun}
\renewcommand{\headrulewidth}{0pt}
\renewcommand{\footrulewidth}{0pt}
%-----------------------------------------------------------------------------

\renewcommand{\sectionmark}[1]{\markright{}} %\thesection\ #1 (this one should be in th brackets, ayoub)
%-----------------------------------------------------------------------------

%-----------------------------------------------------------------------------


\newcommand{\bpf}{\noindent {\ttfamily  PROOF}. $\mbox{}$}
\newcommand{\epf}{\hfill \hbox{\raisebox{.5ex}{$\blacksquare$}$\phantom{.}$}}
%-----------------------------------------------------------------------------

\newtheorem{problemx}{Problem}[section]
% The problem environment takes no source label. Marking a question
% "Test 2" or "2013 examination" turns it into something to be memorised
% rather than a topic to be understood, and it ties the question to one
% course's assessment when the same material is examined under several
% codes. Problems sit under the topic they belong to; that is the only
% label a reader needs.
\newenvironment{problem}{\begin{problemx}}{\end{problemx}}
\newenvironment{solution}
  {\renewcommand{\qedsymbol}{}\begin{proof}[\textit{Solution}.]}
  {\end{proof}}

% tex4ht expands amssymb's \checkmark in TEXT mode and hands MathJax the raw
% internals, so the reader sees \m@th \mathchar"458 instead of a tick.
% \ensuremath forces maths mode, where the token survives.
\let\amscheckmark\checkmark
\renewcommand{\checkmark}{\ensuremath{\amscheckmark}}

\begin{document}
\begin{titlepage}
\begin{figure}[H]
\centering
\includegraphics[scale=0.5]{unza.png}
\end{figure}
\vspace{0.6cm}
\centering
{\LARGE THE UNIVERSITY OF ZAMBIA}\\
\vspace{0.3cm}
{\LARGE School of Natural  Sciences}\\
\vspace{0.3cm}
{\LARGE Department of Mathematics $\&$  Statistics}\\
\vspace{3cm}
{\LARGE\sffamily MAT 5611: STATISTICAL INFERENCE}\\
{\LARGE\sffamily }\par
\vfill
{\Large Lecture Notes: 2018 Academic Year}\par
\vspace{4cm}
{\Large\sffamily Lecturer: Dr. V. NAWA}
\vfill
{\large\bfseries by: NEPHAT JESHURUN MWANZA}\par
\vspace*{0.1cm}
\end{titlepage}









\tableofcontents
\thispagestyle{empty}
\clearpage
\setcounter{page}{1}

%___________________________________________________________--------------------------------------------------------------------------


















\newpage
\chapter{PROPERTIES OF STATISTICS}
\vspace{0.3cm}



\section{Introduction}
Assume that we have designed and collected data $\, X_1\, , \, X_2\, , \, \cdots\cdots\cdots\, , \, X_n\,,\,$ where $n$ is the sample size, is fixed and known.\\
Statistical inference deals with methods of using the outcomes to obtain information on the ``true distribution function" (or ``true parameter") underlying the experiment.\\
The two main topics in statistical inference are estimation and hypothesis testing.
\begin{enumerate}
\item[] In estimation, the observations are used to find an approximate value (estimate) for some numerical characteristic of the true distribution function. This can take the form of one numerical value (point estimation) or a  set of possible values (interval estimation).
\item[] In hypothesis testing, the observations are used to conclude whether or not the true distribution or true parameter belongs to a smaller subset of distribution functions or parameter space respectively.
\end{enumerate}


\begin{notation}
We will sometimes denote the random sample $X_1\, , \, X_2, \, \cdots\,, \, X_n$ drawn from a probability (density) function (p.d.f) $f_{\theta}(x)$ compactly by a random vector $ X = (X_1\, ,\, X_2, \, \cdots\, X_n)$ i.e $$f_{\theta}(x) = f_{\theta}(x_1\, , \, x_2\, , \, \cdots\, , \, x_n)$$
\end{notation}


\textbf{Example}\\
Suppose $\, X = (X_1\, ,\, X_2 , \, \cdots, \, X_n)\,$ is a random sample from a $N(\mu\, ,\, \sigma^2)$ distribution. Then 
$\theta = (\mu\, , \, \sigma^2)$
$$f_{\theta}(x) = \prod^n_{i = 1}\, \dfrac{1}{\sqrt{2\pi \sigma^2}}e^{-\dfrac{1}{2\sigma^2}(x_i -\mu)^2}$$
and $\Omega = \big\{(\mu\, , \, \sigma^2):\, -\infty < \mu < \infty\, , \, \sigma^2 >0\big\}\,$ is the parameter space.\\


\textbf{Note:}
Models are written in the form $\big\{f_{\theta}(x)\, :\, \theta\in \Omega\big\}\,$ where $\Omega$ is the parameter space and $f_{\theta}(x)$ is the p.d.f.


\begin{defn}
A statistic, $T(X)$, is a function of the data which does not depend on the unknown parameter $\theta$.
\end{defn}


e.g $\hspace{0.5cm}\, \, \overline{X} = {\frac{1}{n}\sum^n_{i = 1} x_i } = $ sample mean, \, $S^2 = {\frac{1}{n - 1}\sum^n_{i=1}(X_i - \overline{X})^2} = $ Sample variance,\\ 
$X_{(1)}  = \min(X_1\, ,\, X_2  , \, \cdots , \, X_n)\,$,\, $X_{(n)} = \max(X_1\, ,\, X_2 , \,\cdots , \, X_n)\,$,\\
$(X_{(1)}\, ,\, X_{(2)}, \, \cdots, \, X_{(n)}) = $ order statistics.\\

$\dfrac{\overline{X} - \mu}{\sigma /\sqrt{n}}\,$ is not a statistic if $\mu$ and $\sigma^2$ are unknown.\\




\textbf{Note:}
\begin{enumerate}
\item Although a statistic $T(X)$ , is not a function of $\theta$, its distribution can depend on $\theta$.
\item An estimator is a statistic used to estimate an unknown parameter $\theta$.
\end{enumerate}
i.e $\,$ An estimator is a random variable. e.g $\,\overline{X}$\\

An estimate is an observed value of a random variable e.g $\,\overline{x}$.






\section{Unbiasedness and Mean Square Error}
What desirable properties should our estimator have?
\begin{enumerate}
\item Is it estimating the correct parameter?
\item Is it not consistently too large or too small?
\end{enumerate}




One way of addressing (1) is to study unbiasedness.



\begin{notation}
We will denote an expected value under the assumed parameter value $\theta$ by $E_{\theta}(\cdot)$. Thus
\[E_{\theta}\left[h(x)\right] = 
\begin{cases}
\displaystyle{\int_{\text{all} \, x} h(x)\, f_{\theta}(x)\, dx} & \text{if}\hspace{0.2cm} X\hspace{0.2cm}\text{is continuous}\\\\
\displaystyle{\sum_{\text{all}\, x} h(x)\, f_{\theta}(x)} & \text{if}\hspace{0.2cm} X\hspace{0.2cm} \text{is discrete}\\
\end{cases}
\]
\end{notation}





\begin{exa}
Suppose that $X$ has Cauchy CAU(1,0) distribution show that $E_{\theta}(X)$ does not exists.
\end{exa}


\begin{proof}[\textbf{Solution.}]
$\displaystyle{f_{\theta}(x) = \dfrac{1}{\pi \left(1 + x^2\right)}}\hspace{0.2cm},\hspace{0.5cm} -\infty < x < \infty $
\begin{align*}
E_{\theta}\,|x| &  = \int^{\infty}_{-\infty}\frac{\mid x\mid}{\pi \left(1 + x^2\right)}\, dx \\\\
& = \dfrac{2}{\pi}\int^{\infty}_{0} \frac{x}{1 + x^2}\, dx\\\\
& = \lim_{m \rightarrow \infty} \, \frac{2}{\pi}\int^{\infty}_0\frac{x}{1 + x^2}\, dx\\\\
& = \lim_{m\rightarrow \infty} \frac{1}{\pi}\, \log\left(1 + x^2\right)\Bigg|^m_0\\\\
& = \infty
\end{align*}
\end{proof}









\begin{exa}
Suppose that $X$ is a random variable with p.d.f
$$f_{\theta}(x) = \dfrac{\theta}{x^{\theta + 1}}\hspace{0.3cm}, \hspace{0.5cm} x\geq 1$$
For what values of $\theta$ do $E_{\theta}(X)$ and  $var_{\theta}(X)$ exists?
\end{exa}


\begin{proof}[\textbf{Solution.}]
\begin{align*}
E_{\theta}(X) & = \int_{\text{all}\, x} x\, f_{\theta}(x)\, dx = \int_1^{\infty} x\, \frac{\theta}{x^{\theta + 1}}\, dx
 = \int_1^{\infty} \theta\, x^{-\theta}\, d\theta\\\\
& = \dfrac{\theta}{1 - \theta} \, x^{1-\theta}\Bigg|^{\infty}_1\hspace{0.3cm}, \hspace{0.5cm} \theta \neq 1\\\\
& = \frac{\theta}{1 - \theta}\, (0 - 1)\hspace{0.3cm}, \hspace{0.5cm} 1 - \theta < 0\\\\
& = \frac{\theta}{1 - \theta}\hspace{0.3cm}, \hspace{0.5cm} \theta > 1
\end{align*}

$$var_{\theta}(X) = E_{\theta}\left(X^2\right) - \left[E_{\theta}(X)\right]^2$$

\begin{align*}
E_{\theta}\left(X^2\right) & = \int_1^{\infty}x^2 \, \frac{\theta}{x^{\theta + 1}}\, dx
 = \int_1^{\infty} \theta\, x^{1-\theta}\, dx\\\\
& = \dfrac{\theta}{2 - \theta}\, x^{2 - \theta}\Bigg|^{\infty}_{1}\hspace{0.3cm},\hspace{0.5cm} \theta \neq 2\\\\
& = \frac{\theta}{2 - \theta}\, (0  - 1)\hspace{0.3cm}, \hspace{0.5cm} 2-\theta < 0\\\\
& = \frac{\theta}{\theta - 2}\hspace{0.3cm} , \hspace{0.5cm} \theta > 2
\end{align*}


\begin{align*}
var_{\theta}(X) & = E_{\theta}\left(X^2\right) - \left[E_{\theta}(X)\right]^2\\\\
& = \dfrac{\theta}{\theta - 2} - \left(\frac{\theta}{\theta - 1}\right)^2  = \frac{\theta}{\theta - 2} - \frac{\theta^2}{\left(\theta - 1\right)^2}\hspace{0.5cm}\theta\neq 1\, , \, \theta\neq 2\\\\
& = \frac{\theta(\theta - 1)^2 - \theta^2(\theta - 2)}{(\theta - 1)^2(\theta - 2)}\\\\
& = \frac{\theta\left[\theta^2 - 2\theta + 1 - \theta^2 + 2\theta\right]}{(\theta - 1)^2(\theta - 2)}\\\\
& = \frac{\theta}{(\theta - 1)^2(\theta - 2)}\\\\
\therefore\,  var_{\theta}(X) & = \frac{\theta}{(\theta - 1)^2(\theta - 2)}\hspace{0.3cm}, \hspace{0.5cm} 2<\theta
\end{align*}
\end{proof}



\begin{problem}
If $X\thicksim$ GAM$(\alpha\, , \, \beta)$, show that 
$$E_{\theta}\left(X^p\right) = \frac{\beta^p\, \Gamma\left(\alpha + p\right)}{\Gamma (\alpha)}$$
For what values of $p$ does this expectation exist?
\end{problem}


\begin{problem}
Let $\, m(t) = E\left(e^{tx}\right)\,$ be the moment generating function of a continued random variable $X$. The function $m(-t)$ is often called the Laplace Transform of the probability function $X$. Show that
$$E\left(X^{-p}\right) = \frac{1}{\Gamma (p)}\, \int_0^{\infty}\, m(-t)\, t^{p - 1}\, dt$$
\end{problem}



\begin{defn}
A statistic $T(X)$ is an unbiased estimator of $\theta$ if $E_{\theta}\left(T(X)\right) = \theta\,$ for all $\theta\in \Omega$.
\end{defn}



\begin{exa}
Suppose $X_i\thicksim$ POI$(i\theta)\, , \hspace{0.3cm} i = 1, \, 2, \, \cdots\, n\,$ independently. Determine whether the estimators $\, {T_1 = \frac{1}{n} \sum_{i = 1}^{n} \frac{X_i}{i}}\hspace{0.3cm}$ and $\, {T_2 = \frac{2}{n(n+1)}\sum^n_{i=1}X_i}\,$ are unbiased estimators of $\theta$.
\end{exa}




\begin{proof}[\textbf{Solution.}]
$$
E\left(T_1\right)  \, =\, E\left[\frac{1}{n}\sum_{i = 1}^n\frac{X_i}{i}\right] = \frac{1}{n}\sum_{i = 1}^n\frac{E\left(X_i\right)}{i}\,
 = \,\frac{1}{n}\sum^n_{i = 1}\frac{i\theta}{i}\,
 =\, \frac{1}{n}\, n\theta\, 
 = \,\theta.$$
 
$\therefore\,\, T_1\,$ is an unbiased estimator of $\theta$.

\begin{align*}
E\left(T_2\right) & = E\left( \frac{2}{n(n+1)} \, \sum^n_{i=1} X_i\right) = \frac{2}{n(n+1)}\sum^n_{i = 1} E\left(X_i\right)\\
& = \frac{2}{n(n+1)}\sum_{i = 1}^ni\theta\\
& = \frac{2\theta}{n(n+1)}\sum^n_{i=1}i = \frac{2\theta}{n(n+1)}\cdot \frac{n(n+1)}{2}\\
& = \theta
\end{align*}

$\therefore\,\, T_2\,$ is an unbiased estimator of $\theta$.
\end{proof}






\textbf{Note:}\,
Unbiased estimator is generally not preserved under transformations. For example, if $T$ is an unbiased estimator of $\theta$, then $T^2$ may be not an unbiased estimator of $\theta^2$.






\begin{exa}
Suppose $X_1, \, X_2 \cdots \, X_n$ are uncorrelated random variables with $E(X_i) = \mu$ and $var(X_i) = \sigma^2\, , \hspace{0.3cm} i = 1\, , \, 2\, , \, \cdots\cdots\, , \, n$.
\begin{enumerate}
\item[(i).] Show that $\, \displaystyle{T = \frac{1}{n} \sum^n_{i = 1}\left( X_i - \mu\right)^2}\,$ is an unbiased estimator for $\sigma^2$ if $\mu$ is known.
\item[(ii).] Show that $\,\displaystyle{S^2 = \frac{1}{n - 1}\sum_{i = 1}^n\left(X_i - \mu\right)^2}\,$ is an unbiased estimator of $\sigma^2$ if $\mu$ is unknown.
\item[(iii).] If $\, X_1\, , \, X_2\, \cdots\cdots\, X_n\,$ is a random sample from the $N(\mu\, , \, \sigma^2)$ distribution is unbiased estimator of $\sigma$.
\end{enumerate}
\end{exa}



\begin{proof}[\textbf{Solution.}]
\begin{enumerate}
\item[(i).] 
$\begin{aligned}[t]
E\left(T\right)&  = E\left[\frac{1}{n}\sum^n_{i=1}\left(X_i - \mu\right)^2\right]  = \frac{1}{n}\sum^n_{i=1}E\left(X_i -\mu\right)^2
 = \frac{1}{n}\sum_{i=1}^n\sigma^2
 = \frac{1}{n}\, n\sigma^2
 = \sigma^2.
\end{aligned}$


\item[(ii).]
$\begin{aligned}[t]
E\left(S^2\right) & = E\left[\frac{1}{n - 1}\, \sum_{i = 1}^n \left(X_i - \overline{X}\right)^2\right]
= E\left[\frac{1}{n -1}\sum_{i = 1}^n\left(X_i - \mu - (\overline{X} - \mu)\right)^2\right]\\\\
& = E\left[\frac{1}{n -1}\sum_{i = 1}^n\left((X_i - \mu)^2 + (\overline{X}- \mu)^2 - 2(X_i -\mu)(\overline{X} - \mu\right)\right]\\\\
& = \frac{1}{n -1} E\left[\sum_i(X_i -\mu)^2  + n(\overline{X} - \mu)^2 - 2n(\overline{X} - \mu)^2\right]\\\\
& = \frac{1}{n - 1}\left[\sum_i E(X_i - \mu)^2 - n\underbrace{E(\overline{X} -\mu)^2}_{var(\overline{X})}\right]\\\\
& = \frac{1}{n - 1}\left[\sum_i \sigma^2 - n\, var(\overline{X})\right]\\\\
& = \frac{1}{n - 1} \left[n\, \sigma^2 - n\, \frac{\sigma^2}{n}\right]
\end{aligned}$

\begin{align*}
\implies\hspace{0.3cm} E\left(S^2\right) = \frac{1}{n - 1}\left[n\sigma^2 - \sigma^2\right]
 = \frac{1}{n - 1}\, (n - 1)\sigma^2
 = \sigma^2.\\
\end{align*}


\item[(iii)] $\displaystyle{\frac{(n - 1)S^2}{\sigma^2}\, \thicksim\, \chi^2_{n=1}\, =, GAM\left(\frac{n-1}{2}\, , \, 2\right)}\, $ by 1.2.4
$$\,\,E\left(X^p\right) = \frac{\beta^p\,\Gamma\left(\alpha + p\right)}{\Gamma(\alpha)}$$
with $p = \dfrac{1}{2}\hspace{0.3cm},\hspace{0.3cm} \alpha = \frac{n - 1}{2}\hspace{0.3cm},\hspace{0.3cm} \beta = 2$. We have
\begin{align*}
E\left[\left(\frac{(n - 1)S^2}{\sigma^2}\right)^{1/2}\right] & = \frac{2^{1/2}\, \Gamma\left(\frac{n - 1}{2} + \frac{1}{2}\right)}{\Gamma\left(\dfrac{n - 1}{2}\right)}\\\\
E\left[\frac{\sqrt{n-1}\, S}{\sigma}\right] &  = \frac{\sqrt{2}\,\, \Gamma\left(\dfrac{n}{2}\right)}{\Gamma\left(\dfrac{n - 1}{2}\right)}\\\\
\frac{\sqrt{n - 1}}{\sigma}\, E(S) & = \frac{\sqrt{2}\hspace{0.2cm} \Gamma\left(\dfrac{n}{2}\right)}{\Gamma\left(\dfrac{n - 1}{2}\right)}
\end{align*}

\begin{align*}
\implies\hspace{0.5cm} \frac{\Gamma\left(\dfrac{n - 1}{2}\right)\, \sqrt{n - 1}}{\sqrt{2}\,\, \Gamma\left(\dfrac{n}{2}\right)}\, E(S) & = \sigma\\
\implies\hspace{0.5cm} E\left[\frac{\sqrt{n-1}\,\, \Gamma\left(\dfrac{n - 1}{2}\right)}{\sqrt{2}\,\, \Gamma\left(\dfrac{n}{2}\right)}\right] & = \sigma 
\end{align*}
$\displaystyle{T = \dfrac{\Gamma\left(\dfrac{n-1}{2}\right)}{\Gamma\left(\dfrac{n}{2}\right)}\, \sqrt{\dfrac{n -1 }{2}}}\, S\,$ is an unbiased estimator of $\sigma$.
\end{enumerate}
\end{proof}






\begin{exa}
Suppose $X\thicksim$ BIN$(n,\theta)$
\begin{enumerate}
\item[(i).] Find an unbiased estimator $T(X)$ of $\theta$.
\item[(ii).] Is $\,\dfrac{1}{T(X)}\,$ an unbiased estimator of $\, \dfrac{1}{\theta}\,$?
\end{enumerate}
\end{exa}




\begin{proof}[\textbf{Solution.}]
\begin{enumerate}
\item[(i).] $T(X) = \dfrac{X}{n}\, $ is an unbiased estimator of $\theta$. i.e 
$$E(T) = E\left(\frac{X}{n}\right) = \frac{E(X)}{n} = \frac{n\, \theta}{n} = \theta$$

\item[(ii).]
$\begin{aligned}[t]
E\left(\frac{1}{T}\right) & = E\left(\frac{n}{X}\right)
 = \sum_{\text{all}\, x} \frac{n}{x}\, \binom{n}{x}\, \theta^x\, (1-\theta)^{n-x}
 = \sum^n_{x = 0} \frac{n}{x}\, \binom{n}{x}\, \theta^x\, (1 - \theta)^{n-x}
\end{aligned}$\\
which doesn't exist (we are starting from a zero 0).
\end{enumerate}
\end{proof}





\begin{problem}
Let $X_1\, , \, X_2\, ,\, \cdots\cdots\, X_n\,$ be a random sample from the POI$(\theta)$ distribution. Find
$$E_{\theta}\left(X^{(k)}\right) = E_{\theta}\left[X(X-1)(X-2)\cdots\cdots(X-k+1)\right]$$
the $k^{\text{th}}$ factorial moment of $X$, and thus find an unbiased estimator of $\theta^k\,\, , \,\ k = 1\,, \, 2\, \, \cdots\cdots$
\end{problem}




Suppose $T_1(X)$ and $T_2(X)$ are estimators of $\theta$.


% Diagram  1
\begin{table}[H]
\centering
\begin{tabular}{ccccccc}
\begin{tikzpicture}[=>Stealth]
\draw[smooth,domain = -2.5:2.5, variable = \x]plot(\x, {6*1/sqrt(2*pi)*exp(-1/2*\x*\x)});
\draw[-Stealth](-2.8,0)--(3,0)node[below]{$t_1$};
\draw[-Stealth](-2.8,0)--(-2.8,3.5);
\draw[dashed](0,0)--(0,2.4);
\draw(0,0)node[below]{$\theta$}
(0,2.4)node[above]{$f_{\theta}(t_1) = $ p.f.d of $T_1$}
(0,-1)node[below]{$E(T_1) = \theta$};
\end{tikzpicture}
&&&&&&
\begin{tikzpicture}
\draw[smooth,domain = -2.5:2.5, variable = \x] plot(\x, {6*1/sqrt(pi)*exp(-1/2*\x*\x)});
\draw[-Stealth](-2.8,0)--(3,0)node[below]{$t_2$};
\draw[-Stealth](-2.8,0)--(-2.8,3);
\draw[dashed](0,0)--(0,3.4);
\draw(0,0)node[below]{$\theta + \frac{1}{n}$}
(0,3.4)node[above]{$g_{\theta}(t_2) = $ p.d.f of $T_2$}
(0,-1)node[below]{$E(T_2) = \theta + \dfrac{1}{n}$};
\end{tikzpicture}\\
\end{tabular}
\end{table}







$var_{\theta}(T_1)$  is large $var_{\theta}(T_2)$ is small which estimator is better?\\

To determine whether a given estimator $T(X)$ does well for estimating $\theta$, we consider a loss function or distance between the estimator and the true value denoted by $L\left(\theta\, , \, T(X)\right)$. This loss function is averaged over all possible values of the data to obtain the risk
$$Risk = E_{\theta}\big\{L(\theta\, , \, T(X)\big\}$$
A good estimator is one with small risk. One particular loss function is 
$$L\left(\theta\, , \, T(X)\right) = \left(T(X) - \theta\right)^2$$
which is called the squared error loss function whose corresponding risk, called mean squared error (MSE) is given by 
$$MSE_{\theta}(T) = E_{\theta}\left[(T(X) - \theta)^2\right]$$
Another loss function is $$L\left(\theta\, , \, T(X)\right) = \begin{vmatrix}
T(X) - \theta\\
\end{vmatrix}$$
which is called the absolute error loss function whose corresponding risk, called the mean absolute error, is given by 
$$Risk = R\left(\theta\, , \,T\right) = E_{\theta}\begin{vmatrix}
T(X) - \theta\\
\end{vmatrix}$$







\begin{problem}
Show that $\,\, MSE_{\theta} (T) = var_{\theta}(T) + \left(Bias_{\theta}(T)\right)^2\,$ where $Bias_{\theta}(T) = E_{\theta}\left(T(X) - \theta\right)$.
\end{problem}




\begin{exa}
Let $\, X_1\, , \,X_2\, , \, \cdots\cdots\, X_n\,$ be a random sample from a UNIF$(0,\theta)$ distribution. Compute the MSE's of the following three estimators of $\theta$:
$$T_1 = 2\overline{X}\hspace{0.4cm},\hspace{0.4cm} T_2  = X_{(n)}\hspace{0.4cm},\hspace{0.4cm} T_3 = (n + 1)X_{(1)}$$
where $X_{(n)} = \max(X_1\, , \, X_2\, \cdots\cdots\, X_n)$ and $X_{(1)} = \min(X_1\, , \, X_2\, \cdots\cdots\, X_n)$
\end{exa}




\begin{proof}[\textbf{Solution.}]
$MSE(T_1) = var(T_1) + (bias(T_1))^2$
\begin{align*}
bias (T_1) & = E(T_1) - \theta\\
& = E\left(2\overline{X}\right) - \theta = 2\, E(\overline{X}) - \theta\\
& = 2E(X) - \theta\\
& = 2\left(\frac{\theta}{2}\right) - \theta = 0.
\end{align*}

\begin{align*}
MSE(T_1) & = var(T_1) = var(2\overline{X})\\
& = 4\, var(\overline{X}) = \frac{4}{n}\, var(X)\\
& = \frac{4}{n}\, \frac{\theta^2}{12}\\
& = \frac{\theta^2}{3n}.
\end{align*}


$MSE(T_2) = var(X_{(n)}) + (bias(X_{(n)})^2$\\

$Bias(X_{(n)}) = E(X_{(n)}) - \theta$\\


Let $Y = X_{(n)}$
\begin{align*}
F(y) & = P(Y\leq y) = P(X_{(n)} \leq y)\\
& = P(X_1\leq y\, , \, X_2\leq y\, , \,\cdots\cdots\, ,\, X_n \leq y)\\
& = \left[F_X(y)\right]^n
\end{align*}
$$f(y) = n\, \left[F_X(y)\right]^{n-1}\, f_X(y)$$

\begin{align*}
\text{So}\hspace{0.8cm}F_X(y)  & = P(X\leq x) = \int_0^x\dfrac{1}{\theta}\, \, dt
 = \dfrac{t}{\theta}\Bigg|_0^x
 = \dfrac{x}{\theta}
\end{align*}


$$f(y) = n\left(\frac{y}{\theta}\right)^{n-1}\, \dfrac{1}{\theta} = \frac{n\, y^{n-1}}{\theta^n}\hspace{0.5cm}, \hspace{0.3cm} 0 < y < \theta$$



\begin{align*}
E\left(X_{(n)}\right) & = E(Y)
 = \int_0^{\theta} \frac{y\, n\, y^{n-1}}{\theta^n}\, dy
 = \frac{n\theta}{n + 1}.
\end{align*}




$$Bias(T_2) = \frac{n\theta}{n+ 1} - \theta = -\frac{\theta}{n + 1}$$



$$var(T_2) = E\left(T^2_2\right) - \left(E(T_2)\right)^2$$




\begin{align*}
E\left(T_2^2\right)  = E(Y^2)
 = \int^{\theta}_0y^2\, \frac{n y^{n-1}}{\theta^n}\, dy
 = \frac{n\, \theta^2}{n + 2}.
\end{align*}


$$var(T_2) = \frac{n\, \theta^2}{n+2} - \left(\frac{n\, \theta}{n + 1}\right)^2 = \frac{n\, \theta^2}{(n + 1)^2\, (n + 2)}.$$




\begin{align*}
MSE(T_2) & = \frac{n\, \theta^2}{(n + 1)^2\, (n + 2)} + \frac{\theta^2}{(n + 1)^2}
 = \frac{2\,\theta^2}{(n + 1)\, (n+ 2)}.
\end{align*}






$$MSE(T_3) = var(T_3) + \left(bias(T_3)\right)^2$$


\begin{align*}
bias(T_3) & = E(T_3) - \theta\,
 = \,E\left((n + 1)\,X_{(1)}\right) - \theta\,
 = \, (n+1)\, E(X_{(1)}) - \theta.
\end{align*}




Let $Y = X_{(1)}$, then
\begin{align*}
F(y) & = P\left(X_{(1)}\leq y\right)\\\\
& = 1 - P\left(X_{(1)} > y\right)\\\\
& = 1 - P\left(X_1 > y\, , \, X_2>y\, , \, \cdots\cdots\, , \, X_n>y\right)\\\\
& = 1 - \left(1 - F_X(y)\right)^n
\end{align*}



\begin{align*}
f(y) & = n\, \left(1 - F_X(y)\right)^n\, f_X(y)\\\\
& = n \, \left(1 - \frac{y}{\theta}\right)^{n-1}\, \frac{1}{\theta}\\\\
& = \frac{n}{\theta}\left(1 - \frac{y}{\theta}\right)^{n-1}\hspace{0.5cm},\hspace{0.3cm} 0 < y < \theta
\end{align*}



\begin{align*}
E\left(X_{(1)}\right) & = \int_0^{\theta}y\, \frac{n}{\theta}\, \left(1 - \frac{y}{\theta}\right)^{n-1}\, dy\\\\
&\hspace{0.6cm}\text{Let}\hspace{0.3cm} u = 1 - \frac{y}{\theta}\\\\
& = \frac{\theta}{n + 1}.
\end{align*}







\begin{align*}
E(T_3) & = (n + 1)\, E\left(X_{(1)}\right)\,
 = \, (n + 1)\, \frac{\theta}{n + 1}
 = \theta.
\end{align*}



\begin{align*}
MSE(T_3) & = var(T_3)\\\\
& = var\left[(n+1)\, \left(X_{(1)}\right)\right]\\\\
& = (n+1)^2\, var\left(X_{(1)}\right)
\end{align*}



$$var\left(X_{(1)}\right) = E\left(X_{(1)}^2\right) - \left(E(X_{(1)})\right)^2$$



\begin{align*}
E\left(X_{(1)}^2\right) & = E\left(Y^2\right)\,
 =\, \int_0^{\theta} y^2\, \frac{n}{\theta}\, \left(1 - \frac{y}{\theta}\right)^{n-1}\, dy\,
 = \,\frac{2\, \theta^2}{(n + 1)\, (n + 2)}.
\end{align*}



$$var\left(X_{(1)}\right) = \frac{2\, \theta^2}{(n + 1)\, (n + 2)} - \frac{\theta^2}{(n + 1)^2}\, = \, \frac{n\, \theta^2}{(n + 1)^2 \, (n + 2)}$$


$$MST(T_3) = (n + 1)^2\, \frac{n\, \theta^2}{(n + 1)^2\, (n + 2)} = \frac{n\, \theta^2}{n+ 2}$$


i.e $\, \, T_2$ is the best estimator.
\end{proof}








\begin{problem}
Let $\, X_1\, , \, X_2\, , \, \cdots\cdots\, , \, X_n\,$ be a random sample from a UNIF$(0\, , \, \theta)$ distribution. Consider the following estimators of $\theta$:
\begin{align*}
T & = X_{(1)}  + X_{(n)}\\
R & = X_{(n)} - X_{(1)}
\end{align*}
which estimator is unbiased?
 Compare the MSE's of these estimators.
\end{problem}





\begin{problem}
Let $\, X_1\, , \, X_2\, , \, \cdots\cdots\, , \, X_n\,$ be a random sample from the $N(\mu\, , \, \sigma^2)$ distribution. Consider the following estimators of $\sigma^2$:
$$T_1= \frac{n - 1}{n}S^2\hspace{0.8cm}, \hspace{0.8cm} T_2 = \frac{n - 1}{n + 1}\, S^2\hspace{0.8cm}\text{where}\hspace{0.5cm}S^2 = \frac{1}{n - 1}\sum^n_{i = 1}\left(X_i - \overline{X}\right)^2$$
Compare the MSE's of these estimators by grouping them as a function of $\sigma^2$ for $n = 5$.
\end{problem}



\begin{exa}
Let $\, X\thicksim N(\theta\, ,\, 1)$. Compare the MSE's of the three estimators
$$T_1 = X\hspace{0.8cm}, \hspace{0.8cm} T_2 = \frac{X}{2}\hspace{0.8cm},\hspace{0.8cm} T_3 = 0$$
\end{exa}






\begin{proof}[\textbf{Solution.}]
$MSE(T_1) = var(T_1) + \left(bias(T_1)\right)^2$\\

$E(T_1) = E(X) = \theta\hspace{0.8cm}$ unbiased.


$$MSE(T_1) = var(T_1) = var(X) = 1.$$


$MSE(T_2) = var(T_2) + \left(bias(T_2)\right)^2$\\

$E(T_2) = \frac{\theta}{2}\,$ and $\, var(T_2) = \frac{1}{4}$

\begin{align*}
bias(T_2) & = E(T_2) - \theta = \frac{\theta}{2} - \theta\,= -\,\frac{\theta}{2}
\end{align*}

$$MSE(T_2) = \dfrac{1}{4} + \dfrac{\theta^2}{4}$$






$MSE(T_3) = var(T_3) + \left(bias(T_3)\right)^2$\\
We have that $\,E(T_3) = 0$\,\, and \, $var(T_3) = 0$ which implies that \,$bias(T_3) = E(T_3) - \theta = -\theta$
$$MSE(T_3) = \theta^2$$



% Diagram 2
\begin{figure}[H]
\centering
\begin{tikzpicture}[=>Stealth]
\draw[-Stealth](-3,0)--(3.5,0)node[below]{$\theta$};
\draw[-Stealth](0,-0.5)--(0,3.5);
\draw(-2.5,1)--(3,1)
(0,0)parabola(2,3)
(0,0)parabola(-2,3)
(0,0.5)parabola(2.5,2)
(0,0.5)parabola(-2.5,2);
\draw(0,3)node[]{$-$}
(0,1.2)node[left,scale=0.9]{1}
(-0.2,3)node[left,scale=0.9]{4}
(2.5,-0.08)--(2.5,0.08)
(-2.5,-0.08)--(-2.5,0.08)
(-2.5,-0.08)node[below,scale=0.9]{-2}
(2.5,-0.08)node[below,scale=0.9]{2};
\draw(3,0.8)node[above]{$MSE(T_1)$}
(2.4,2)node[right]{$MSE(T_2)$}
(2,3)node[above,right]{$MSE(T_3)$};
\end{tikzpicture}
\end{figure}


\vspace{0.3cm}
i.e $\, $ No estimator outperforms all other estimators even the trivial estimator $T_3$ is better for some values of $\theta$.\\
In order to achieve a best estimator, we need to restrict ourselves to a specific class of estimators and select the best within the class. One possible restriction is to the class of unbiased estimators.
\end{proof}






\begin{defn}
An estimator $T(X)$ is said to be uniformly minimum variance unbiased estimator $(UMVUE)$ of the parameter $\theta$ if 
\begin{enumerate}
\item[i.] it is an unbiased estimator of $\theta$
\item[ii.] among all unbiased estimators of $\theta$ it has the smallest $MSE$ and therefore the smallest variance.
\end{enumerate}
\end{defn}





\section{Sufficiency}
A sufficient statistic for a parameter $\theta$ is a statistic that, in a certain sense, captures all the information about $\theta$ contained in the sample. Any additional information in the sample besides the value of the sufficient statistic does not contain any more information about $\theta$.



\begin{sufp}
If $T(X)$ is a sufficient statistic  for the model $\, \{f_{\theta}(x)\, :\, \theta\in \Omega\}\,$ then any inference about $\theta$ should depend on the sample only through the value of $T(X)$ i.e if $x$ and $y$ are two sample points such that $T(x) = T(y)$, then the inference about $\theta$ should be the same whether $X = x\,$ or $\, Y = y\,$ is used.
\end{sufp}


\begin{defn}
A statistic is sufficient for the model $\, \{f_{\theta}(x)\, :\, \theta\in \Omega\}\,$  if the distribution of the data $\, X_1\, , \, X_2\, , \, \cdots\cdots\, , X_n\,$ given $T = t$ does not depend on the unknown parameter $\theta$.
\end{defn}




\begin{exa}
Let $\, X_1\, , \, X_2\, , \, \cdots\cdots\, , \, X_n\,$  be a random sample from the $POI(\theta)$ distribution. Show that $\, T = \sum^n_i X_i\, $is sufficient for this model.
\end{exa}





\begin{proof}[\textbf{Solution.}]
$\displaystyle{f_{\theta}(x) = \frac{e^{-\theta}\, \theta^x}{x!}}\hspace{1cm} x  = 0\, , \, 1\, , \, 2\, ,\, \cdots\cdots\cdots$
\begin{align*}
f_{\theta}(x_1\, , \, x_2\, , \, \cdots\cdots\, , \, x_n) \, & = \, \prod^n_{i = 1}\, \frac{e^{-\theta}\, \theta^{x_i}}{x_i!}\,
 = \, \frac{e^{n \theta}\, \theta^{\sum x_i}}{\prod^n_{i= 1}\, x_i!}\,\hspace{1.4cm} x_i = 0\, , \, 1\, , \, 2\, , \, \cdots
\end{align*}

$$T \, = \, \displaystyle{\sum^n_{i =1}X_i}\, \thicksim\, POI(n\theta)$$

$$f_{\theta}(t) = \, \frac{e^{-n\,\theta}\, \left(n\, \theta\right)^t}{t!}\hspace{1.5cm} t = 0\, , \, 1\, ,\, 2\, , \, \cdots\cdots\cdots$$


\begin{align*}
f\left(x/T = t\right) & = \, P\left(X_1 = x_1\, , \, \cdots\cdots\, X_n = x_n\big/\, T = t\right)\\\\
& = \, \frac{P\left(X_1 = x_1\, , \, \cdots\cdots\, , \, X_n = n\, , \, \sum^n_{i = 1}X_i = t\right)}{P_{\theta}\left(T = t\right)}\\\\
& = \, \frac{e^{-n\,\theta}\, \theta^{\sum x_i}}{\prod\,x_i!}\, \times\, \frac{t!}{e^{-n\,\theta}\, \left(n\theta\right)^t}\hspace{1cm}\text{if}\hspace{0.5cm} \sum^n_{i = 1} x_i = t
\end{align*}
\begin{align*}
& = \, \frac{\theta^{\sum x_i}}{\prod^n_{i = 1}\, x_i!}\, \times \, \frac{t!}{n^t\, \theta^t}\hspace{1.8cm}\text{if}\hspace{0.5cm} \sum^n_{i = 1} x_i = t\\\\
& = \, \frac{t!}{\prod^n_{i = 1}\, x_i!}\, \left(\frac{1}{n}\right)^t\\\\
& = \, \frac{t!}{x_1!\, x_2!\, \cdots\cdots\, x_n!}\, \left(\frac{1}{n}\right)^t\hspace{0.3cm}, \hspace{1.5cm}\text{if}\hspace{0.2cm} \sum^n_{i = 1} x_i = t
\end{align*}

$$\text{i.e}\hspace{1.5cm} X\big/\, T = t\, \thicksim\, MULT\left(t\, , \, \frac{1}{n}\, , \, \cdots\cdots\, \frac{1}{n}\right)$$

Since $f\left(x\, /\, T = t\right)$ does not depend on $\theta, {T = \sum^n_{i = 1}X_i}$ is a sufficient statistics for this model.
\end{proof}




\begin{exa}
Let $\, X_1\, , \, X_2\, , \, \cdots\cdots\, , \, X_n\,$  be a random sample from $\, X\thicksim N(\theta\, , \, \sigma^2)\,$ where $\sigma^2$ is known. Show that $\, T = \sum^n_{i = 1} X_i \, $ is sufficient for $\, \theta$.
\end{exa}



\begin{proof}[\textbf{Solution.}]
$\displaystyle{f_{\theta}(x) = \frac{e^{-1/2\sigma^2\left(x - \theta\right)^2}}{\sqrt{2\, \pi\, \sigma^2}}}\hspace{1cm},\hspace{1.5cm} -\infty < x < \infty$

\begin{align*}
f_{\theta}(x_1\, , \, \cdots\cdots\, , \, x_n) \, & = \, \prod^n_{i = 1}\, \frac{e^{-1/2\sigma^2\left(x_i - \theta\right)}}{\sqrt{2\, \pi\, \sigma^2}}
 = \frac{e^{-1/2\sigma^2\, \sum^n_{i = 1}\left(x_i - \theta\right)^2}}{\left(2\pi\, \sigma^2\right)^{n/2}}
\end{align*}

$$ T = \sum^n_{i = 1} X_i \, \thicksim\, N\left(n\theta\, , \, n\sigma^2\right)$$

$$f_{\theta}(t) = \, \frac{e^{\frac{-1}{2n\sigma^2}\, (t - n\theta)^2}}{\sqrt{2\pi n\sigma^2}}$$



\begin{align*}
f(x_1\, , \, x_2\, , \, \cdots\cdots\, , x_n/\, T = t) \, & = \, \frac{f_{\theta}(x_1\, , \, x_2\, , \, \cdots\cdots\, , \, x_n)}{f_{\theta}(t)}\hspace{1.5cm}\text{if}\hspace{0.5cm} \sum x_i = t
\end{align*}
\begin{align*}
& = \, \frac{\frac{e^{\frac{-1}{2\sigma^2}\sum (x_i - \theta)^2}}{\left(2\pi\, \sigma^2\right)^{n/2}}}{\dfrac{e^{\frac{-1}{2n\sigma^2}(t -n\theta)^2}}{\sqrt{2\pi n\sigma^2}}}\\\\
& = \, \frac{\left(2\pi n\sigma^2\right)^{1/2}\, e^{\frac{-1}{2\sigma^2}\left(\sum x_i^2 - 2\theta \sum x_i + n\theta^2\right)}}{\left(2\pi\, \sigma^2\right)^{n/2}\, e^{\frac{-1}{2n\sigma^2}\left(t^2 - 2nt\theta + n^2\theta^2\right)}}\\\\
& = \, \frac{\sqrt{n}}{\left(\sqrt{2\pi\, \sigma^2}\right)^{n - 1}}\, e^{\frac{-1}{2\sigma^2}\left(\sum x_i^2 - t^2\right)}\hspace{1.5cm}\text{if}\hspace{0.5cm} \sum x_i = t\\\\
& = \, \frac{\sqrt{n}}{\left(\sqrt{2\pi\, \sigma^2}\right)^{n - 1}}\, e^{\frac{-1}{2n\sigma^2}\, n\left(\sum x_i^2 - \frac{t^2}{n}\right)}\hspace{1.5cm} \text{if}\hspace{0.5cm} \sum x_i = t
\end{align*}
which does not depend on $\theta$. Therefore $\, {T = \sum^n_{i = 1} X_i}\,\, $ is sufficient for $\theta$.
\end{proof}








\begin{problem}
Let $\, X_1\, , \, X_2\, , \, \cdots\cdots\, , \, X_n\,$  be a random sample from $X\thicksim BIN(1,\theta)\,$ and let $\, T = \sum^n_{i = 1} X_i$
\begin{enumerate}
\item[(a).] Show that $T$ is a sufficient statistic for this model.
\item[(b).] Explain how you would generate data with the same distribution as the original data using only the value of the sufficient statistic.
\item[(c).] Let $\, U = \begin{cases}
U\left(X_1\right) = 1\, & \text{if}\hspace{0.3cm} X_1 = 1\\
0 & \text{otherwise}\\
\end{cases}
$\\
Find $\, E_{\theta}\left(U\right)\, $ and $\, E\left(U\big/\, T  = t\right)$
\end{enumerate}
\end{problem}



To use the above definition, we must guess a statistic $T(X)$ to be sufficient, however, the next theorem allows us to find a sufficient statistic by simple inspection of  the pdf.


\begin{thm}
Suppose $X$ has probability (density) function $\, \{f_{\theta}(x)\, :\, \theta\in \Omega\}\,$ and $\, T(X)\,$ is a statistic. Then $\, T(X)\,$ is a sufficient statistic for $\, \{f_{\theta}(x)\, :\, \theta\in \Omega\}\,$ if and only if there exist two non negative functions $g(\cdot)$ and $h(\cdot)$ such that 
$\, f_{\theta}(x) = g\left(T(x);\, \theta\right)\, h(x)\,$
for all $\, x \, , \, \theta\in \Omega$.
\end{thm}


\newpage
\begin{proof}
``Discrete Case"\\
Suppose $T$ is a sufficient statistic for the model, 
\begin{align*}
f_{\theta}(x)\, & = \, P\left(X_1 = x_1\, ,\, X_2 = x_2\, , \,  \cdots\cdots\cdots\,  , X_n = x_n\right)\\\\
& = \, \underbrace{P\left(X_1 = x_1\, , \, X_2 = x_2\, , \, \cdots\cdots X_n = x_n\, , \, T = t\right)\,}_{h(x)}\, \underbrace{P(T = t)}_{g(x)}
\end{align*}
$h(x)$ is a function of $x$ and does not depend $\theta$ by definition and $g(x)$ depends only $\theta$ and $x$.\\
``Only if"\\
Suppose that $\, f_{\theta}(x) = g\left(T(x);\, \theta\right)\, h(x)\,$ for all $x\, , \, \theta\in \Omega$. We need to show that $f(x/ T = t)$ does not depend on  $\theta$.
\begin{align*}
P_{\theta}\left(T = t\right) \, & = \, \sum_{x:\, T(x)  = t} \, g\left(T(x);\, \theta\right)\, h(x)\,
 = \, g\left(T(x);\, \theta\right)\, \sum_{x:\, T(x) = t}\, h(x)
\end{align*}

\begin{align*}
\therefore\hspace{0.5cm} f\left(x/\, T = t\right) \, & = \, \frac{P_{\theta}(X = x)}{P(T = t)}\hspace{1.5cm} \text{if}\hspace{0.5cm} T(x) = t\\
& = \, \frac{g\left(T(x); \theta\right)\, h(x)}{g\left(T(x);\, \theta\right)\, \sum_{x:\, T(x) = t} \, h(x)} \hspace{1.5cm} \text{if}\hspace{0.5cm} T(x) = t\\
& = \, \frac{h(x)}{\sum_{x:\, T(x) = t}\, h(x)}
\end{align*}
which does not depend on $\, \theta$.\\

$\therefore\hspace{0.3cm}T(X)\,$ is a sufficient statistic.
\end{proof}











\begin{exa}
Let $\, X_1\, , \, X_2\, , \, \cdots\cdots\, , \, X_n\,$  be a random sample from the $N(\mu, \sigma^2)$ distribution. Show that $\, T = \left( \overline{X}\, , \, S^2\right)\, $ is a sufficient statistic for this model.
\end{exa}


\begin{proof}[\textbf{Solution.}]
$\displaystyle{f_{\theta}(x) = \, \frac{1}{\sqrt{2\pi \sigma^2\,}}\, e^{\frac{-1}{2\sigma^2}(x -\mu)^2}}\hspace{1cm}, \hspace{1cm}\theta = \left(\mu, \sigma^2\right)$

\begin{align*}
f_{\theta}(x_1\, , \, x_2\, , \, \cdots\cdots\, , \, x_n) & = \, \prod^n_{i = 1}\, \frac{1}{\sqrt{2\pi\sigma^2\,}}\, e^{\frac{-1}{2\sigma^2}(x_i  - \mu)^2}\\\\
& = \, \frac{1}{\left(2\pi\sigma^2\right)^{n/2}}\,e^{\frac{-1}{2\sigma^2}\sum^n_{i=1}(x_i -\mu)^2}\\\\
& = \, \frac{1}{\left(2\pi\sigma^2\right)^{n/2}}\, e^{\frac{-1}{2\sigma^2}\left(\sum^n_{i=1}x_i^2 - 2\mu\sum x_i + n\mu^2\right)}\\\\
& = \, \underbrace{\dfrac{1}{\left(2\pi\right)^{n/2}}}_{h(x)}\,\, \underbrace{\dfrac{1}{\left(\sigma^2\right)^{n/2}}\, e^{\frac{-1}{2}\, \left(\sum_i x_i^2 - 2\mu \sum x_i + n\mu^2\right)}}_{g\left(T(x);\, \theta\right)}
\end{align*}


i.e $\, {T_1(X) = \left(\sum^n_{i = 1}X_i\, , \, \sum^n_{i = 1} X_i^2\right)}\,\,$ is a sufficient statistic for this model.

\begin{align*}
f_{\theta}(x_1\,, \, x_2\, , \, \cdots\cdots\, , \, x_n) & = \frac{1}{\left(2\pi\sigma^2\right)^{n/2}}\, e^{\frac{-1}{2\sigma^2}\sum^n_{i = 1}\left(x_i - \overline{x} + \overline{x} - \mu\right)^2}\\\\
& = \, \frac{1}{\left(2\pi\sigma^2\right)^{n/2}}\, e^{\frac{-1}{2\sigma^2}\sum^n_{i = 1}\left[(x_i - \overline{x})^2 + \underbrace{2(x_i - \overline{x})(\overline{x} - \mu)}_0 + (\overline{x} - \mu)^2\right]}\\\\
& = \, \frac{1}{\left(2\pi\sigma^2\right)^{n/2}}\, e^{\frac{-1}{2\sigma^2}\left(\sum^n_{i = 1} (x_i - \overline{x})^2 \, + \, n(\overline{x} - \mu)^2\right)}\\\\
& = \, \frac{1}{\left(2\pi\sigma^2\right)^{n/2}}\, e^{\frac{-1}{2\sigma^2}\left((n - 1)S^2 \, + \, n(\overline{x} - \mu)^2\right)}\\\\
& = \, \underbrace{\frac{1}{\left(2\pi\right)^{n/2}}}_{h(x)}\,\, \underbrace{\frac{1}{\left(\sigma^2\right)^{n/2}}\, e^{\frac{-1}{2\sigma^2}\left((n - 1)S^2\, + \, n(\overline{x} - \mu)^2\right)}}_{g\left(T(x);\, \theta\right)}
\end{align*}
i.e $\, \, T(x) = \left(\overline{X}\, , \, S^2\right)\, $ is a sufficient statistic.
\end{proof}

%%%%


\textbf{Note:}\\
Sufficient statistics are not unique.




\begin{exa}
Let $\, X_1\, , \, X_2\, , \, \cdots\cdots\, , \, X_n\,$  be a random sample from the $\, WEI(1,\theta)\,$ distribution.\\
Find a sufficient statistic for this model.
\end{exa}


\begin{proof}[\textbf{Solution.}]
$\displaystyle{f_{\theta}(x) = \theta\, x^{\theta - 1}\, e^{-x^{\theta}}}$
\begin{align*}
f_{\theta}(x_1\, , \, x_2\, , \, \cdots\cdots\, , \, x_n) \, & = \, \theta^n\, \left(\prod\, x_i\right)^{\theta - 1}\, e^{-\sum^n_{i = 1} x_i^{\theta}}\\\\
& = \, \underbrace{\theta^n\, \left(\prod^n_{i = 1}\, x_{(i)}\right)^{\theta}\, e^{-\sum^n_{i = 1} x_i^{\theta}}}_{g\left(T(x);\, \theta\right)}\,\, \underbrace{\left(\prod^n_{i = 1}x_i\right)^{-1}}_{h(x)}
\end{align*}


$T(X) = \left(X_{(1)}\, , \, X_{(2)}\, , \, \cdots\cdots\, , \, X_{(n)}\right)\,$ is a sufficient statistic for this model.
\end{proof}





\begin{exa}
Let $\, X_1\, , \, X_2\, , \, \cdots\cdots\, , \, X_n\,$  be a random sample from the $UNIF(0,\theta)$ distribution. Find a sufficient statistics for $\theta$.
\end{exa}



\begin{proof}[\textbf{Solution.}]
$\displaystyle{f_{\theta}(x) = \frac{1}{\theta}}\hspace{1cm}, \hspace{1cm} 0 < x < \theta$

\begin{align*}
f_{\theta}(x_1\,, \, x_2\, , \, \cdots\cdot\, , \, x_n ) \, & = \, \frac{1}{\theta^n}\hspace{1.5cm}, \hspace{0.5cm} 0 < x_1\, \cdots\cdots\, x_n < \theta\\\\
& = \, \frac{1}{\theta^n}\hspace{1.5cm} , \hspace{0.5cm} 0 < x_{(1)} < x_{(n)} < \theta\\\\
& = \, \frac{1}{\theta^n}\hspace{1.5cm}, \hspace{0.5cm} x_{(1)} > 0\hspace{0.5cm} \text{and}\hspace{0.5cm} x_{(n)} < \theta\\\\
& = \, \frac{1}{\theta^n}\, I\left(x_{(1)} > 0\right)\, I\left(x_{(n)} < \theta\right)\\\\
& = \, \underbrace{\frac{1}{\theta^n}\, I\left(x_{(n)} < \theta\right)}_{g\left(T(x);\, \theta\right)}\, \, \underbrace{I\left(x_{(1)} > 0\right)}_{h(x)}
\end{align*}

$\therefore\, \, T(X) = X_{(n)}\,$ is sufficient statistic for $\theta$.
\end{proof}








\textbf{Note}
$$I(x\in A) = \begin{cases}
1 & \text{if}\hspace{0.3cm} x\in A\\\\
0 & \text{otherwise}\\
\end{cases}$$







\begin{problem}
Let $\, X_1\, , \, X_2\, , \, \cdots\cdots\, , \, X_n\,$  be a random sample from  $X\thicksim N(\mu, \sigma^2)\,$ with $\mu$ known. Find a sufficient statistic for $\sigma^2$.
\end{problem}



\begin{problem}
Let $\, X_1\, , \, X_2\, , \, \cdots\cdots\, , \, X_n\,$  be a random sample from the $EXP(1,\theta)$ distribution. Find a sufficient statistic for this model.
\end{problem}





\begin{coro}
Let $\, X_1\, , \, X_2\, , \, \cdots\cdots\, , \, X_n\,$  be a random sample from the model $\, \{f_{\theta}(x)\, :\, \theta\in \Omega\}$. If $T(X)$ is a sufficient statistic for the model then so is any one-to-one function of $T(X)$.
\end{coro}


\begin{proof}
Let $\, T^*(X) = r \left(T(x)\right)\, $ where $r$ is a one-to-one function with inverse $r^{-1}$. Then
\begin{align*}
f_{\theta}(x) & = \, g\left(T(x), \theta\right)\, h(x)\\\\
& = \, g\left(r^{-1}\left(T^*(x)\right)\right)\, h(x)\\\\
& = \, g^*\left(T^*(x)\right)\, h(x)
\end{align*}
where $\, g^*\left(T(x), \theta\right) = \, g\left[r^{-1}\, \left(T^*(x)\right)\right]$\\

$\therefore\,\, $ by the factorisation  Theorem $T^*(x)$ is a sufficient statistic.
\end{proof}









\section{Minimal Sufficiency}
Recall that the purpose of a sufficient statistic is to achieve data reduction without loss of information about the parameter $\theta$. Since there are many sufficient statistics in a problem, we prefer one that achieves most data reduction such that a statistic is a minimal sufficient statistic.



\begin{exa}
A sufficient statistic $T(X)$ is called a minimal sufficient statistic for $\, \{f_{\theta}(x)\, :\, \theta\in \Omega\}\,$ if for any other sufficient statistic $U(X)$, there exists a function $g(\cdot)$ such that $\, T(X) = g\left(U(X)\right)$.\\
i.e A minimal sufficient statistic can be recovered from any other sufficient statistic.
\end{exa}





\begin{thm}
Suppose there exists a function $T(X)$ such that for the model $\, \{f_{\theta}(x)\, :\, \theta\in \Omega\}\,$ an for any two sample points $x$ and $y$, the  ratio $\, \dfrac{f_{\theta}(x)}{f_{\theta}(y)}\,$ is constant as a function of $\theta$ if and only if $\, T(x) = T(y)$. Then $T(X)$ is a minimal sufficient statistic for $\theta$.
\end{thm}


In each of the previous four example $POI(\theta)\, , \, \, N(\mu, \sigma^2)\, , \, \, WEI(1,\theta)\, $ and $\, UNIF(0, \theta)$. We found a minimal sufficient statistic as it is not possible to find a further reduction of the data and still maintain the property of sufficiency.





\begin{exa}
Let $\, X_1\, , \, X_2\, , \, \cdots\cdots\, , \, X_n\,$ be a random sample from a distribution with pdf 
$$f_{\theta}(x) \, =\, \theta\, x^{\theta - 1}\hspace{0.5cm}, \hspace{0.5cm} 0 < x < 1\hspace{0.5cm}, \hspace{0.5cm} \theta > 0$$
Find a minimal sufficient statistic for $\, \{f_{\theta}(x)\, :\, \theta\in \Omega\}$
\end{exa}




\begin{proof}[\textbf{Solution.}]
\begin{align*}
f_{\theta}(x_1\, ,\, \cdots\cdots\, , \, x_n) \, & = \, \prod^n_{i = 1}\, \theta\, x_i^{\theta - 1}\,
 = \,\theta^n\, \left(\prod^n_{i = 1} x_i \right)^{\theta - 1}.
\end{align*}
Let $x = (x_1\, , \, \cdots\cdots\, x_n)\, $ and $\, y = (y_1\, , \, \cdots\cdots\, , \, y_n)\,$ be two possible outcomes, then
\begin{align*}
\frac{f_{\theta}(x)}{f_{\theta}(y)} = \frac{\theta^n\, \left(\prod^n_{i = 1} x_i\right)^{\theta - 1}}{\theta^n\, \left(\prod^n_{i = 1}y_i \right)^{\theta - 1}}\, = \,\left( \frac{\prod^n_{i = i}\, x_i}{\prod^n_{i = 1}y_i}\right)^{\theta - 1}
\end{align*}

This ratio does not depend on $\theta$ if and only if (iff) $\,\prod^n_{i = 1}\, x_i \, = \, \prod^n_{i = 1} \, y_i\,$, therefore 
$\, T(X) = \prod^n_{i = 1}\, x_i\,$ is a minimal sufficient statistic for this model.
\end{proof}






\begin{exa}
Let $\, X_1\, , \, X_2\, , \, \cdots\cdots\, ,\, X_n\, $ be a random sample from the $N(\theta, \theta^2)$ distribution. Find a minimal sufficient statistic for $\, \{f_{\theta}(x)\, :\, \theta\in \Omega\}$.
\end{exa}


\begin{proof}[\textbf{Solution.}]
\begin{align*}
f(x_1\, , \, x_2\, , \, \cdots\cdots\, , \, x_n)& = \, \prod^n_{i = 1} \frac{1}{\sqrt{2\pi \theta^2\,}}\, e^{\frac{-1}{2\theta^2}(x_i - \theta)^2}\\\\
& = \, \frac{1}{\left(2\pi\, \theta^2\,\right)^{n/2}}\, e^{\frac{-1}{2\theta^2}\, \sum^n_{i = 1}(x_i - \theta)^2}\\\\
& = \, \frac{1}{\left(2\pi\theta^2\right)^{n/2}}\, e^{\frac{-1}{2\theta^2}\, \left(\sum x_i^2 - 2\theta \sum x_i + n\theta^2\right)}
\end{align*}

Let $\, x = (x_1\, , \, x_2\, , \, \cdots\cdots\, , \, x_n)\,$ and $y = (y_1\,  ,\, y_2\, , \, \cdots\cdots\, , \, y_n)\,$ be points, then
\begin{align*}
\frac{f_{\theta}(x)}{f_{\theta}(y)} \, & = \, \frac{\dfrac{1}{\left(2\pi\theta^2\right)^{n/2}}\, e^{\frac{-1}{2\theta^2}\, \left(\sum x_i^2 - 2\theta\, \sum x_i + n\theta^2 \right)}}{\dfrac{1}{\left(2\pi\theta^2\right)^{n/2}}\, e^{\frac{-1}{2\theta^2}\left(\sum y_i^2 - 2\theta \sum y_i + n\theta^2\right)}}\\\\
& =\, e^{\frac{-1}{2\theta^2}\, \left(\sum x_i^2 - \sum y_i^2 - 2\theta\left(\sum x_i - \sum y_i\right)\right)}
\end{align*} 

This  ratio does not depend on $\theta$ if and only if $\, \sum^n_{i = 1}\, x^2_i \, = \, \sum^n_{i = 1}\, y^2_i\,$ and $\, \sum^n_{i = 1} \, x_i \, = \, \sum^n_{i = 1}\, y_i$\\

$\therefore\,\, T(X) = {\left(\sum^n_{i = 1} X_i \, , \, \sum^n_{i = 1}X^2_i\right)}\,\,$ is a minimal sufficient statistic for this model.
\end{proof}









\begin{exa}
Let $\, X_1\, , \, X_2\, , \, \cdots\cdots\, , \, X_n\,$ be a random sample from the $UNIF(0,\theta)$ distribution. Find a minimal sufficient statistic for $\, \{f_{\theta}(x):\, \theta\in \Omega\}\,$.
\end{exa}


\begin{proof}[\textbf{Solution.}]
\begin{align*}
f_{\theta}(x_1\, , \, x_2\, , \, \cdots\cdots\, , \, x_n)\, & = \, \frac{1}{\theta^n}\hspace{1cm},\hspace{0.5cm} 0 < x_1\,, \cdots\cdots\, x_n < \theta\\\\
& = \, \frac{1}{\theta^n}\hspace{1cm}, \hspace{0.5cm} 0 < x_{(1)}\, , \, x_{(n)} < \theta\\\\
& = \, \frac{1}{\theta^n}\, I\left(x_{(1)} > 0\right)\, I\left(x_{(n)} < \theta\right)
\end{align*}

Let $\, x = (x_1\, , \, x_2\, , \, \cdots\cdots\, ,\, x_n)\,$ and $\, y = (y_1\, , \, y_2\, , \, \cdots\cdots\, , \, y_n)\,$ be two possible outcomes, then
\begin{align*}
\frac{f_{\theta}(x)}{f_{\theta}(y)} \, & = \, \frac{\dfrac{1}{\theta^n}\, I\left(x_{(1)} > 0\right)\, I\left(x_{(n)} < \theta\right)}{\dfrac{1}{\theta^n}\, I\left(y_{(1)} > 0\right)\, I\left(y_{(n)}< \theta\right)}\\\\
& = \, \frac{I\left(x_{(1)} > 0\right)\, I\left(x_{(n)} < \theta\right)}{I\left(y_{(1)} > 0\right)\, I\left(y_{(n)} < \theta\right)}
\end{align*}
If $\, x_{(n)} = y_{(n)}\,$ then the ratio does not depend on $\theta$.\\\\


Suppose $\, x_{(n)} > y_{(n)}$.\\
If $\theta$ is such that $\theta > x_{(n)} > y_{(n)}\, , \, $ then 
$$\dfrac{f_{\theta}(x)}{f_{\theta}(y)}\, = \, \frac{I\left(x_{(1)} > 0\right)}{I\left(y_{(1)} > 0\right)}$$

If $\theta$ is such that $x_{(n)} > \theta > y_{(n)}\,$
$$\frac{f_{\theta}(x)}{f_{\theta}(y)}\, = \, 0$$
Therefore the value of $\, \dfrac{f_{\theta}(x)}{f_{\theta}(y)}\,$ depends on $\theta$ if $x_{(n)}  > y_{(n)}$.

A similar argument holds for the case $\, x_{(n)} < y_{(n)}$.\\
$\therefore\,\, $ the ratio does not depend on $\theta$ if and only if $x_{(n)} = y_{(n)}$.\\
Therefore $\,\, T(X) = X_{(n)}\,$ is a minimal sufficient statistics for the model.
\end{proof}












\begin{problem}
Let $\, X_1\, , \, X_2\, , \, \cdots\cdots\, , \, X_n\,$ be a random sample from the $UNIF(\theta, \theta + 1)\,$ distribution. Find  a minimal sufficient statistic for $\, \{f_{\theta}(x)\, :\, \theta \in \Omega\}$.
\end{problem}





\begin{problem}
Let $\, X_1\, , \, X_2\, , \, \cdots\cdots\, , \, X_n\,$ be a random sample from the $EXP(1, \theta)$ distribution. Show that $X_{(1)}$ is sufficient statistic for the model and find the conditional distribution for $\, X_1\, , \, X_2\, , \, \cdots\cdots\, , \, X_n\,$ given $X_{(1)} = t$. Prove that $X_{(1)}$ is a minimal sufficient statistic.
\end{problem}





\begin{problem}
Let $\, X_1\, , \, X_2\, , \, \cdots\cdots\, , \, X_n\,$ be a random sample from the $CAU(1,\theta)$ distribution. Find a minimal sufficient for $\, \{f_{\theta}(x)\, :\, \theta\in \Omega\}$.
\end{problem}








\section{Completeness}
The property of completeness is useful for determining the uniqueness of estimators and finding $UMVUES$.



\begin{defn}
A statistic $T$ or a family of distributions. $T$ is complete if for every function $h(x)\, , \, \hspace{0.1cm} E_{\theta}\left[h(T)\right] = 0\,$,  for all $\theta \in \Omega$ implies $\, P_{\theta}\left[h(T) = 0\right] = 1,\,\,$ for all $\theta \in \Omega$.
\end{defn}

i.e $\, \, E_{\theta}\left[h(T)\right] = 0\, $ implies $\, h(T) = 0\,$ with probability one.







\begin{exa}
Let $\, X_1\, , \, X_2\, , \, \cdots\cdots\, , \, X_n\,$ be a random sample from the  $N(\theta, 1)$ distribution. Consider $T(X) = \left(X_1\, ,\, \sum^n_{i = 1} X_i\right)\,$ prove that $T(X)$ is sufficient for $\, \{f_{\theta}(x)\, :\, \theta\in \Omega\}$ but not complete.
\end{exa}



\begin{proof}[\textbf{Solution.}]
\begin{align*}
f_{\theta}(x_1\,, \, x_2\, , \, \cdots\cdots\, , \, x_n)\, & = \, \prod^n_{i = 1} \frac{1}{\sqrt{2\pi}}\, e^{\frac{-1}{2}(x_i - \theta)^2}\\\\
& = \, \frac{1}{\left(2\pi\right)^{n/2}}\, e^{\frac{-1}{2}\, \sum^n_{i = 1}\left(x^2_i - 2\theta\, x_i + \theta^2\right)}\\\\
& = \, \frac{1}{\left(2\pi\right)^{n/2}}\,\, e^{\frac{-1}{2}\left(\sum x_i^2 - 2\theta \sum x_i + n\theta^2\right)}\\\\
& = \, \underbrace{e^{\frac{-1}{2}\left(n\theta^2 - 2\theta \sum x_i\right)}}_{g\left(T(x);\, \theta\right)}\, \, \underbrace{\dfrac{1}{\left(2\pi\right)^{n/2}}\, e^{\frac{-1}{2}\sum x_i^2}}_{h(x)}\\\\
& = \, e^{\frac{-1}{2}\, \left[n\theta^2 - 2\theta\left( x_1 + \sum^n_{i = 2} x_i\right)\right]}\,\, \frac{1}{\left(2\pi\right)^{n/2}}e^{\frac{-1}{2}\, \sum x_i^2}
\end{align*}

where $\, \, T(X) \, = \, {\left(X_1\, , \, \sum^n_{i = 1} X_i\right)}$.\\
$\therefore\,\,$ By the Factorisation Theorem $T(X)$ is a sufficient statistic.\\

Consider $\,\, h(T) = X_1 - \overline{X}$. Then
\begin{align*}
E\left[h(T)\right] & = E\left(X_1 - \overline{X}\right)\\
& = \, E\left(X_1\right) \, - \, E\left(\overline{X}\right)\\
& = \, \theta\, -\, \theta\\
& = \, 0\hspace{0.7cm}\text{for}\hspace{0.3cm} \theta\in \Omega
\end{align*}



\begin{align*}
P_{\theta}\left(h(T) = 0\right]  & = P\left(X_1 - \overline{X} = 0\right)\\
& = 0\\
& \neq 1
\end{align*}
$\therefore\,\, T\,$ is not complete.
\end{proof}


\begin{exa}
Let $X_1, \, X_2, \, \cdots\, , \, X_n$ be a  random sample from the Bernoulli$(\theta)$ distribution and let $T(X) = \sum^n_{i = 1} X_i$. Show that $T(X)$ is complete and sufficient for $\{f_{\theta}(x); \,\theta \in \Omega\}$.
\end{exa}


\begin{proof}[\textbf{Solution.}]
$$f_{\theta}(x_1,\, x_2, \, \cdots\, , \, x_n) = \prod^n_{i = 1} \theta^{x_i}\, \left(1 - \theta\right)^{1 - x_i} \, = \, \underbrace{\theta^{\sum x_i}\, \left(1  - \theta\right)^{1 - \sum x_i}}_{g\left(T(X);\theta\right)}\, \underbrace{1}_{h(x)}.$$
Therefore by the factorization Theorem $T(X) = \sum^{n}_{i = 1}X_i$ is a sufficient statistic.\\
Now $$T = \sum^n_{i = 1}X_i\, \thicksim\, BIN(n,\theta)$$ suppose $E_{\theta}\left[h(T)\right] = 0$ for all $\theta$. Then
\begin{align*}
E_{\theta}\left[h(T)\right] & = \sum^n_{t = 0} h(t)\, \binom{n}{t}\, \theta^t\, \left(1 - \theta\right)^{n - t}\\
&  = 0\\
& = \sum^n_{t = 0} h(t)\, \binom{n}{t}\, \left(\frac{\theta}{1 - \theta}\right)^{t}\, \left(1 - \theta\right)^{n} = 0
\end{align*}
$$\sum^n_{t = 0} h(t)\, \binom{n}{t}\, \left(\frac{\theta}{1 - \theta}\right)^t = 0\hspace{0.3cm} \text{since}\hspace{0.2cm} (1 - \theta)^n \neq 0, \, \, 0 < \theta < 1.$$
$$\sum^n_{t = 0}h(t)\, \binom{n}{t}\, r^t = 0\, ;\hspace{0.3cm} \text{where}\hspace{0.2cm}  r = \frac{\theta}{1 - \theta}\, , \, \,\, 0 < r < \infty.$$
This is a polynomial of degree $n$  in $r$ with coefficients $h(t)\, \binom{n}{t}$ for $t = 0, \, 1, \, \cdots\, , \, n$. For the polynomial to be 0, each coefficient of $r$ must be 0.\\
Since none of $\binom{n}{t}$ terms is 0, this implies $h(t) = 0, \, t = 0, \, 1, \, \cdots\, , \, n$ which means that $P_{\theta}\left[h(T) = 0\right] = 1$ for all $\theta \in \Omega$. Therefore $T$ is a complete statistic.
\end{proof}


\begin{problem}
Let $X_1, \, X_2, \, \cdots\, , \, X_n$ be a random sample from the $UNIF(0,\theta)$ distribution. Show that $T(X) = X_{(n)}$ is a complete statistic for $\{f_{\theta}(x);\, \theta \in \Omega\}$.
\end{problem}

\begin{problem}
Let $X_1, \, X_2, \, \cdots\, , \, X_n$ be a random sample from the $N(\theta, a\theta^2)$ distribution where $a>0$ is a known constant and $\theta >0$. Show that the minimal sufficient statistic is not complete.
\end{problem}

\begin{thm}
If $T(X)$ is a complete and sufficient statistic for the model $\{f_{\theta}(x);\, \theta \in \Omega\}$, then $T(X)$ is a minimal sufficient statistic.
\end{thm}


\begin{problem}
The converse of the above theorem is not true. Let $X_1, \, X_2, \, \cdots\, , \, X_n$ be a random sample from the $UNIF(\theta - 1, \theta + 1)$ distribution. Show that $T(X) = \left(X_{(1)}, X_{(n)}\right)$ is minimal sufficient. Show also that for the nonzero function $$h(T) = \frac{X_{(n)} - X_{(1)}}{2} - \frac{n - 1}{n  + 1}$$ $E_{\theta}\left[h(T)\right] =0$ for all $\theta \in \Omega$ and therefore $T$ is not a complete statistic.
\end{problem}


\begin{problem}
Let $X_1, \, X_2, \, \cdots\, , \, X_n$ be a random sample from the $UNIF(0,\theta)$ distribution. prove that $T(X) = X_{(n)}$ is a minimal sufficient statistic for this model.
\end{problem}


\begin{problem}
Let $X_1, \, X_2, \, \cdots\, , \, X_n$ be a random sample from the $EXP(1, \theta)$ distribution. Show that $T(X) = X_{(1)}$ is a minimal sufficient statistic for the model.
\end{problem}

\begin{thm}
For any random variables $X$ and $Y,\,$ $E_{\theta}(X) = E_{\theta}\left[E_{\theta}(X/Y)\right]$ and\\ $var(X) = E_{\theta}\left[var_{\theta}(X/Y)\right] \, + \, var_{\theta}\left[E_{\theta}(X/Y)\right].$
\end{thm}

\begin{thm}
If $T(X)$ is a complete statistic, then there is at most on function of $T$ that provides an unbiased estimator of the parameter $\tau(\theta)$.
\end{thm}


\begin{thm}
If $T(X)$ is a complete sufficient statistic for the model $\{f_{\theta}(x);\, \theta\in \Omega\}$  and $E_{\theta}\left[g(T)\right] = \tau(\theta)$, then $g(T)$ is the $UMVUE$ of $\tau(\theta)$.
\end{thm}


\begin{proof}
Since $T(X)$ is a complete sufficient statistic by Theorem 1.5.11, $g(T)$ is unique. Suppose $U = U(X)$ is an unbiased estimator of $\tau(\theta)$. Let $g_1(T) = E\left(U/T\right)$ (doesn't depend on $\theta$)
\begin{align*}
E_{\theta}\left[g_1(T)\right] & = E_{\theta}\left[E\left(U/T\right)\right]\\
& = E(U)\\
& = \tau(\theta).
\end{align*}
By Theorem 1.5.11 $g_1(T) = g(T)$. Now 
\begin{align*}
var_{\theta}(U) & = \underbrace{E_{\theta}\left[var(U/T)\right]}_{\geq 0}\,  + \, var_{\theta}\Big[\underbrace{E(U/T)}_{g(T)}\Big]\\
& \geq var\left[g(T)\right]
\end{align*}
i.e $\, var_{\theta}(U) \geq var\left[g(T)\right]$. Therefore $g(T)$ is the $UMVUE$.
\end{proof}



\begin{exa}
Let $X_1, \, X_2, \, \cdots\, , \, X_n$ be a random sample from the Bernoulli$(\theta)$ distribution. Find the $UMVUE$ of $\theta^2$.
\end{exa}

\begin{proof}[\textbf{Solution.}]
We already know that $T(X) = \sum^n_{i = 1} X_i$ is a complete sufficient statistic.\\
Consider $g(T) = \overline{X}$. Then $E(\overline{X}) = \theta$
$$var\left(\overline{X}\right) = E\left(\overline{X}^2\right) - \left(E(\overline{X})\right)^2$$
$$\implies\,\, \frac{var(X)}{n} = E\left(\overline{X}^2\right) - \theta^2$$
$$\frac{\theta (1 - \theta)}{n} + \theta^2 = E\left(\overline{X}^2\right)\, \implies\, \frac{\theta - \theta^2 + n\theta^2}{n} = E\left(\overline{X}^2\right)$$
$$\implies\,\, \theta + (n - 1)\theta^2 = n \, E\left(\overline{X}^2\right)$$
$$(n - 1)\theta^2 = n\, E\left(\overline{X}^2\right) - \theta$$
$$n\, E\left(\overline{X}^2\right) - E(\overline{X})  = (n - 1)\, \theta^2$$
$$\implies\,\, \frac{E\left[n\, \overline{X} - \overline{X}\right]}{n - 1} = \theta^2$$
$$\implies\,\, E\left[\frac{\overline{X}(n\, \overline{X} - 1)}{n - 1}\right] = \theta^2$$
Therefore $\, \dfrac{\overline{X}(n \, \overline{X} - 1)}{n - 1}\,$ is the $UMVUE$ of $\theta^2$.\\

or
$$E\left[\frac{\sum X_i\, \left(\sum X_i - 1\right)}{n(n - 1)}\right] = \theta^2$$

$$E\left[\frac{T(T - 1)}{n (n - 1)}\right] = \theta^2$$
Therefore $\, \dfrac{T (T - 1)}{n (n - 1)}\,$ is the $UMVUE$ of $\theta^2$.
\end{proof}


\begin{exa}
Let $X_1, \, X_2, \, \cdots\, , \, X_n$ be a random sample from the $UNIF(0,\theta)$ distribution. Find the $UMVUE$ of $\tau(\theta) = \theta$.
\end{exa}

\begin{proof}[\textbf{Solution.}]
$T = X_{(n)}$ is a complete sufficient statistic (Problem 1.5.4). By Example 1.2.12
$$E(T) = \frac{n}{n + 1}\, \theta\,\, \implies\,\, \frac{n + 1}{n}\, E(T) = \theta$$
$$\implies\,\, E\left(\frac{(n + 1)}{n}\, T\right) = \theta.$$
Therefore $\, \dfrac{n + 1}{n}\, T\,$ is the $UMVUE$ of $\tau(\theta) = \theta$.
\end{proof}

\begin{problem}
Let $X_1, \, X_2, \, \cdots\, , \, X_n$ be a random sample from the Bernoulli$(\theta)$ distribution. Find the $UMVUE$ of $\tau(\theta) = \theta(1 -\theta)$.
\end{problem}

\begin{problem}
Let $X_1, \, X_2, \, \cdots\, , \, X_n$ be a random sample from the $EXP(1,\theta)$ distribution. Find the $UMVUE$ of $\tau(\theta) = \theta^2$.
\end{problem}

\begin{thm}[Lehmman - Scheffe']
If $T(X)$ is a complete sufficient statistic for the model $\{f_{\theta}(x):\, \theta \in \Omega\}$ and $U(X)$ is any unbiased estimator of $\tau(\theta)$, then $E\left(U/T\right)$ is the $UMVUE$ of $\tau(\theta)$.
\end{thm}




\section{The Exponential Family}

\begin{defn}
Suppose $X = (X_1, \, \cdots\, , \, X_p)$ has a (joint) probability (density) function of the form 
$$f_{\theta}(x) = C(\theta)\, \exp\left\{\sum^k_{j = 1} q_j(\theta)\, T_j(x)\right\}\, h(x)\hspace{0.5cm}\cdots\cdots\hspace{0.3cm} (1.1)$$
for functions $q_j(\theta), \, T_j(x), \, h(x), \, C(\theta)$. Then we say that $f_{\theta}(x)$ is a member of the exponential family of densities. We call 
$(T_1(X), \, \cdots\, , \, T_k(X))$ the natural  sufficient statistic.
\end{defn}


Note that the natural sufficient statistic is not unique. e.g multiplication of $T_j$ by a constant and division of $q_j$ by the same constant results in the same function $f_{\theta}(x)$.


\begin{exa}
Show that $T(X) = (T_1(X), \, \cdots\, , \, T_k(X))$ is a sufficient statistic for the model $\{f_{\theta}(x):\, \theta\in \Omega\}$ where $f_{\theta}(x)$ has the form $(1.1)$.
\end{exa}

\begin{proof}[\textbf{Solution.}]
$$f_{\theta}(x) = \underbrace{C(\theta)\, \exp\left\{\sum^k_{j = 1} q_j(\theta)\, T_j(x)\right\}}_{g\left(T(X), \theta\right)}\, \underbrace{h(x)}_{h(x)}$$
where $T(X) = (T_1(X), \, \cdots\, , \, T_k(X))$. Therefore by the Factorization Theorem $T(X)$ is a sufficient statistic.
\end{proof}



\begin{exa}
Show that the $BIN(n,\theta)$ distribution has an exponential family distribution and find the natural sufficient statistic.
\end{exa}

\begin{proof}[\textbf{Solution.}]
$X\, \thicksim\, BIN(n,\theta)$
\begin{align*}
f_{\theta}(x) & = \binom{n}{x}\, \theta^x\, (1 - \theta)^{n-x}\hspace{0.3cm} , \hspace{0.6cm} x = 0, \, 1, \, \cdots\, , \, n\\
& = (1 - \theta)^n\, \left(\frac{\theta}{1 - \theta}\right)^x\, \binom{n}{x}\\
& = (1 - \theta)^n\, \exp\left\{\log\left(\frac{\theta}{1 - \theta}\right)^x\right\}\, \binom{n}{x}\\
& = (1 - \theta)^n\, \exp\left\{x\, \log\left(\frac{\theta}{1 - \theta}\right)\right\}\, \binom{n}{x}\\
& = \underbrace{(1 - \theta)^n}_{C(\theta)}\, \exp\Bigg\{\underbrace{\left[\log\left(\frac{\theta}{1 - \theta}\right)\right]}_{q(\theta}\, \underbrace{x}_{Tx)}\Bigg\} \, \underbrace{\binom{n}{x}}_{h(x)}.
\end{align*}
$\therefore f_{\theta}(x)$ is a member of the exponential family and the natural sufficient statistic is $T(X) = X$.
\end{proof}


\textbf{Note:} The exponential family is closed under repeated independent sampling.



\begin{thm}
Let $X_1, \, \cdots\, , \, X_n$ be a random sample from a distribution with probability (density) function given by $(1.1)$. Then $(X_1, \, \cdots\, , \, X_n)$ also has an exponential form with joint probability (density) function
$$f_{\theta}(x_1, \, \cdots\, , \, x_n) = \left(C(\theta)\right)^n\, \exp\left\{\sum^k_{j = 1} q_j(\theta)\, \sum^n_{i = 1} T_j(x_i)\right\}\, \prod^n_{i = 1}h(x_i)$$
i.e $C$ is  replaced by $C^n$, $T_j(x)$ by $\sum^n_{i = 1} T_j(x_i)$ and $h(x)$ by $\prod^n_{i = 1} (x_i)$.\\
The natural sufficient statistic is $$\left(\sum^n_{i = 1} T_1(X_i), \, \cdots\cdots\, ,\, \sum^n_{i = 1} T_k(X_i)\right).$$ 
\end{thm}



\begin{exa}
Let $(X_1, \, \cdots\, , \, X_n)$ be a random sample from the $POI(\theta)$ distribution. Show that $(X_1, \, \cdots\, , \, X_n)$ is a member of the exponential family.
\end{exa}


\begin{proof}[\textbf{Solution.}]
$X\, \thicksim\, POI(\theta)$
\begin{align*}
f_{\theta}(x)    & = \frac{e^{-\theta}\, \theta^x}{x!}\hspace{0.6cm}, \hspace{0.5cm} x   = 0, \, 1, \, \cdots\hspace{0.3cm} \theta > 0 \\
& = e^{-\theta}\, \exp\left(\log \theta^x\right)\, \frac{1}{x!}\\
& = \underbrace{e^{-\theta}}_{C(\theta)}\, \exp\Bigg\{\underbrace{\left(\log \theta\right)}_{q(\theta)}\, \cdot\, \underbrace{x}_{T(x)}\Bigg\}\, \underbrace{\frac{1}{x!}}_{h(x)}
\end{align*}
$\therefore\, f_{\theta}(x)$ is a member of the exponential family and the natural sufficient statistic is $T(X) = X$.\\
$\therefore\,$ by Theorem 1.6.4 $(X_1, \, \cdots\, , \, X_n)$ has an exponential family (EF) distribution and the natural sufficient statistic is $T(X) = \sum^n_{i = 1} X_i$.\\

OR
\begin{align*}
f_{\theta}(x_1, \, \cdots\, , \, x_n) & = \prod^n_{i = 1} \frac{e^{-\theta}\, \theta^{x_i}}{x_i!}
 = e^{-n\theta}\, \theta^{\sum^n_{i = 1} x_i}\, \prod^n_{i = 1} \frac{1}{x_i!}\\
& = \underbrace{e^{-n\theta}}_{C(\theta)}\, \exp \Bigg\{\underbrace{\left(\log \theta\right)}_{q(\theta)}\, \underbrace{\sum^n_{i = 1} x_i}_{T(X)}\Bigg\}\, \underbrace{\prod^n_{i = 1}\frac{1}{x_i!}}_{h(x)}.
\end{align*}
\end{proof}



\begin{defn}
The canonical form of the exponential family $$f_{\eta}(x) = C(\eta)\, \exp\Bigg\{\sum^k_{j = 1} \eta_j\, T_j(x)\Bigg\}\, h(x)$$ is obtained from equation (1.1) by replacing $q_j(\theta)$ by a new parameter $\eta_j$.\\
The natural parameter space in this form is the set of all values of $\eta$ for which $f_\eta(x)$ is integrable i.e
$$\left\{\eta:\, \int^{\infty}_{-\infty} f_{\eta}(x)\, dx\, < \infty \right\}$$
If $X$ is discrete the integral is replaced by a sum over all $x$ such that $f_{\eta}(x) > 0$.
\end{defn}


\textbf{Note:}\\
If the statistic satisfies a linear constraint, for example $\sum^n_{j = 1} T_j(x)= 0,$ then the number of terms can be reduced, otherwise the parameters $\eta_j$ are not statistically meaningful.




\begin{defn}
$X$ is said to have a regular exponential family distribution if its canonical form is of full rank in the sense that neither the $T_j$ nor the $\eta_j$ satisfy any linear constraints and the natural parameter space contains a $k-$dimensional rectangle.
\end{defn}


\textbf{Note:} By Theorem 1.6.4, if $X_i$ has a regular exponential family distribution then $X = (X_1, \, \cdots\, , \, X_n)$ also has a regular exponential family (REF) distribution.


\begin{exa}
Show that $X\, \thicksim\, BIN(n,\theta)$ has a regular exponential family distribution.
\end{exa}


\begin{proof}[\textbf{Solution.}]
From example 1.6.3
$$f_{\theta}(x) = (1 - \theta)^n\, \exp\left\{\left[\log\left(\frac{\theta}{1 - \theta}\right)\right]\, x\right\}\, \binom{n}{x}.$$
Let $\,  \eta = \log\left(\frac{\theta}{1 - \theta}\right) = q(\theta)$
\begin{align*}
0 &< \theta < 1,\\
0 & < \frac{\theta}{1 - \theta} < \infty\\
-\infty &  < \eta < \infty
\end{align*}
We make $\theta$ the subject 
\begin{align*}
e^{\eta} & = \frac{\theta}{1 - \theta}\\
\theta &  = \frac{e^{\eta}}{1 + e^{\eta}}\\
1 - \theta &   = \frac{1}{1 + e^{\eta}}
\end{align*}
The natural parameter space is 
$$\Omega_{\eta} = \left\{\eta:\, -\infty < \eta < \infty \right\}$$
Therefore the canonical form is
$$f_{\eta}(x)  = \left(\frac{1}{1 + e^{\eta}}\right)^n\, \exp\left\{\eta \, x \right\}\, \binom{n}{x}$$
no linear constraints on $T(X) = X$ or $\eta$ and 
$\Omega_{\eta}$ contains a 1-dimensional rectangle.\\
$\therefore\, X\, \thicksim\, BIN(n,\theta)$ has a REF distribution.
\end{proof}

\begin{thm}
If $X$ has a regular exponential family distribution then $(T_1(X), \, \cdots\, , \, T_k(X))$ is a complete sufficient statistic.
\end{thm}


\begin{exa}
Let $X_1, \, \cdots\, , \, X_n$ be a random sample from the $N(\mu, \sigma^2)$ distribution. Find a  complete sufficient statistic for this model. Find the $UMVUE$s of $\mu$ and $\sigma^2$.
\end{exa}


\begin{proof}[\textbf{Solution.}]
$X\, \thicksim\, N(\mu, \sigma^2)\hspace{0.3cm} , \hspace{0.3cm} \theta = (\mu, \sigma^2)$
$$\Omega = \left\{(\mu, \sigma^2):\, -\infty < \mu < \infty, \, \, \sigma^2 > 0\right\}$$
\begin{align*}
f_{\theta}(x) & = \frac{1}{\sqrt{2\pi \sigma^2}}\, e^{-\frac{1}{2\sigma^2}(x - \mu)^2}\\
& = \frac{1}{\sigma^2}\, e^{-\frac{1}{2\sigma^2}(x^2 - 2\mu\, x + \mu^2)}\, \frac{1}{\sqrt{2\pi}}\\
& = \underbrace{\frac{1}{\sigma^2}}_{C(\theta)}\, e^{-\frac{\mu^2}{2\sigma^2}}\, \exp\Bigg\{\underbrace{-\frac{1}{2\sigma^2}}_{q_1(\theta)}\, \underbrace{x^2}_{T_1(x)}\, + \, \underbrace{\frac{\mu}{\sigma^2}}_{q_2(\theta}\, \underbrace{x}_{T_2(x)}\Bigg\}\, \underbrace{\frac{1}{\sqrt{2\pi}}}_{h(x)}
\end{align*}
$\therefore\, f_{\theta}(x)$ is a member of the exponential family with natural sufficient statistic\\ $\, T(X) = (X, X^2)$.\\

Let
\begin{align*}
\eta_1 & = q_1(\theta)  = -\frac{1}{2\sigma^2}\\\\
\eta_2 & = q_2(\theta) = \frac{\mu}{\sigma^2}
\end{align*}
The natural parameter space is $$\Omega_{\eta} = \left\{(\eta_1, \eta_2):\, \eta_1 < 0, \, \, -\infty < \eta_2 < \infty\right\}$$
\begin{figure}[H]
\centering
\begin{tikzpicture}
\draw[-Stealth](-2,0)--(2,0)node[below]{$\eta_1$};
\draw[-Stealth](0,-2)--(0,2)node[above]{$\eta_2$};
\path[pattern = north west lines](-2,1.9)--(0,1.9)--(0,-1.9)--(-2,-1.9)--cycle;
\draw(-0.5,1)--(-0.5,-1)--(-1.5,-1)--(-1.5,1)--cycle;
\end{tikzpicture}
\end{figure}
\begin{enumerate}
\item[-] No linear constraints on $T_j \, '$s or $\eta_j\, '$s.
\item[-] $\Omega_{\eta}$ contains a 2-dimensional rectangle
\end{enumerate}
$$f_{\eta}(x) = C(\eta)\, \exp\left\{\eta_1\, x^2 + \eta_2\, x\right\}\, \frac{1}{\sqrt{2\pi}}$$
Therefore $X$ has a REF distribution.\\
$\therefore\, (X_1, \, \cdots\, , \, X_n)$ has REF distribution with natural sufficient statistics $$T(X) = \left(\sum^n_{i = 1}X_i \, , \, \sum^n_{i = 1} X_i^2\right).$$
$T(X)$ is a completer sufficient statistic (Theorem 1.6.9).\\
$\implies\,\left(\overline{X}\, , \, S^2\right)$ is also a complete sufficient statistic since $E_{\theta}(\overline{X}) = \mu$ and $E_{\theta}(S^2) = \sigma^2$. Therefore  $\overline{X}$ is a $UMVUE$ of $\mu$ and $S^2$ is a $UMVUE$ of $\sigma^2$.
\end{proof}





\begin{exa}
Show that $X\thicksim N(\theta, \theta^2)$ does not have a regular exponential family distribution.
\end{exa}

\begin{proof}[\textbf{Solution.}]
$f_{\theta}(x)  = \frac{1}{\sqrt{2\pi\theta^2}}\, e^{-\frac{1}{2\theta^2}(x  - \theta)^2}$
$$\Omega = \left\{\theta:\, \theta \neq 0\right\}$$
\begin{align*}
f_{\theta}(x) & = \frac{1}{\sqrt{2\pi\, \theta^2}}e^{-\frac{1}{2\theta^2}(x - \theta)^2}\\
& = \frac{1}{|\theta|}\, e^{-\frac{1}{2\theta^2}(x^2 - 2\theta x + \theta^2)}\, \frac{1}{\sqrt{2\pi}}\\
& = \underbrace{\frac{1}{|\theta|}}_{C(\theta)}\, \exp\Bigg\{\underbrace{-\frac{1}{2\theta^2}}_{q_1(\theta)}\, \underbrace{x^2}_{T_2(x)}\, + \, \underbrace{\frac{1}{\theta}}_{q_2(x)}\, \underbrace{x}_{T_2(x)}\Bigg\}\, \underbrace{\frac{e^{-\frac{1}{2}}}{\sqrt{2\pi}}}_{h(x)}
\end{align*}
$\therefore\, \, f_{\theta}(x)$ is  a member of the REF with natural sufficient statistic $T(X) = (X, X^2)$.\\

In Canonical form 
\begin{align*}
\eta_1 & = q_1(\theta) = -\frac{1}{2\theta^2}\\
\eta_2 & = q_2(\theta) = \frac{1}{\theta}
\end{align*}
We have that $\eta_1 = -\frac{1}{2\theta^2}$ and $\eta_2 = \frac{1}{\theta}$.
$$\implies\hspace{0.5cm} \eta_1 = -\frac{1}{2}\, \eta^2_2$$
$$\Omega_{\eta} = \left\{(\eta_1, \eta_2):\, \eta_1 = -\frac{1}{2}\, \eta^2_2\, , \, \, \eta_2 \neq 0\right\}.$$
\begin{figure}[H]
\centering
\begin{tikzpicture}
\draw[-Stealth](-2,0)--(2,0)node[right]{$\eta$};
\draw[-Stealth](0,-2)--(0,2);
\draw(0,0)parabola(2,-1.5)
(0,0)parabola(-2,-1.5);
\draw(0,0)circle(0.2);
%(0,0)node[shape = circle, draw, scale = 2, fill = white]{};
\end{tikzpicture}
\end{figure}
\begin{enumerate}
\item[-] Points in $\Omega_{\eta}$ lie on the curve $\eta_1 = -\frac{1}{2}\eta^2_2$.
\item[-] $\Omega_{\eta}$ doesn't contain a 2-dimensional rectangle.
\end{enumerate}
$\therefore\, f_{\theta}(x)$ doesn't have a REF distribution.
\end{proof}
\vspace*{10pt}




\begin{exa}
Let $(X_1, \, X_2, \, X_3)\thicksim MULT(n, \theta_1, \theta_2, \theta_3)$. Find a $UMVUE$ of $\tau(\theta) = \theta_1\theta_2$.
\end{exa}

\begin{proof}[\textbf{Solution.}]
$$f_{\theta}(x_1, x_2, x_3) = \frac{n!}{x_1!\, x_2!\, x_3!}\theta_1^{x_1}\, \theta_2^{x_2}\, \theta_3^{x_3}$$
$$x_i = 0, \, 1, \, \cdots$$
$$x_1  + x_2 + x_3  = n$$
$$0 < \theta_i < 1$$
$$\theta_1 + \theta_2 + \theta_3 = 1$$
\begin{align*}
f_{\theta}(x_1, x_2,x_3) & = \exp\left\{\log\left(\theta_1^{x_1}\, \theta_2^{x_2}\, \theta_3^{x_3}\right)\right\}\, \frac{n!}{x_1!\, x_2!\, x_3!}\\
& = \exp\left\{\log \theta_1^{x_1} + \log\theta_2^{x_2} + \log \theta_3^{x_3}\right\}\, \frac{n!}{x_1!\, x_2!\, x_3!}\\
& = \exp\left\{\left(\log \theta_1\right)\, x_1  + \left(\log \theta_2\right)\, x_2 + \left(\log\theta_3\right)\, x_3\right\}\, \frac{n!}{x_1!\, x_2!\, x_3!}
\end{align*}
Therefore $\, f_{\theta}(x_1, x_2, x_3)$ is a member of the EF with natural sufficient statistic\\ $T(X) = (X_1, \, X_2, \, X_3) = (T_1, \, T_2, \, T_3).$\\
$T_1 + T_2 + T_3 = n $ - a linear constraint on $T_j\,'$s. We therefore rewrite it as 
\begin{align*}
& f_{\theta}(x_1,x_2)  =  \frac{n!}{x_1!\, x_2!\, (n - x_1 - x_2)!}\, \theta_1^{x_1}\, \theta_2^{x_2}\, (1 - \theta_1 - \theta_2)^{n - x_1 - x_2}\\
& = (1 - \theta_1 - \theta_2)^n\, \exp\left\{\log\left[\left(\frac{\theta_1}{1 - \theta_1 - \theta_2}\right)^{x_1}\, \left(\frac{\theta_2}{1 - \theta_1 -\theta_2}\right)^{x_2}\right]\right\}\, \frac{n!}{x_1!\, x_2!\, (n - x_1 - x_2)!}\\
& = (1 - \theta_1 -\theta_2)^n\, \exp\left\{x_1\, \log\left(\frac{\theta_1}{1 - \theta_1 - \theta_2}\right)  + x_2\, \log \left(\frac{\theta_2}{1 - \theta_1 - \theta_2}\right)\right\}\, \frac{n!}{x_1!\, x_2!\, (n - x_1 - x_2)!}.
\end{align*}
$\therefore\, f_{\theta}(x_1, x_2)$ is a member of the EF with natural sufficient statistic $T = (X_1, X_2) = (T_1, T_2)$.\\
In Canonical form
$$\eta_1 = \log\left(\frac{\theta_1}{1 - \theta_1 -\theta_2}\right)$$
$$\eta_2 = \log\left(\frac{\theta_2}{1 - \theta_1 - \theta_2}\right)$$
$$\Omega_{\eta}  = \left\{(\eta_1, \eta_2):\, -\infty < \eta_1 < \infty\, , \, -\infty < \eta_2 < \infty\right\}.$$
\begin{enumerate}
\item[-] No linear constraints on $T_j\,'$s or $\eta_j\, '$s.
\item[-] $\Omega_{\eta}$ contains a 2-dimensional rectangle.
\end{enumerate}
$$f_{\eta}(x_1, x_2) = \left(\frac{1}{1  + e^{\eta_1} + e^{\eta_2}}\right)^n\, \exp\left\{\eta_1\, x_1  + \eta_2\, x_2\right\}\, h(x)$$ is a REF distribution.\\
Therefore $T(X) = (X_1, X_2)$  is a complete sufficient statistic.
$$E(X_1) = n\theta_1\, , \, E(X_2) = n\theta_2\, ,\, Cov(X_1, X_2) = -n\theta_1\theta_2.$$
$$COV(X_1,X_2) = E(X_1X_2) - E(X_1)E(X_2)$$
$$-n\theta_1\theta_2  + n\theta_1n\theta_2 - E(X_1X_2)$$
$$(n^2 - n)\theta_1\theta_2 = E(X_1X_2)$$
$$\frac{E(X_1X_2)}{n(n-1)} = \theta_1 \theta_2$$
$$E\left(\frac{X_1X_2}{n(n- 1)}\right) = \theta_1\theta_2$$
$\therefore\, T(X) = \frac{X_1X_2}{n(n-1)}\, $ is a $UMVUE$ of $\tau(\theta) = \theta_1\theta_2$.
\end{proof}




\begin{exa}
Let $X_1, \, \cdots\, , \, X_n$ be a random sample from the $POI(\theta)$ distribution. Find the $UMVUE$ of $\tau(\theta) = e^{-\theta}$.
\end{exa}


\begin{proof}[\textbf{Solution.}]
$$f_{\theta}(x) = \frac{e^{-\theta}\, \theta^x}{x!}\,\hspace{0.2cm} , \hspace{0.3cm} x  = 0, \, 1,\, 2, \, \cdots\cdots$$
Consider $\, P(X_1  = 0) = e^{-\theta} = \tau(\theta)$.\\

Let 
$\hspace{0.3cm} U = U(X_1) = \begin{cases}
1 & \text{if}\, \, \hspace{0.3cm} X_1 = 0\\
0 & \text{otherwise}\\
\end{cases}
$
\begin{align*}
E(U) & = 1\times P(X_1 = 0) + 0\times P(X_1  = 0)\\
& = P(X_1 = 0)\\
&  = e^{-\theta}.
\end{align*}
$POI(\theta)$ belongs the the REF distribution. Therefore $T = \sum^n_{i = 1} X_i$ is a complete sufficient statistic.
$$i.e\hspace{0.5cm} f_{\theta}(x) = e^{-\theta}\, \exp\left\{\left(\log \theta\right)\, x \right\}\, \frac{1}{x!}$$
$$\eta = \log \theta$$
$$\Omega  = \left\{\eta:\, -\infty < \eta < \infty\right\}$$
\begin{enumerate}
\item[-] $\Omega_{\eta}$ contains  a 1-dimensional rectangle.
\item[-] No constraints on $T$ or $\eta$
\end{enumerate}
By Lehmaann-Scheff`es theorem $E\left[U/T = t\right]$ is the $UMVUE$ of $\tau(\theta)$
\begin{align*}
E\left[U/T = t\right] & = 1\times P(X_1 = 0/T = t)
 = \frac{P\left(X_1 = 0\,, \, \sum^n_{i = 1} X_i = t\right)}{P(T = t)}\\
& = \frac{P\left(X_1 = 0\, , \, \sum^n_{i = 1} X_i  = t - 0\right)}{e^{-n\theta}(n\theta)^t/t!}\\
& = \frac{P(X_1 = 0)\, P\left(\sum^n_{i = 1} X_i = t\right)}{e^{-n\theta}(n\theta)^t/t!}\\
& = \frac{e^{-\theta}\, \frac{e^{-(n - 1)\theta}((n-1)\theta)^t}{t!}}{e^{-n\theta}(n\theta)^t/t!}\\
& = \frac{e^{-n\theta}\, (n - 1)^t}{e^{-n\theta}\, n^t} = \left(\frac{n - 1}{n}\right)^t\, , \hspace{0.3cm} t = 0, \, 1, \, 2, \, \cdots
\end{align*}
$\therefore\, h(T) = \left(1 - \dfrac{1}{n}\right)^T\,$ is the  $UMVUE$ of $\tau(\theta)  = e^{-\theta}$.
\end{proof}\vspace*{0.5pt}






\begin{exa}
Let $X_1, \, \cdots\, , \, X_n$ be a random sample from the $N(\theta,1)$ distribution. Find the $UMVUE$ of $\Phi(C - \theta)  = P(X\leq C)$ for some constant $C$ where $\Phi$ is the standard normal cumulative function.
\end{exa}

\begin{proof}[\textbf{Solution.}]
$\tau(\theta) = \Phi(C - \theta) = P(X\leq C)$\\

Let $\hspace{0.3cm} U(X_1) = \begin{cases}
1 &  \text{if}\hspace{0.2cm} X\leq C\\
0 & \text{otherwise}
\end{cases}$
$$E(U) = P(X\leq C)$$
$N(\theta,1)\,$ is a REF distribution. Therefore $T = \sum^n_{i = 1} X_i \,$ is a complete sufficient statistic.\\
By Lehmann-Scheffe`s Theorem $E(U/T = t)$ is the $UMVUE$ of $\tau(\theta)$.
$$f(x/T= t) = \frac{f(x , t)}{f_T(t)}  = \frac{f(x, t-x)}{f_T(t)} = \frac{f(x)\, f(t -x)}{f_T(t)}.$$
$$T = \sum^n_{i = 1} X_i\, \thicksim\, N(n\theta, n)$$
$$T - X  = \sum^{n - 1}_{i = 1} X_i \, \thicksim \, N\left((n - 1)\theta, n - 1\right)$$
\begin{align*}
& f(X/T  = t)  = \frac{\frac{1}{\sqrt{2\pi}}\exp\left\{\frac{-1}{2}(x - \theta)^2\right\}\, \cdot\, \frac{1}{\sqrt{2\pi(n - 1)}}\exp\left\{\frac{-1}{2(n-1)}(t - x) - (n - 1)\theta)^2\right\}}{\frac{1}{\sqrt{2\pi n}}\exp\left\{\frac{-1}{2n}(t - n\theta)^2\right\}}\\
& = \frac{\sqrt{n}}{\sqrt{2\pi(n - 1)}}\exp\left\{\frac{-1}{2}\left((x - \theta)^2 + \frac{1}{n - 1}(t - x - (n - 1)\theta)^2 - \frac{1}{n}(t - n\theta)^2\right)\right\}\\
&  = \frac{\sqrt{n}}{\sqrt{2\pi (n - 1)}}\, \exp\left\{\frac{-1}{2}\left((x  - \theta)^2 + \frac{(t - n\theta)^2}{n - 1} - \frac{2(t - \theta n)(x - \theta)}{n - 1} + \frac{(x  - \theta)^2}{n - 1} - \frac{(t - n\theta)^2}{n}\right)\right\}\\
& = \frac{\sqrt{n}}{\sqrt{2\pi(n - 1)}}\, \exp\left\{\frac{-n}{2(n - 1)}\left(x^2 - \frac{2tx}{n} + \frac{t^2}{n^2}\right)\right\}\\
& = \frac{\sqrt{n}}{\sqrt{2\pi(n - 1)}}\, \exp\left\{\frac{-n}{2(n-1)}\left(x - \frac{t}{n}\right)^2\right\}\\
& =  \frac{1}{\sqrt{2\pi\left(\frac{n - 1}{n}\right)}}\,\exp\left\{\frac{-1}{2\left(\frac{n - 1}{n}\right)}\left(x - \frac{t}{n}\right)^2\right\} 
\end{align*}
$$X/ T = t\, \thicksim\, N\left(\frac{t}{n}\, , \, \frac{n - 1}{n}\right)$$
\begin{align*}
E\left(U/T =  t\right) & = E\left[I(X\leq C)/ T = t\right]\\
& = P\left(X \leq C / T = t\right)\\
& = P\left(\frac{X - t/n}{\sqrt{\frac{n - 1}{n}}} \, \leq\, \frac{C - t/n}{\sqrt{\frac{n - 1}{n}}}\right)\\
&  = P\left(Z \, \leq \frac{C - t/n}{\sqrt{\frac{n - 1}{n}}}\right)\\
& = \Phi\left(\frac{C - t/n}{\sqrt{\frac{n - 1}{n}}}\right)
\end{align*}
$\therefore\,\, \Phi\left(\frac{C - t/n}{\sqrt{\frac{n - 1}{n}}}\right)$ is the $UMVUE$ of $\tau(\theta) = P(X\leq C)$.
\end{proof}\vspace*{0.4pt}










\begin{problem}
Let $X_1, \, X_2, \, \cdots\, , \, X_n$ be a random sample for the pdf $$f_{\theta}(x) = \theta\, x^{\theta - 1}\, , \,\hspace{0.3cm} 0 < x < 1, \, \, \theta > 0.$$
Show that the geometric mean $$\left(\prod^n_{i = 1} X_i \right)^{1/n}$$ is a complete sufficient statistic and find the $UMVUE$ of $\theta$.
$$\textit{Hint:} - \log X_i \, \thicksim \, Exp\left(\frac{1}{\theta}\right)$$
\end{problem}


\begin{problem}
Let $X_1, \, X_2, \, \cdots\, , \, X_n$ be a random sample from the $GAM(\alpha, \beta)$ distribution and $\theta = (\alpha , \beta)$. Find the $UMVE$ of $\tau(\theta) = \alpha\beta$.
\end{problem}



\begin{problem}
Let $X\thicksim NBIN(k, \theta)$. Find the $UMVUE$ of $\theta$. $$\textit{Hint:}\hspace{0.3cm}\, \text{Find}\,\,\, \, E\left(\frac{1}{X + K - 1}\right).$$
\end{problem}


\begin{problem}
Let $X_1, \, X_2, \, \cdots\, , \, X_n$ be a random sample from the $N(\theta, 1)$ distribution. Find the $UMVUE$ of $\tau(\theta) = \theta^2$.
\end{problem}




\begin{problem}
Let $X_1, \, X_2, \, \cdots\, , \, X_n$ be a random sample from the $N(0, \theta)$ distribution. Find the $UMVUE$ of $\tau (\theta) = \theta^2$.
\end{problem}



\begin{problem}
Let $X_1, \, X_2, \, \cdots\, , \, X_n$ be a random sample from the $POI(\theta)$ distribution. Find the $UMVUE$ for $\tau(\theta) = \left(1 + \theta\right)\, e^{-\theta}$.\\ \, \textit{Hint:}\, \, Find $P(X \leq 1)$.
\end{problem}


\begin{problem}
Let $X_1, \, X_2, \, \cdots\, , \, X_n$ be a random sample from the $POI(\theta)$ distribution. Find the $UMVUE$ for $\tau(\theta) = e^{-2\theta}$.
$$\textit{Hint:}\hspace{0.3cm} E_{\theta}\left[(-1)^{X_1}\right]$$
\end{problem}


\begin{problem}
Let $X_1, \, X_2, \, \cdots\, , \, X_n$ be a random sample from the $GAM(2, \theta)$ distribution. Find the $UMVUE$ of $\tau_1(\theta) = \frac{1}{\theta}$ and the $UMVUE$ of $\tau_2(\theta) = P\left(X_1 > c\right)$ where $c>0$ is a constant.
\end{problem}





\section{Ancillary Statistics}
Information on the unknown parameter $\theta$ is provided by the sensitivity of the distribution of a statistic to changes in the parameter.\\ \vspace*{0.4pt}
A statistic that has no sensitivity at all in  distribution to the parameter, which would appear to contain little information about $\theta$, is called an ancillary statistics.

\begin{defn}
$U(X)$ is an ancillary statistic if its distribution does not depend on the unknown parameter $\theta$.
\end{defn}

The following theorem allows us to use the properties of completeness and ancillarity to prove independence of two statistics without finding their joint distribution.

\begin{thm}[Basu]
Let $X$ have probability (density) function $\{f_{\theta}(x);\, \theta \in \Omega\}$. Let $T(X)$ be  a complete sufficient statistic. Then $T(X)$ is independent of every ancillary statistic $U(X)$.
\end{thm}



\begin{proof}
We need to show that $P\left(U(X) \in B, \, T(X)\in C\right) = P\left(U(x)\in B\right)\, P\left(T(X)\in C\right)$ for all sets $B, C$ and $\theta \in \Omega$.\\
Let $g(t) = P\left(U(X) \in B / T(X)  = t\right)  - P\left(U(X) \in B\right)$ for all $t\in A$ where $P\left(T\in A\right) = 1$. By sufficiency $P\left(U(X) \in B / T(X) = t\right)$ does not depend on $\theta$ and by ancillarity $P\left(U(X)\in B\right)$ also does not depend on $\theta$. Therefore $g(T)$ is a statistic.\\
Let $$I\left(U(X) \in B\right) = 
\begin{cases}
1 & \text{if}\,\,\,\, U(X) \in B\\
0 & \text{otherwise}\\
\end{cases}$$
then $E\left[I(U(X) \in B)\right] = P\left(U(X)\in B\right)$, $\, E\left[I(U(X) = B) / T = t \right] = P\left(U(X) \in B / T = t\right)$ and $g(t) = E\left[I\left(U(X) \in B / T  = t\right)\right] - E\left[I\left(U(X) \in B\right)\right]$. Then 
\begin{align*}
E\left[g(t)\right] & = E\left\{E\left(I(U(X) \in B / T = t\right)\right\} - E\left[I\left(U(X) \in B\right)\right]\\
& = E\left[E\left(U(X) \in B\right)\right] - E\left[I\left(U(X) \in B\right)\right]\\
& = 0.
\end{align*}
Since $T$ is complete this implies that $P_{\theta}\left[g(T) = 0\right] = 1$ for all $\theta \in \Omega$. Therefore\\ $P\left(U(X)\in B / T(X)  = t\right) = P\left(U(X) \in B\right)$ for all $t \in A$ and $B$.\\
Suppose $T$ has probability density function $g_{\theta}(t)$, then 
\begin{align*}
P\left(U(X)\in B, T(X) \in C\right) & = P\left(U(X)\in B / T(X) = t\right)\, P\left(T(X) \in C\right)\\
& =  \int_C P\left(U(X) \in B / T = t\right)\, g_{\theta} (t) \, dt\\
& = \int_CP\left(U(X) \in B\right)\, g_{\theta}(t)\, dt\\
& = P\left(U(X)\in B\right)\, P\left(T(X)\in C\right)
\end{align*}
for all sets $B, C$ and $\theta \in \Omega$ as required.
\end{proof}


\begin{exa}
Let $X_1, \, \cdots\, , \, X_n$ be a random sample from the $EXP(\theta)$ distribution. Show that $T(X) = \sum^n_{i = 1} X_i$ and $U(X) = \left(\frac{X_1}{T}, \, \cdots\, , \, \frac{X_n}{T}\right)$ are independent random variables. Find $E(X_1/T)$.
\end{exa}

\begin{proof}[\textbf{Solution.}]
$X_i \, \thicksim\, EXP(\theta)$
$$T = \sum^n_{i = 1} X_i \, \thicksim\, GAM(n,\theta)$$
$$Z_i = \frac{X_i}{\theta} \, \thicksim\, EXP(1)\,\,\hspace{0.2cm}\text{is ancillary}.$$
$$Z = \frac{T}{\theta} \, \thicksim \, GAM(n ,1)\, \, \text{is ancillary now}$$
$$\frac{X_i}{T} = \frac{X_i/\theta}{T/\theta} = \frac{Z_i}{Z}\hspace{0.3cm}\text{is ancillary}$$
$$\therefore\, \hspace{0.2cm} U(X) = \left(\frac{X_1}{T}, \, \cdots\, , \, \frac{X_n}{T}\right) = \left(\frac{Z_1}{Z}, \, \cdots\, , \, \frac{Z_n}{Z}\right)\hspace{0.3cm} \text{is ancillary statistic}$$
$X_i\thicksim EXP(\theta)$ is a REF distribution.\\
$\therefore\, T = \sum^n_{i = 1} X_i\,$ is a complete sufficient statistic. Therefore by Basus' theorem $U(X)$ and $T$ are independent.\\
Now $\, E(X_1) = \theta$
$$E\left(\frac{X_1}{T}\, \cdot\, T\right) = \theta$$
$$E\left(\frac{X_1}{T}\right)\, E(T) = \theta$$
$$E\left(\frac{X_1}{T}\right)\, n\theta = \theta$$
$$\implies\hspace{0.3cm} E\left(\frac{X_1}{T}\right) = \frac{1}{n}.$$
\end{proof}



\begin{exa}
Let $X_1, \, \cdots\, , \, X_n$ be a random sample from the $N(\mu, \sigma^2)$ distribution. Prove that $\overline{X}$ and $S^2$ are independent.
\end{exa}

\begin{proof}[\textbf{Solution.}]
$X_i \thicksim N(\mu, \sigma^2)$ assume $\sigma^2$ is known, then $\overline{X}$ is a complete sufficient statistic.
$$Z_i = X_i - \mu \, \thicksim \, N(0, \sigma^2)\,\hspace{0.2cm}\text{is ancillary}$$
$\overline{Z} = \overline{X} - \mu \, \thicksim\, N(0, \sigma^2/n)$ is ancillary then $\sum^n_{i = 1} (Z_i - \overline{Z})^2$ is ancillary
$$\sum^n_{i = 1} (Z_i - \overline{Z})^2 = \sum^n_{i = 1} (X_i - \mu  - (\overline{X} - \mu)^2) = \sum^n_{i = 1} (X_i - \overline{X})^2\,\hspace{0.2cm} \text{is ancillary}.$$
Therefore $\, S^2  = \frac{1}{n - 1}\sum^n_{i = 1}(X_i  - \overline{X})^2$ is ancillary.\\ 
Therefore by Basus' theorem $\overline{X}$ and $S^2$ are independent.  Since this argument holds for all values of $\sigma^2$.\\
$\therefore\, \overline{X}$ and $S^2$ are independent for all $\mu$ and $\sigma^2$.
\end{proof}\vspace*{0.4pt}



\begin{problem}
Let $X_1, \, \cdots\, , \, X_n$ be a random sample from the $GAM(\alpha, \beta)$ distribution. Show that $T(X) = \sum^n_{i = 1} X_i$ and $U(X) = \left(\frac{X_1}{T}, \, \cdots\, , \, \frac{X_n}{T}\right)$ are independent random variables. Find $E\left(\frac{X_1}{T}\right)$.
\end{problem}


\textbf{Reference:} Dennis D. Boss and Jacqueline M. Hughes - Oliver (19998).\textit{The American Statistician}, 52(3), 218-221.
%---------------------------













\newpage
\chapter{MAXIMUM LIKELIHOOD ESTIMATION}
\section{Introduction}
Let $X_1, \, \cdots\, , \, X_n$ be a random sample from the model $\{f_{\theta}(x):\, \theta \in \Omega\}$ then
\begin{enumerate}
\item the likelihood function of $X_1, \, \cdots\, , \, X_n$ is given by $$L(\theta)  = L(\theta, x) = \prod^n_{i = 1} f_{\theta}(x_i)$$
\item the log likelihood function of $X_1, \, \cdots\, , \, X_n$ is give by $$\mathcal{L}(\theta) = \mathcal{L}(\theta,x) = \log L(\theta, x)$$
\item the score function is given by $$S(\theta) = S(\theta, x) = \frac{\partial}{\partial\theta}\mathcal{L}(\theta,x)$$
\item the function $$I(\theta)  = I(\theta,x) = \frac{-\partial}{\partial\theta}S(\theta,x) = \frac{-\partial^2}{\partial\theta^2}\mathcal{L}(\theta,x)$$
is called the information function of $X_1, \, \cdots\, , \, X_n$.
\item the function $J(\theta) = E_{\theta}\left[I(\theta, X)\right]$ is called the Fisher or expected information function
\item a value $\widehat{\theta} = \widehat{\theta}(x)$ which maximizes the likelihood function $L(\theta)$ is called a maximum likelihood estimate (ML-estimate) for $\theta$. $\, \widehat{\theta}  = \widehat{\theta}(x)$ is called a maximum likelihood estimator (ML-estimator) for $\theta$.
\end{enumerate}



\textbf{Note:} In many cases, $\widehat{\theta}$ is found by solving the score equation $S(\theta)  = 0$ and checking that $I(\widehat{\theta}) > 0$.


\begin{exa}
If $X_1, \, \cdots\, , \, X_n$ is a random sample from the model $\{f_{\theta}(x):\, \theta\in \Omega\}$. Show that both the Score function and information function are sums of $n$ terms of identical form.
\end{exa}

\begin{proof}[\textbf{Solution.}]
$$L(\theta) = \prod^n_{i = 1}f_{\theta}(x_i)$$
$$\mathcal{L}(\theta) =  \log \prod^n_{i = 1} f_{\theta}(x_i) = \sum^n_{i = 1}  \log f_{\theta}(x_i)$$
\begin{align*}
S(\theta) & = \frac{\partial}{\partial\theta}\mathcal{L}(\theta)
 = \frac{\partial}{\partial\theta}\sum^n_{i = 1} \log f_{\theta}(x_i)
 = \sum^n_{i=1} \frac{\partial}{\partial\theta}\log f_{\theta}(x_i)
 = \sum^n_{i = 1} \frac{\frac{\partial}{\partial\theta} \, f_{\theta} (x_i)}{f_{\theta}(x_i)}
\end{align*}
a sum of $n$ terms of identical form.
\begin{align*}
I(\theta) &   = \frac{-\partial}{\partial\theta}\, S(\theta)\\
& =\frac{-\partial}{\partial\theta}\sum^n_{i = 1} \frac{\frac{\partial}{\partial\theta} \, f_{\theta}(x_i)}{f_{\theta}(x_i)}\\
& = -\sum^n_{i  = 1} \frac{f_{\theta}(x_i)\, \frac{\partial^2}{\partial\theta^2}f_{\theta}(x_i) \, - \, \frac{\partial}{\partial\theta}f_{\theta}(x_i)\, \frac{\partial}{\partial\theta}f_{\theta}(x_i)}{\left(f_{\theta}(x_i)\right)^2}\\
& = \sum^n_{i = 1}\frac{\left(\frac{\partial}{\partial\theta}f_{\theta}(x_i)\right)^2\, - \, f_{\theta}(x_i)\, \frac{\partial^2}{\partial\theta^2}\, f_{\theta}(x_i)}{\left(f_{\theta}(x_i)\right)^2}\\
&  = \sum^n_{i = 1}\left[\left(\frac{\frac{\partial}{\partial\theta}f_{\theta}(x_i)}{f_{\theta}(x_i)}\right)^2 \, - \, \frac{\frac{\partial^2}{\partial\theta^2}f_{\theta}(x_i)}{f_{\theta}(x_i)}\right]
\end{align*}
a sum of $n$ terms of identical form.
\end{proof}


\begin{exa}
Let $X_1, \, \cdots\, , \, X_n$ be a random sample from the $POI(\theta)$ distribution. Find the Fisher information and compare it to the variance of the M.L estimator $\widehat{\theta}$.
\end{exa}

\begin{proof}[\textbf{Solution.}]
\begin{align*}
f_{\theta}(x) & = \frac{e^{-\theta}\, \theta^x}{x!}
 = \prod^n_{i = 1}\frac{e^{-\theta}\, \theta^{x_i}}{x_i!} = \frac{e^{-n\theta}\, \theta^{\sum^n_{i = 1} x_i}}{\prod^n_{i = 1} x_i!}.
\end{align*}
$$\mathcal{L}(\theta) = \log L(\theta) = -n\theta + \sum^n_{i = 1} x_i - \log\prod^n_{i = 1} x_i!$$
$$S(\theta) = \frac{\partial \mathcal{L}}{\partial\theta} = -n + \frac{\sum^n_{i = 1} x_i}{\theta}$$
$S(\theta) = 0$
$$-n + \frac{\sum^n_{i = 1} x_i}{\theta} = 0$$
$$\implies\hspace{0.5cm} n = \frac{\sum^n_{i  = 1} x_i}{\theta}\hspace{0.5cm}\implies\hspace{0.5cm} \theta = \frac{1}{n} \sum^n_{i = 1} x_i$$
$$\widehat{\theta} = \overline{X}.$$

$$I(\theta) = -\frac{\partial^2\mathcal{L}}{\partial\theta^2}  = -\left(\frac{-\sum^n_{i =1} x_i}{\theta^2}\right) = \frac{\sum^n_{i = 1} x_i}{\theta^2}.$$
\begin{align*}
J(\theta) & = E\left[I(\theta, X)\right]  = E\left[\frac{1}{\theta^2}\sum^n_{i = 1} X_i\right]
 = \frac{1}{\theta^2}\, \sum^n_{i = 1} E(X_i)
 = \frac{1}{\theta^2}\, \sum^n_{i = 1}\theta
 = \frac{n\theta}{\theta^2}
 =\frac{n}{\theta}.
\end{align*}

$$var(\widehat{\theta}) = var(\overline{X}) = \frac{\theta}{n}.$$

$$\text{i.e}\hspace{0.3cm} var(\widehat{\theta}) = \frac{1}{J(\theta)} = \left[J(\theta)\right]^{-1}.$$
\end{proof}

\textbf{Note:} If $X_1, \, \cdots\, , \, X_n$ is a random sample from $f_{\theta}(x)$ then 
\begin{align*}
J(\theta) & = E\left[-\frac{\partial^2}{\partial\theta^2}\, \sum^n_{i = 1} \log\, f_{\theta}(X_i)\right]
  = \sum^n_{i  = 1} E\left[\frac{-\partial^2}{\partial\theta^2} \, \log f_{\theta}(X_i)\right]
 = n\, E\left[\frac{-\partial^2}{\partial\theta^2}\, \log f_{\theta}(X)\right].
\end{align*}

E.G: In the above example
$$L_1(\theta) = \frac{e^{-\theta}\, \theta^x}{x!}$$

$$\mathcal{L}_1 (\theta) = -\theta + x\, \log x - \log x!$$

$$S_1(\theta) = \frac{\partial\mathcal{L}_1}{\partial\theta} = -1 + \frac{x}{\theta}.$$

$$I_1(\theta) = \frac{-\partial^2\mathcal{L}_1}{\partial\theta^2} = -\left(\frac{-x}{\theta^2}\right) = \frac{x}{\theta^2}.$$
\begin{align*}
J_1(\theta) & =  E\left[I_1(\theta, X)\right) = E\left(\frac{X}{\theta^2}\right)
 =\frac{E(X)}{\theta^2}
 = \frac{\theta}{\theta^2}
 =\frac{1}{\theta}.
\end{align*}

$J_1(\theta)  = 1/\theta$
$$\therefore\hspace{0.5cm} J(\theta)  = n\, J_1(\theta) = \frac{n}{\theta}\hspace{0.3cm}\text{Ans}.$$\vspace*{0.2pt}



\begin{problem}
Suppose that $X\thicksim BIN(n,\theta)$ and we observe $X$. Find the M.L estimator of $\theta$, the score function and the Fisher information.
\end{problem}

\begin{problem}
Suppose that $X\thicksim NB(k,\theta)$ and we observe $X$. Find the M.L estimator of $\theta$, the Score function and the Fisher information.
\end{problem}

\begin{exa}
Let $X_1, \, \cdots\, , \, X_n$ be a random sample from the $UNIF(0,\theta)$ distribution. Find the M.L estimator of $\theta$.
\end{exa}

\begin{proof}[\textbf{Solution.}]
$f_{\theta}(x) = \frac{1}{\theta}\hspace{0.3cm}, \hspace{0.3cm} 0 < x < \theta$
\begin{align*}
L(\theta) & = \frac{1}{\theta^n}\hspace{0.3cm} , \hspace{0.3cm} 0< x_{(1)}\, , \, \cdots\, ,\, x_{(n)} < \theta\\
&  =  \frac{1}{\theta^n}\hspace{0.3cm}, \hspace{0.3cm}0 < x_{(1)} < x_{(n)} < \theta\\
& = \frac{1}{\theta^n}\, I(X_{(1)}> 0)\, I(X_{(n)} < \theta)
\end{align*}
\begin{figure}[H]
\centering
\begin{tikzpicture}
\draw[-Stealth](-0.6,0)--(4,0)node[right]{$\theta$};
\draw[-Stealth](0,-1)--(0,2.5)node[above]{$L(\theta)$};
\draw[smooth, domain = 0.5:3.5, variable = \x]plot(\x, {1/\x});
\draw(0.5,2)node[shape = circle, scale = 0.01cm, fill = black]{}
(0.5,-0.08)--(0.5,0.08)
(0.5,-0.08)node[below]{$x_{(n)}$};
%\draw(1,1.4)parabola[bend at end](3.5,.02);
\end{tikzpicture}
\end{figure}
Therefore $\, \widehat{\theta} = X_{(n)}$.
\end{proof}




\begin{problem}
Let $X_1, \, \cdots\, , \, X_n$ be a random sample from a distribution with pdf $$f_{\theta}(x) = \theta\, x^{\theta - 1}\hspace{0.3cm} , \hspace{0.3cm} 0 < x < 1\hspace{0.2cm}, \hspace{0.3cm} \theta > 0.$$
Find the M.L estimator of $\theta$,  the scored function and the Fisher information.
\end{problem}




\begin{problem}
Suppose $X_1, \, \cdots\, , \, X_n$  is a random sample from the $UNIF(\theta, \theta + 1)$ distribution. Show that the M.L estimator of $\theta$  is not unique.
\end{problem}



\begin{thm}[INVARIANCE PROPERTY]
If $\widehat{\theta}$ is the M.L estimator of $\theta$ and $\tau(\theta)$ is a function of $\theta$, then $\tau(\widehat{\theta})$ is the M.L estimator of $\tau(\theta)$.
\end{thm}


\begin{exa}
Let $X_1, \, \cdots\, , \, X_n$  be a random sample from the $PIO(\theta)$ distribution. Find the M.L estimator of $\tau(\theta)  = e^{-\theta}$.
\end{exa}

\begin{proof}[\textbf{Solution.}]
From example 2.1.2, $\widehat{\theta} = \overline{X}$.\\
By theorem 2.1.8 $\widehat{\tau} = e^{-\overline{X}}\, $ is the M.L estimator of $\tau(\theta) = e^{-\theta}$.
\end{proof}\vspace*{0.2pt}



\section{Properties of the Score and Information}
Consider a family of probability density functions $\{f_{\theta}(x):\, \theta\in \Omega\}$. Let $A = \{x:\, f_{\theta}(x) > 0\}$. Then $\int_Af_{\theta}(x) \, dx = 1$ and therefore $$\int_A\frac{\partial}{\partial\theta}\, f_{\theta}(x)\, dx = \frac{\partial}{\partial\theta}\int_Af_{\theta}(x)\, dx = 0$$
provided that the integral can be interchanged with the derivative. Models that permit this interchange and calculations of the Fisher information are called regular models.




\begin{thm}
If $X = (X_1, \, \cdots\, , \, X_n)$ is a random sample from a regular model $\{f_{\theta}(x):\, \theta\in \Omega\}$, then 
\begin{enumerate}
\item[(i)] $E_{\theta}\left[S(\theta, X)\right] = 0$
\item[(ii)] $var_{\theta}\left[S(\theta, X)\right] = E_{\theta}\left[\left(S(\theta,X)\right)^2\right] = J(\theta)$.
\end{enumerate}
\end{thm}\vspace*{0.1pt}

\begin{proof}[\textbf{Solution.}]
\begin{enumerate}
\item[(i)] $E_{\theta}\left[S(\theta,X)\right] = E_{\theta}\left[\displaystyle{\sum^n_{i=1}}\dfrac{\partial}{\partial \theta} \, \log \, f_{\theta}(X_i)\right] = \displaystyle{\sum^n_{i= 1} E_{\theta} \left[\frac{\partial}{\partial\theta} \, \log\, f_{\theta}(X_i)\right]}$.
Now 
\begin{align*}
E\left[\frac{\partial}{\partial\theta}\, \log\, f_{\theta}(X_i)\right] & = \int_A\frac{\partial}{\partial\theta}\, \log\, f_{\theta}(x)\, f_{\theta}(x)\, dx\\
& = \int_A\frac{\frac{\partial}{\partial\theta}\, f_{\theta}(x)}{f_{\theta(x)}}\, \cdot\, f_{\theta}(x)\, dx\\
& = \int_A\frac{\partial}{\partial\theta} f_{\theta}(x)\, dx\\
& = \frac{\partial}{\partial\theta}\, \int_Af_{\theta}(x)\, dx\\
& = \frac{\partial}{\partial\theta}\, A(1)\\
& = 0.
\end{align*}

\item[(ii)]$\begin{aligned}[t]
var\left[S(\theta,X)\right]  = E\left[\left(S(\theta, X)\right)^2\right] - \left(E\left[S(\theta,X)\right]\right)^2
&   = E\left[\left(S(\theta,X)\right)^2\right] - 0^2\\
& = E\left[\left(S(\theta,X)\right)^2\right]\hspace{1cm}\cdots\cdots\hspace{0.5cm} (*)
\end{aligned}$\\

$E\left[\left(S(\theta,X)\right)^2\right] = E\left\{\left[\sum^n_{i = 1}S_1(\theta,X_i)\right]^2\right\}\,$ but 
\end{enumerate}
\begin{align*}
S_1(\theta, X_i)  = & \frac{\partial}{\partial\theta}\, \log\, f_{\theta}(X_i)\\
 = &  E\left\{\sum^n_{i = 1}\left(S_1(\theta,X_i)\right)^2 \, + \, \sum_{i\neq j}S_1(\theta,X_i)\, S_1(\theta,X_j)\right\}\\
 = &  \sum^n_{i = 1} E\left(S_1(\theta,X_i)\right)^2 \, + \, \sum_{i\neq j}\underbrace{E\left(S_1(\theta,X_i)\right)}_0\,\underbrace{E\left(S_1(\theta,X_j)\right)}_0\\
& \text{Since}\, X_i \, \text{and}\, X_j\, \text{are independent}\\
& = \sum^n_{i = 1}E\left[\frac{\frac{\partial}{\partial\theta}\, f_{\theta}(X_i)}{f_{\theta}(X_i)}\right]^2.
\end{align*}
Recall Example 2.1.1
$$I(\theta,X) = \sum^n_{i = 1} \left[\left(\frac{\frac{\partial}{\partial\theta}\, f_{\theta}(X_i)}{f_{\theta}(X_i)}\right)^2 \, - \, \frac{\frac{\partial^2}{\partial\theta^2}\, f_{\theta}(X_i)}{f_{\theta}(X_i)}\right]$$
The Fisher information
$$J(\theta,X)   = \sum^n_{i = 1} E\left[\left(\frac{\frac{\partial}{\partial\theta}\, f_{\theta}(X_i)}{f_{\theta}(X_i)}\right)^2 \right]\, - \,\sum^n_{i = 1} E\left[ \frac{\frac{\partial^2}{\partial\theta^2}\, f_{\theta}(X_i)}{f_{\theta}(X_i)}\right]$$
but
\begin{align*}
E\left(\frac{\frac{\partial^2}{\partial\theta^2}\, f_{\theta}(X_i)}{f_{\theta}(X_i)}\right) &   = \int_A\frac{\frac{\partial^2}{\partial\theta^2}\, f_{\theta}(x)}{f_{\theta}(x)}\, \cdot\, f_{\theta}(x)\, dx\\
&  = \int_A \frac{\partial^2}{\partial\theta^2}\, f_{\theta}(x)\, dx\\
& = \frac{\partial^2}{\partial\theta^2}\, \int_Af_{\theta}(x)\, dx\\
& = \frac{\partial^2}{\partial\theta^2}\, (1)
 = 0.
\end{align*}
$$\therefore\,\,\, J(\theta) = \sum^n_{i = 1} E\left[\left(\frac{\frac{\partial}{\partial\theta}\, f_{\theta}(X_i)}{f_{\theta}(X_i)}\right)^2\right].$$

$$\therefore\,\,\, var_{\theta}\left(S(\theta,X)\right) = E_{\theta}\left[\left(S(\theta,X)\right)^2\right] = J(\theta).$$
\end{proof}




\begin{exa}
Let $X = (X_1, \, \cdots\, , \, X_n)$ be a random sample from the $POI(\theta)$ distribution, then
$$L(\theta) = \prod^n_{i = 1}\frac{e^{-\theta}\, \theta^{x_i}}{x_i!} = \frac{e^{-n\theta}\, \theta^{\sum^n_{i  = 1} x_i}}{\prod^n_{i = 1} x_i!}.$$

$$\mathcal{L}(\theta) = \log\, L(\theta) = -n\theta + \sum^n_{i = 1} x_i\, \log \theta  - \log\prod^n_{i = 1} x_i!.$$

$$S(\theta) = \frac{\partial\mathcal{L}}{\partial\theta} = -n + \frac{\sum^n_{i = 1} x_i}{\theta}.$$

$$I(\theta) = \frac{-\partial^2\mathcal{L}}{\partial\theta^2} = -\left(\frac{-\sum^n_{i = 1} x_i}{\theta^2}\right) = \frac{\sum^n_{i = 1} x_i}{\theta^2}.$$

\begin{align*}
J(\theta) & = E\left[I(\theta, X)\right] = E\left[\frac{1}{\theta^2} \, \sum^n_{i = 1} X_i\right]\\
& =  \frac{1}{\theta^2} \, \sum^n_{i = 1} E(X_i)\\
& = \frac{1}{\theta^2}\, \cdot\, n\theta\\
 & = \frac{n}{\theta}.
\end{align*}
Now
\begin{align*}
E\left[S(\theta,X)\right] & = E\left[-n + \frac{1}{\theta}\, \sum^n_{i = 1} X_i\right]
 = -n  + \frac{1}{\theta}\, \sum^n_{i = 1} E(X_i)\\
& = -n  + \frac{1}{\theta}\cdot n\theta\\
& =0.
\end{align*}
\begin{align*}
var\left[S(\theta,X)\right] & = var\left[-n + \frac{1}{\theta}\, \sum^n_{i = 1} X_i\right]
 = \frac{1}{\theta^2}\, \sum^n_{i = 1} var(X_i)\\
& = \frac{1}{\theta^2}\, \cdot\, n\theta\\
& = \frac{n}{\theta}\\
& = J(\theta).
\end{align*}
\end{exa}



\begin{multiparameter}
Suppose $\theta = (\theta_1, \, \cdots\, , \, \theta_k)^t$, then the score function $S(\theta)$ is a $k$-dimensional vector whose $i^{\text{th}}$ component is the derivative of $\mathcal{L}(\theta)$ with respect to the $i^{\text{th}}$ component of $\theta$ $$S(\theta) = \left[\frac{\partial\mathcal{L}(\theta)}{\partial\theta_1}\, , \, \cdots\, , \, \frac{\partial \mathcal{L}(\theta)}{\partial\theta_k}\right]^t.$$
The information function $I(\theta)$ is a $k\times k$ matrix whose $(i,j)^{\text{th}}$ element is minus the derivative of $\mathcal{L}(\theta)$ with respect  to the $i^{\text{th}}$ and $j^{\text{th}}$ components of $\theta$, 
$$I(\theta) = \left[I_{ij}(\theta)\right]_{k\times k} = \left[\frac{-\partial^2\mathcal{L}(\theta)}{\partial\theta_i\, \partial\theta_j}\right]_{k\times k} \hspace{0.2cm} , \hspace{0.5cm} i,j = 1, \, 2, \, \cdots\, , \, k$$
The Fisher information is a $k\times k$ matrix whose elements are componentwise expected of the information matrix
$$J_{ij}(\theta) = E\left[I_{ij}(\theta,X)\right]\hspace{0.2cm}  , \hspace{0.3cm} i,j = 1, \, 2, \, \cdots\, , \, k$$
For a regular family of distributions $E\left[S(\theta, X)\right] = (0, \, \cdots\, , \, 0)^t$ and $$J(\theta) = var\left[S(\theta,X)\right] = E\left[S(\theta, X)\, \left[S(\theta,X)\right]^t\right]$$
\textbf{Note:} The invariance property of the M.L estimator also holds in the multiparameter case.
\end{multiparameter}\vspace*{0.1pt}



\begin{exa}
Let $X_1, \, \cdots\, , \, X_n$ be a random sample from the $N(\mu, \sigma^2)$ distribution. Find the ML  estimator of $\theta = (\mu, \sigma^2)^t$ the score function and the Fisher information.
\end{exa}


\begin{proof}[\textbf{Solution.}]
Let $\, \theta_1  = \mu, \, \, \theta_2 = \sigma^2, \, \, \theta = (\theta_1, \theta_2)^{t}$
$$L(\theta)  = \prod^n_{i = 1}\frac{1}{\sqrt{2\pi\theta^2}}\exp\left\{-\frac{1}{2\theta_2}(x_i  - \theta_1)^2\right\}  = \frac{1}{\left(2\pi\theta_2\right)^{n/2}}\, \exp\left\{-\frac{1}{2\theta_2}\, \sum^n_{i = 1} (x_i - \theta_1)^2\right\}$$

$$\mathcal{L}(\theta) = \frac{-1}{2\theta_2} \, \sum^n_{i = 1} (x_i - \theta)^2 - \frac{n}{2}\, \log\, 2\pi \theta_2$$

$$S(\theta) = \begin{pmatrix}
\dfrac{\partial \mathcal{L}}{\partial \theta_1}\\\\
\dfrac{\partial \mathcal{L}}{\partial\theta_2}\\
\end{pmatrix}$$

$$\frac{\partial \mathcal{L}}{\partial\theta_1} = \frac{-1}{2\theta_2}\sum^n_{i = 1} 2(x_i - \theta_1)\, (-1) = \frac{1}{2\theta_2}\sum^n_{i = 1} 2(x_i - \theta_1)$$

$$\frac{\partial\mathcal{L}}{\partial\theta_2} = \frac{1}{2\theta^2_2} \sum^n_{i = 1}(x_i  - \theta_1)^2 - \frac{n}{2}\cdot\frac{2\pi}{2\pi\,\theta_2} = \frac{1}{2\theta^2_2}\sum^n_{i = 1} (x_i - \theta_1)^2 - \frac{n}{2\theta_2}.$$

$$S(\theta)  = \begin{pmatrix}
\dfrac{1}{\theta_2}\sum^n_{i = 1}(x_i - \theta_1)\\\\
\dfrac{1}{2\theta_2^2}\sum^n_{i = 1} (x_i - \theta_1)^2 - \dfrac{n}{2\theta_2}\\
\end{pmatrix}.$$

$S(\theta) = 0$

$$\frac{1}{\theta^2}\, \sum^n_{i = 1} (x_i - \theta_1) = 0$$

$$\widehat{\theta}_1 = \frac{1}{n}\sum^n_{i =1} X_i = \overline{X}.$$

$$\frac{1}{2\theta^2_2} \sum^n_{i = 1} (X_i - \theta_1)^2 - \frac{n}{2\theta_2} = 0.$$

$$\frac{1}{n}\sum^n_{i = 1} (X_i - \overline{X})^2 = \widehat{\theta}_2.$$

$\therefore\, \widehat{\theta} = \left(\overline{X}\, \, , \, \, \frac{1}{n}\sum^n_{i = 1}(X_i - \overline{X})^2\right)^t\,$ is the ML estimator of $\theta = (\mu , \sigma^2)^t.$

$$I(\theta) = \begin{pmatrix}
-\dfrac{\partial^2\mathcal{L}}{\partial\theta_1^2} & - \dfrac{\partial^2\mathcal{L}}{\partial\theta_1\, \partial\theta_2}\\\\
- \dfrac{\partial^2\mathcal{L}}{\partial\theta_2\, \partial\theta_1} & -\dfrac{\partial^2\mathcal{L}}{\partial\theta_2^2}\\
\end{pmatrix}.$$

$$-\frac{\partial^2\mathcal{L}}{\partial\theta_1^2} = -\frac{1}{\theta_2}\, \sum^n_{i = 1}(-1) = \frac{n}{\theta_2}.$$

$$-\frac{\partial^2\mathcal{L}}{\partial\theta_1\, \partial\theta_2} = -\left(-\frac{1}{\theta^2_2}\sum^n_{i = 1}(x_i - \theta)\right) = \frac{1}{\theta^2_2}\sum^n_{i = 1} (x_i - \theta_1).$$

$$-\frac{\partial^2\mathcal{L}}{\partial\theta^2_2} = - \left(-\frac{1}{\theta^3_2} \sum^n_{i = 1} (x_i - \theta_1)^2 + \frac{n}{2\theta^2_2}\right) = \frac{1}{\theta^3_2} \sum^n_{i = 1} (x_i - \theta_1)^2 - \frac{n}{2\theta^2_2}.$$

$$I(\theta) = \begin{pmatrix}
\dfrac{n}{\theta_2} && \dfrac{1}{\theta^2_2}\sum^n_{i = 1}(X_i - \theta_1)\\\\
\dfrac{1}{\theta^2_2}\sum^n_{i = 1}(X_i - \theta_1) &&  \dfrac{1}{\theta^3_2}\sum^n_{i = 1} (X_i - \theta_1)^2 - \dfrac{n}{2\theta^2_2}\\
\end{pmatrix}.$$

$$J(\theta)  = E\left[I(\theta)\right] = \begin{pmatrix}
\dfrac{n}{\theta_2} && \dfrac{1}{\theta^2_2}\sum^n_{i = 1}(E(X_i) - \theta_1)\\\\
\dfrac{1}{\theta^2_2}\sum^n_{i = 1}(E(X_i) - \theta_1) &&  \dfrac{1}{\theta^3_2}\sum^n_{i = 1} E(X_i - \theta_1)^2 - \dfrac{n}{2\theta^2_2}\\
\end{pmatrix}.$$


$$\sum^n_{i = 1} \left(E(X_i)  - \theta_1\right) = \sum^n_{i = 1} (\theta_1  - \theta_1) = 0$$
\begin{align*}
\frac{1}{\theta_2^3}\sum^n_{i = 1} E\left(X_i - \theta_1\right)^2 - \frac{n}{2\theta^2_2} &  = \sum^n_{i = 1} var(X_i)  - \frac{n}{2\theta^2_2}
 = \frac{1}{\theta^3_2}\, n\theta_2 - \frac{n}{2\theta_2^2}\\
& = \frac{n}{\theta^2_2} - \frac{n}{2\theta_2^2}\\
& =  \frac{n}{2\theta_2^2}.
\end{align*}

$$\implies\hspace{0.5cm} J(\theta) = \begin{pmatrix}
\dfrac{n}{\theta_2^2} & 0\\\\
0 & \dfrac{n}{2\theta_2^2}\\
\end{pmatrix} = \begin{pmatrix}
\dfrac{n}{\sigma^2} & 0 \\\\
0 & \dfrac{n}{2\sigma^4}\\
\end{pmatrix}.$$
i.e $\hspace{0.3cm} \left(\overline{X}\, , \, \frac{1}{n}\sum^n_{i = 1}(X_i - \overline{X})^2\right)$
$$\left(\overline{X}\, ,\, \dfrac{n - 1}{n}\, S^2\right).$$
\end{proof}



\begin{problem}
Let $(X_1, \, X_2, \, X_3) \thicksim MULT(n, \, \theta_1, \, \theta_2, \, \theta_3)$. Find the M.L estimator of $\theta = (\theta_1, \theta_2)^t$, the score function and the Fisher information.
\textit{Hint:}\, use 
$$f_{\theta}(x_1, x_2, x_3)  = \frac{n!}{x_1!\, x_2!\, (n - x_1 - x_2)!}\, \theta_1^{x_1}\, \theta_2^{x_2}\, \left(1 - \theta_1 - \theta_2\right)^{n - x_1 - x_2}$$
\end{problem}




\begin{newtons}
Suppose that the M.L estimate is determined by the likelihood equation $S(\theta) = 0$. \\
It frequently happens that an analytic solution for $\theta$ cannot be obtained. If we begin with an approximate value for the parameter, $\theta^{(0)}$, we may update that value as follows:
$$\theta^{(i + 1} = \theta^{(i)} + \frac{S\left(\theta^{(i)}\right)}{I\left(\theta^{(i)}\right)}\hspace{0.3cm} , \hspace{0.5cm} i = 0, \, 1, \, \cdots$$
and provided convergence of $\theta^{(i)}$,  $i\longrightarrow \infty$ obtains, it  converges to the score equation above.\\
In the multiparameter case, where $S(\theta)$ is a vector and $I(\theta)$ is a matrix, then Newton's method becomes
$$\theta^{(i + 1)} = \theta^{(i)} + \left[I\left(\theta^{(i)}\right)\right]^{-1}\, S\left(\theta^{(i)}\right)\, , \hspace{0.2cm} i = 0, \, 1, \, \cdots$$
This algorithm is also called the Newton-Raphson method. A similar algorithm is obtained $I(\theta)$ replaced by the Fisher information $J(\theta)$.
\end{newtons}
\vspace*{4pt}



\textbf{Note:} This is based on Taylor series: i.e
$$f(x) = f(x_0) + \frac{(x - x_0)}{1!}\, f'(x_0) +  \frac{(x - x_0)^2}{2!} \, f''(x_0) + \cdots\cdots$$

$$f(x) \approx f(x_0) + (x - x_0)\, f'(x_0)$$
We solve the equation $f(x)  = 0$
$$f(x_0) + (x - x_0) \, f'(x_0) = 0$$
$$\frac{f(x_0)}{f'(x_0)}  + x -x_0 = 0$$
$$x = x_0 - \frac{f(x_0)}{f'(x_0)}$$
i.e
$$\theta^{(i + 1)} = \theta^{(i)} + \frac{S\left(\theta^{(i)}\right)}{I\left(\theta^{(i)}\right)}.$$
\vspace*{4pt}






\begin{exa}
Suppose $X_1, \, \cdots\, , \, X_n$ is a random sample from the distribution with pdf $$f_{\theta}(x) = \frac{\alpha\, \beta}{\left(1 + \beta \, x\right)^{\alpha + 1}}\, , \hspace{0.3cm} x > 0;\, \, \alpha , \, \beta >0.$$
Explain how you find the M.L estimate of $\theta = (\alpha, \beta)^t$. Find the Fisher information.
\end{exa}



\begin{proof}[\textbf{\textit{Solution.}}]
$$L(\theta) = \prod^n_{i = 1} \theta \, x_i^{\theta - 1}\, e^{-x_i^{\theta}} =  \theta^n\, \prod^n_{i = 1} x_i^{\theta - 1}\, e^{-\sum^n_{i = 1} x_i^{\theta}}$$

$$\mathcal{L}(\theta) = n\, \log \theta  + (\theta - 1)\, \log \prod^n_{i = 1} x_i - \sum^n_{i  =1} x_i^{\theta}.$$


$$S(\theta) = \frac{n}{\theta} + \sum^n_{i = 1}\log x_i - \sum^n_{i = 1} x_i^{\theta}\, \log x_i.$$
$S(\theta) = 0\,\,$ cannot be solved explicitly

$$I(\theta) = -\frac{\partial^2\mathcal{L}}{\partial\theta^2} = -\left(-\frac{n}{\theta^2} - \sum^n_{i = 1} x^{\theta}_i \, \left(\log x_i\right)^2\right) = \frac{n}{\theta^2} + \sum^n_{i = 1} x_i^{\theta}\, \left(\log x_i\right)^2.$$
Solve for $\theta$ iteratively using 
$$\theta^{(i + 1)} = \theta^{(i)} + \frac{S\left(\theta^{(i)}\right)}{I\left(\theta^{(i)}\right)}\hspace{0.2cm} , \hspace{0.6cm} i = 0, \, 1, \, 2, \, \cdots$$
iterate until convergence to obtain $\widehat{\theta}$.
\end{proof}



\begin{exa}
Suppose $X_1, \, \cdots\, , \, X_n$ is a random sample from the $GAM(\alpha, \beta)$ distribution. Explain how you would obtain the $M.L$ estimator of $\theta = (\alpha ,\beta)^t$. Find the Fisher information.
\end{exa}


\begin{proof}[\textbf{\textit{Solution.}}]
$$L(\theta) = \prod^n_{i = 1}\frac{x_i^{\alpha - 1}\, e^{-\frac{x_i}{\beta}}}{\beta^{\alpha}\, \Gamma(\alpha)} = \frac{\left(\prod^n_{i = 1} x_i\right)^{\alpha - 1}\, e^{-\frac{1}{\beta}\sum^n_{i = 1} x_i}}{\beta^{n\alpha}\, \left(\Gamma (\alpha)\right)^n}.$$

$$\mathcal{L}(\theta) = (\alpha - 1)\, \log \prod^n_{i = 1} x_i -  \frac{1}{\beta} \sum^n_{i = 1} x_i - n\alpha \, \log\beta - n\, \log \Gamma(\alpha).$$

$$S(\theta) = \begin{pmatrix}
 \dfrac{\partial \mathcal{L}}{\partial \alpha}\\\\  \dfrac{\partial \mathcal{L}}{\partial \beta}\\
\end{pmatrix}\hspace{2cm}
\begin{matrix}
\dfrac{\partial \mathcal{L}}{\partial \alpha}  =  \sum^n_{i = 1} \log x_i - n\, \log \beta - n\, \Psi(\alpha)\\\\
\dfrac{\partial \mathcal{L}}{\partial \beta}  =  \dfrac{1}{\beta^2}\, \sum^n_{i = 1} x_i - \dfrac{n\alpha}{\beta}\\
\end{matrix}$$
i.e $\, \Psi (z) = \frac{d}{dz}\, \log \Gamma (z) = $ diagamme function.\\

$ S(\theta) = \begin{pmatrix}
0\\ 0\\
\end{pmatrix}\,$ cannot be solved explicitly.

$$I(\theta)= \begin{pmatrix}
-\dfrac{\partial^2\mathcal{L}}{\partial \alpha^2} & -\dfrac{\partial^2\mathcal{L}}{\partial \alpha\, \partial\beta}\\\\
-\dfrac{\partial^2\mathcal{L}}{\partial \beta\, \partial\alpha} & -\dfrac{\partial^2\mathcal{L}}{\partial \beta^2}\\
\end{pmatrix}$$

$-\dfrac{\partial^2\mathcal{L}}{\partial \alpha^2} =  n\, \Psi'(\alpha)$, \hspace{1cm}
$-\dfrac{\partial^2\mathcal{L}}{\partial \alpha\, \partial\beta} = -\dfrac{\partial^2\mathcal{L}}{\partial \beta\, \partial\alpha} = \dfrac{n}{\beta}$\hspace{0.3cm} and\\

$-\dfrac{\partial^2\mathcal{L}}{\partial \beta^2} = -\left(-\dfrac{2}{\beta^3} \sum^n_{i = 1} x_i + \dfrac{n\alpha}{\beta^2}\right)  = \dfrac{2}{\beta^3}\, \sum^n_{i  = 1} x_i - \dfrac{n\alpha}{\beta^2}$.\\




$$I(\theta) = \begin{pmatrix}
n\, \Psi'(\alpha) &  \dfrac{n}{\beta}\\\\
\dfrac{n}{\beta} & \dfrac{2}{\beta^3}\sum^n_{i = 1} x_i - \dfrac{n\alpha}{\beta^2}\\
\end{pmatrix}.$$

Solve for $\theta$ iteratively using $$\theta^{(i + 1)} = \theta^{(i)} + \left[I\left(\theta^{(i)}\right)\right]^{-1}\, S\left(\theta^{(i)}\right)$$
iterate until convergence to obtain $\widehat{\theta} = \begin{pmatrix}
                                                            \widehat{\alpha}\\
                                                            \widehat{\beta}\\
                                                          \end{pmatrix}.$
                                                          
\begin{align*}
J(\theta) = E\left[I(\theta,X)\right] & = \begin{pmatrix}
n\, \Psi'(\alpha) &  \dfrac{n}{\beta}\\\\
\dfrac{n}{\beta} & \dfrac{2}{\beta^3}\sum^n_{i = 1} E(X_i) - \dfrac{n\alpha}{\beta^2}\\
\end{pmatrix} = \begin{pmatrix}
n\, \Psi'(\alpha) &  \dfrac{n}{\beta}\\\\
\dfrac{n}{\beta} & \dfrac{2}{\beta^3}\, n\alpha\beta - \dfrac{n\alpha}{\beta^2}\\
\end{pmatrix}\\\\
& =\begin{pmatrix}
n\, \Psi'(\alpha) &  \dfrac{n}{\beta}\\\\
\dfrac{n}{\beta} &  \dfrac{n\alpha}{\beta^2}\\
\end{pmatrix} = n\,\begin{pmatrix}
 \Psi'(\alpha) &  \dfrac{1}{\beta}\\\\
\dfrac{1}{\beta} &  \dfrac{\alpha}{\beta^2}\\
\end{pmatrix}.
\end{align*}
\end{proof}










\begin{problem}
Suppose $X_1, \, \cdots\, , \, X_n$ is a random sample from the distribution with pdf  $$f_{\theta}(x) = \frac{\alpha\beta}{\left(1 + \beta x\right)^{\alpha + 1}}\hspace{0.4cm}, \hspace{0.4cm} x > 0;\hspace{0.4cm} \alpha, \, \beta > 0$$
Explain how you find the $M.L$ estimate of $\theta = (\alpha, \beta)^t$. Find the Fisher information.
\end{problem}



\section{The Information Inequality}
Suppose we consider estimating a parameter $\tau(\theta)$, where $\theta$ is a scalar, using an unbiased estimator $T(X)$. A lower bound for the variance of $T(X)$ can be found by the information inequality.

\begin{information inequality}
Suppose $T(X)$ is an unbiased estimator of the parameter $\tau(\theta)$ in a regular statistical model $\{f_{\theta}(x):\theta \in \Omega\}$. Then $$var\left(T\right) \geq \frac{\left[\tau'(\theta)\right]^2}{J(\theta)}.$$
Equality holds in and only if $f_{\theta}(x)$ is regular exponential family with natural sufficient statistic $T(X)$.
\end{information inequality}


\textbf{Note:}
\begin{enumerate}
\item If equality holds then $T(X)$ is called an efficient estimator of $\tau(\theta)$.
\item The number $\frac{\left[\tau'(\theta)\right]^2}{J(\theta)}$ is called the Cramer- Rao lower bound (C.R.L.B).
\item The ratio of the C.R.L.B to the variance of an unbiased estimator is called the efficiency of the estimator.
\end{enumerate}

\begin{proof}
$E(T) = \tau(\theta)$

$$\int_A T(X)\, f_{\theta}(x)\, dx  = \tau (\theta)$$

$$\frac{\partial}{\partial \theta}\int_A T(X)\, f_{\theta}(x)\, dx   = \tau'(\theta)$$


$$\int_A T(X)\, \frac{\partial}{\partial \theta}f_{\theta}(x)\, dx  = \tau'(\theta)$$

$$\int_A T(X)\, \frac{\frac{\partial}{\partial \theta}f_{\theta}(x)}{f_{\theta}(x)}\, \cdot\, f_{\theta}(x)\, dx  = \tau'(\theta)$$

$$\int_A T(X)\, S(\theta, X)\, f_{\theta}(x) \, dx = \tau'(\theta)$$
$$E\left[T(X)\, S(\theta,X)\right] = \tau'(\theta)$$
$$cov\left[T(X),S(\theta,X)\right] = E\left[T(X)\, S(\theta, X)\right] - E\left[T(X)\right]\, \underbrace{E\left[S(\theta,X)\right]}_0$$
Therefore $\, Cov\left[T(X), S(\theta, X)\right] = \tau'(\theta)$.\\

By the Cauchy-Schwartz inequality $\, \left[Cov(X,Y)\right]^2 \leq var(X)\, var(Y).$
$$-1 \, \leq \, \frac{Cov(X,Y)}{\sqrt{var(X)\, var(Y)}}\, \leq 1.$$
Therefore
$$\left(Cov\left[T(X),S(\theta,X)\right]\right)^2 \leq var(T(X))\, \cdot \, var\left[S(\theta ,X)\right]\hspace{0.3cm}\cdots\hspace{0.3cm} (2.1)$$
\begin{align*}
\left[\tau'(\theta)\right]^2 & \geq var(T(x))\, J(\theta)\\
var(T(X))\, & \geq \, \frac{\left[\tau'(\theta)\right]^2}{J(\theta)}
\end{align*}
Equality in $(2.1)$ holds iff $T$ and $S(\theta,X)$ are linear functions of each other. i.e
$$S(\theta, x) = A_1(\theta) \, T(x) + A_2(\theta)$$
$$\int S(\theta , x)\, d\theta = \int\left(A_1(\theta)\, T(X) + A_2(\theta)\right)\, d\theta$$
$$\int \frac{\partial}{\partial \theta}\log \, f_{\theta}(x)\, dx  = C_1(\theta)\, T(X)  + C_2(\theta) + \underbrace{C_3(X)}_{\text{constant}}$$
$$\log\, f_{\theta}(x) = C_1(\theta) \, T(x) + C_2(\theta) + C_3(x)$$
\begin{align*}
f_{\theta}(x) & = e^{C_1(\theta) \, T(x) + C_2(\theta) + C_3(x)}
= \underbrace{e^{C_2(\theta)}}_{C(\theta)} \, \exp\Big\{\underbrace{C_1(\theta)}_{q(\theta)}\, \underbrace{T(x)}_{T(X)}\Big\}\, \underbrace{e^{C_3(x)}}_{h(x)}
\end{align*}
which is an exponential family with natural sufficient statistic $T(X)$.
\end{proof}


\textbf{Note:} If $\tau(\theta) = \theta$, then $$var(T) \geq \frac{1}{J(\theta)}.$$



\begin{exa}
Let $X_1, \, X_2, \, \cdots\, , \, X_n$ be a random sample from the $POI(\theta)$ distribution. Show that the $UMVUE$ of $\theta$ achieves the Cramer-Rao lower bound.
\end{exa}

\begin{proof}[\textbf{\textit{Solution.}}]
 $\tau(\theta) = \theta \, \implies\, \tau'(\theta) = 1$.\\
 $T(X) = \sum^n_{i = 1}X_i\,$ is a complete sufficient statistic.  Now $E\left(\overline{X}\right)  = \theta$, therefore $\overline{X}$ is the $U.M.V.U.E$ of $\theta$.
 $$L(\theta) = \frac{e^{-\theta}\, \theta^x}{x!}.$$
 
 $$\mathcal{L}(\theta) = -\theta + x\, \log \theta  - \log x!.$$
 
 $$S(\theta) = -1 + \frac{x}{\theta}.$$
 
 $$I(\theta) = -\left(-\frac{x}{\theta^2}\right) = \frac{x}{\theta^2}.$$
 
 $$J_1(\theta) = E\left[I(\theta,X)\right] = E\left(\frac{X}{\theta^2}\right) = \frac{\theta}{\theta^2} = \frac{1}{\theta}.$$
 $$\therefore\hspace{0.3cm} J(\theta) = \frac{n}{\theta}.$$
 
 $$C.R.L.B = \frac{1}{J(\theta)} = \frac{\theta}{n}.$$
 
 $$var\left(\overline{X}\right) = \frac{var(X)}{n}   = \frac{\theta}{n} = CRLB.$$
 Therefore the $UMVUE$ attains the $CRLB$.\\
 
 OR\\
 $X\thicksim POI(\theta)$ is the $REF$ distribution with natural sufficient statistic $\sum^n_{i  = 1} X_i$. $\overline{X}$ is also a natural sufficient statistic. Therefore $var(\overline{X}) = CRLB$.
\end{proof}





\begin{exa}
Suppose $X_1, \, \cdots\, , \ X_n$ is a random sample from the distribution with pdf $$f_{\theta}(x) = \theta\, x^{\theta - 1}\hspace{0.3cm} \, 0 < x < 1, \hspace{0.3cm} \theta > 0.$$
\begin{enumerate}
     \item[(a)] Show that the variance of the $UMVUE$ of $\theta$ does not achieve the $CRLB$.  What is the efficiency of the $UMVUE$.
     \item[(b)] Find the $CRLB$ for variances of unbiased estimators of $\, \tau(\theta) = \theta (1  - \theta)$.
\end{enumerate}
\end{exa}

\begin{proof}[\textbf{\textit{Solution.}}]
\begin{enumerate}
\item[(a)] $f_{\theta}(x) = \theta \, x^{\theta - 1}\, = \underbrace{\theta}_{C(\theta}\, \exp\Big\{\underbrace{\theta}_{q(\theta)}\, \underbrace{x}_{T(X)}\Big\}\, \underbrace{\frac{1}{x}}_{h(x)}.$\\
Let $\eta = \theta$ $$f_{\eta}(x) = \eta \, \exp\left\{\eta\, \log x\right\}\, \frac{1}{x}.$$
Therefore $f_{\theta}(x)$ is a $REF$ distribution and hence $\sum^n_{i = 1}\log x_i\,$ is a compete sufficient statistic.
$$-\log X_i \thicksim EXP\left(\frac{1}{\theta}\right)$$

$$T = - \sum^n_{i = 1} \log X_i \thicksim GAM\left(n, \frac{1}{\theta}\right).$$
by  problem 1.2.4
$$E\left(X^p\right) = \frac{\beta^p\, \Gamma\left(\alpha + p\right)}{\Gamma(\alpha)}$$

$$E\left(T^{-1}\right) = \frac{\left(\frac{1}{\theta}\right)^{-1}\, \Gamma(n  - 1)}{\Gamma(n)} = \frac{\theta\, \Gamma(n - 1)}{(n - 1)\, \Gamma(n - 1)}  = \frac{\theta}{n - 1}.$$

$$\therefore\, \, E\left[\frac{n - 1}{T} \right] = \theta.$$
Therefore $\,\frac{n - 1}{T}\,$ is the $UMVUE$ of $\theta$.\\

Now $$L_1(\theta) = \theta\, x^{\theta - 1}$$

$$\mathcal{L}_1(\theta) = \log\theta + (\theta - 1)\, \log x$$

$$S_1(\theta) = \frac{1}{\theta} + \log x$$

$$I_1(\theta) = -\left(-\frac{1}{\theta^2}\right) = \frac{1}{\theta^2}$$

$$J_1(\theta) = \frac{1}{\theta^2}$$

$$\therefore\hspace{0.3cm} J(\theta) = n\, J_1(\theta) =  \frac{n}{\theta^2}.$$

$$CRLB = \frac{\left[\tau'(\theta)\right]^2}{J(\theta)} = \frac{1}{J(\theta)} = \frac{\theta^2}{n}.$$

Now $$var\left[\frac{n - 1}{T}\right] = (n - 1)^2\, var\left(\frac{1}{T}\right)$$

$$var\left(\frac{1}{T}\right) = E\left(\frac{1}{T}\right)^2 - \left(\left(\frac{1}{T}\right)\right)^2.$$

$$E\left(\frac{1}{T}\right)^2 = E\left(T^{-2}\right) = \frac{\left(\frac{1}{\theta}\right)^{-2}\, \Gamma(n - 2)}{\Gamma(n)} = \frac{\theta^2\, \Gamma(n - 2)}{(n - 1)\, (n - 2)\, \Gamma(n - 2)} = \frac{\theta^2}{(n - 1)\, (n - 2)}.$$

$$var\left(\frac{1}{T}\right) = \frac{\theta^2}{(n - 1)(n - 2)} - \frac{\theta^2}{(n - 1)^2}.$$
\begin{align*}
\therefore\,\, var\left(\frac{n - 1}{T}\right) & = (n - 1)^2\, \theta^2\, \left[\frac{1}{(n - 1)(n - 2)} - \frac{1}{(n  - 1)^2}\right] = \left(\frac{n - 1}{n - 2} - 1\right)\, \theta^2\\
&  = \left(\frac{n - 1 - n + 2}{n  - 2}\right)\, \theta^2\\
& = \frac{\theta^2}{n - 2} > \frac{\theta^2}{n}.
\end{align*}
Therefore the variance of the $UMVUE$ of $\theta$ does not achieve the $CRLB$.\\


OR\\
$T = -\sum^n_{i = 1} \log X_i$ is a natural sufficient statistic. But $g(T) = \frac{n - 1}{T}$ is not a natural sufficient statistic. Therefore $g(T) = \frac{n - 1}{T}$ does not achieve the $CRLB$.
$$\textit{efficiency} = \frac{CRLB}{var[g(T)]} = \frac{\theta^2/n}{\theta^2/n  - 2}  = \frac{n - 2}{n}.$$
i.e $\, \frac{n - 2}{n}\longrightarrow 1, \, g(T)$ is an asymptotically efficient estimator.


\item[(b)] $\tau(\theta) = \theta(1 - \theta) = \theta - \theta^2$
$$CRLB = \frac{\left[\tau'(\theta)\right]^2}{J(\theta)} = \frac{\left(1 - 2\theta\right)^2}{n/\theta^2} = \frac{\theta^2\, (1 - 2\theta)^2}{n}.$$
\end{enumerate}
\end{proof}



\begin{problem}
For each of the following, determine whether the variance of the $UMVUE$ of $\theta$ based on a random sample $X_1, \, \cdots\, , \, X_n$ achieve the Cramer Rao Lower bound and find the efficiency of the $UMVUE$. For each case find the $CRLB$ for variances of unbiased estimators of $\theta^2$.
\begin{multicols}{4}
\begin{itemize}
\item[(a)] $N(\theta,\mu)$
\item[(b)] $BIN(1,\theta)$
\item[(c)] $N(0,\theta^2)$
\item[(d)] $N(0,\theta)$
\end{itemize}
\end{multicols}
\end{problem}
\vspace*{4pt}



\begin{the multiparameter case} If $\theta = (\theta_1, \, \cdots\, , \, \theta_k)^t$ is a vector, $\tau(\theta)$ is a function of $\theta$ and $T(X)$ is any unbiased estimator of $\tau(\theta)$ in a regular model, then the information inequality becomes  $$var_{\theta}(T)  \geq D(\theta)^t\, J^{-1}(\theta)\, D(\theta)$$ for all $\theta \in \Omega$ where $$D(\theta) = \left(\frac{\partial \tau}{\partial\theta_1}\, , \, \cdots\, , \, \frac{\partial \tau}{\partial \theta_k}\right)^t.$$
\end{the multiparameter case}



\section{Limiting Distributions (Review)}
\begin{defn}[convergence in probability]
A sequence of random variables  $X_1, \, \cdots\, , \, X_n$ converges in probability to a constant $C$ if for each $\varepsilon > 0$
$$\lim_{n\rightarrow \infty} P\left(\left|X_n - C\right| > \varepsilon\right) = 0$$
we write $X_n \underset{P}{\longrightarrow} C$.
\end{defn}




\begin{chebyshev's inequality}
 Suppose $X$ is a random variable with $E(X) = \mu$ and $var(X) = \sigma^2 < \infty$, then for any $K > 0$
 $$P\left(\left|X - \mu\right| > K\right)  < \frac{\sigma^2}{K^2}.$$
\end{chebyshev's inequality}


\begin{weak law of large numbers}[WLLN]
Suppose $X_1, \, \cdots\, , \, X_n$ is a random sample from a distribution with $E(X_i)   = \mu$ and $var(X_i) = \sigma^2 < \infty$ then  $$\overline{X_n} = \frac{1}{n}\sum^n_{i = 1} X_i \underset{P}{\longrightarrow} \mu.$$
\end{weak law of large numbers}




\begin{defn}[Convergence in distribution]
Let $X_n$ be a sequence of random variables. We say $X_n$ converges is distribution to a random variable $X$ if $$\lim_{n\rightarrow \infty}P\left(X_n \leq x\right) = P(X \leq x)$$ for all values of $x$ at which the right hand side $\, F(x) = P(X\leq x)$ is continuous. We write $X_n \underset{D}{\longrightarrow} X$.
\end{defn}





\begin{central limit theorem}
Suppose $X_1, \, \cdots\, , \, X_n$ is a  random sample from a distribution with  $E(X_i) = \mu$ and $var(X_i) = \sigma^2 < \infty$. Then 
$$Y_n  = \frac{\sum^n_{i = 1} X_i - n\mu}{\sqrt{n\, \sigma^2}} = \frac{\sqrt{n}\left(\overline{X} - \mu\right)}{\sigma} \, \underset{P}{\longrightarrow} N(0,1).$$
\end{central limit theorem}



\begin{limit theorems}
\begin{enumerate}
\item If $X_n \underset{P}{\longrightarrow} a$ where $a$ is a constant and $g$ is a real-valued function which is continuous at $a$, then $$g(X_n) \underset{P}{\longrightarrow} g(a).$$

\item If $X_n \longrightarrow X$ and $g$ is a real-valued continuous function, then $$g(X_n) \underset{D}{\longrightarrow} g(x).$$

\item If $X_n \underset{P}{\longrightarrow} X$ then $X_n \underset{D}{\longrightarrow} X$.

\item (Slutsky) If $X_n \underset{P}{\longrightarrow} a$ and $Y_n \underset{D}{\longrightarrow} Y$, then 
   \begin{enumerate}
        \item[(a)] $X_n + Y_n \, \underset{D}{\longrightarrow}\, a + Y$
        \item[(b)] $X_nY_n \, \underset{D}{\longrightarrow}\, aY$
        \item[(c)] $\frac{X_n}{Y_n} \, \underset{D}{\longrightarrow} \, \frac{Y}{a};\hspace{0.3cm} a\neq 0$
   \end{enumerate}
   
\item Suppose $a$ and $b>0$ are constants and $$n^b\left(X_n - a\right) \, \underset{D}{\longrightarrow}\, X.$$ Let $g$ be a real-valued function that is differentiable and whose derivative $g'$ is continuous at $a$, then  $$n^b\, \left[g(\lambda^n) - g(a)\right] \, \underset{D}{\longrightarrow}\, g'(a)\, X.$$
\end{enumerate}
\end{limit theorems}





\section{Large Sample properties of $ML$ estimators}
Large sample theory or the asymptotic behavior of estimators as the sample size $n \rightarrow \infty$ is a form of optimality for evaluating estimators.

\begin{defn}
Consider a sequence of estimators $T_n$ where the subscript $n$ indicates that the estimator has been obtained from the data $(X_1, \, \cdots\, , \, X_n)$ with sample size $n$. Then the sequence is said to be a consistent sequence of estimators of $\tau(\theta)$ if $T_n \underset{P}{\longrightarrow} \tau(\theta)$ for all $\theta \in \Omega$.
\end{defn}




\begin{thm}
Suppose  $X_1, \, \cdots\, , \, X_n$ is a random sample from a regular statistical model $\{f_{\theta}(x):\, \theta \in \Omega\}$. 
$$S_1(\theta, X) = \frac{\partial}{\partial\theta}\log f_{\theta}(x)$$
be a score function for a sample of size one. Then with probability tending to 1 as $n\longrightarrow \infty$, the likelihood equation
$$\sum^n_{i = 1} S_1(\theta, X_i) = 0$$
has a root $\widehat{\theta_n}$ such that $\widehat{\theta_n}$ converges in probability to $\theta_0$ , the true value of the parameter as $n \longrightarrow \infty$.\\
i.e The $ML$ estimator is consistent $\left(\widehat{\theta_n}\, \underset{P}{\longrightarrow} \theta_0\right)$.
\end{thm}



\begin{thm}
Let $X_1, \, \cdots\, , \, X_n$ be a random sample from a regular statistical model $\{f_{\theta}(x):\, \theta \in \Omega\}$. Suppose $\widehat{\theta_n}$ is a consistent root of the likehood equation as in Theorem 2.5.2.\\

Let $$J_1(\theta) = E\left[-\frac{\partial^2}{\partial\theta^2}\log f_{\theta}(x)\right]$$ be the Fisher information for a sample of  size one. Then 
$$\sqrt{n}\left(\widehat{\theta_n} - \theta_0\right)\, \underset{D}{\longrightarrow}\, Y \thicksim N\left(0, \frac{1}{J_1(\theta_0)}\right)$$
where $\theta_0$ is the true value of the parameter.\\
This result may also be written as $$\sqrt{n\, J_1(\theta_0)}\, \left(\widehat{\theta_n} - \theta_0\right) = \sqrt{J(\theta_0)}\, \left(\widehat{\theta_n} - \theta_0\right) \, \underset{D}{\longrightarrow}\, N(0,1).$$
By the Limiting Theorem it also follows that
 $$\sqrt{n}\, \left(\tau\left(\widehat{\theta_n}\right)  - \tau(\theta_0)\right) \, \longrightarrow\, W\thicksim N\left(0, \frac{(\tau'(\theta_0))^2}{J_1(\theta_0)}\right).$$
\end{thm}



\textbf{Note:} The theorem asserts:
\begin{enumerate}
\item the $ML$ estimator is asymptotically unbiased.
\item the asymptotic variance of the $ML$ estimator approaches the $CRLB$.
\end{enumerate}



\begin{defn}
A random variable $Q(X,\theta)$ whose distribution does not depend on $\theta$ or any unknown parameters is called a \textit{Pivotal quantity}.
\end{defn}

An asymptotic pivotal quantity is a function of the data and the parameter $\theta$ whose distribution converges to a distribution independent of $\theta$ as $n \longrightarrow \infty$. e.g  If $X_1, \, \cdots\, , \, X_n$ is a  random sample from a $N(\theta, \sigma^2)$ distribution where $\sigma^2$ is known, then  $$T = \frac{\sqrt{n}\, \left(\overline{X} - \theta\right)}{\sigma}$$ is a pivotal quantity.\\


If $X_1, \, \cdots\, , \, X_n$ is a random sample of from a distribution with mean $\theta$ and known variance $\sigma^2$ then the asymptotic distribution of $T$ is $N(0,1)$ by the CLT and $T$ is an asymptotic pivotal quantity pivotals can be used to construct confidence intervals as follows write 
$$P(C_1 < Q(X,\theta) < C_2) = 1 - \alpha$$
provided $Q$ is monotone, rewrite the statement.
$$P\left(C_1(X) < \theta < C_2(X)\right) =  1- \alpha$$
i.e $[C_1(X) , C_2(X)]$ is a random C.I with confidence coefficient $1 - \alpha$.\\

In cases, where an exact pivotal cannot be found, we can use the asymptotic distribution of $\widehat{\theta_n}$. An approximate C.I. can be constructed based on the asymptotic pivotals
$$\sqrt{J(\widehat{\theta_n})}\, \left(\widehat{\theta_n} - \theta_0\right) \, \underset{D}{\longrightarrow} Z\thicksim N(0,1)$$ or 
$$\sqrt{I(\widehat{\theta_n})}\, \left(\widehat{\theta_n} - \theta_0\right)\, \underset{D}{\longrightarrow} Z\thicksim (0,1).$$
\vspace*{4pt}


\begin{exa}
Let $X_1, \, \cdots\, , \, X_n$ be a random sample from the distribution with pdf
$$f_{\theta}(x) = \theta\, x^{\theta - 1}\,\, , \hspace{0.5cm} 0< x < 1.$$
\begin{enumerate}
\item[(a)] Find the $ML$ estimator of $\theta$.
\item[(b)] Find an approximate $95\%$ CI for $\theta$ based on the asymptotic distribution of $\widehat{\theta_n}$. 
\item[(c)] Find the exact $95\%$ C.I for $\theta$.
\end{enumerate}
\end{exa}


\begin{proof}[\textbf{\textit{Solution}}]
\begin{enumerate}
\item[(a)] $f_{\theta}(x) = \theta\, x^{\theta  - 1}$
$$L(\theta) = \prod^n_{i= 1} \theta\, x_i^{\theta - 1} = \theta^n\, \left(\prod^n_{i = 1} x_i\right)^{\theta - 1}.$$

$$\mathcal{L}(\theta)  = n\, \log\theta + (\theta - 1)\, \log\prod^n_{i = 1} x_i.$$

$$S(\theta) = \frac{n}{\theta} + \log\prod^n_{i = 1} x_i.$$
$S(\theta) = 0$
$$\frac{n}{\theta} + \sum^n_{i = 1} \log x_i = 0$$

$$\widehat{\theta_n} = \frac{n}{-\sum^n_{i  = 1} \log x_i}$$
$\widehat{\theta} = \frac{n}{T}$ where $T = -\sum^n_{i = 1} \log X_i$ is the $ML$ estimator of $\theta$.


\item[(b)] $I(\theta) = -\left(-\dfrac{n}{\theta^2}\right) = \dfrac{n}{\theta^2}$
$$J(\theta) = E\left[I(\theta,X)\right] = \dfrac{n}{\theta^2}.$$

$$\sqrt{J\left(\widehat{\theta_n}\right)}\, \left(\widehat{\theta_n} - \theta\right) \, \underset{D}{\longrightarrow}\, Z\thicksim N(0,1)$$

$$\frac{\widehat{\theta_n}  - \theta}{\dfrac{1}{\sqrt{J\left(\widehat{\theta_n}\right)}}}\, \underset{D}{\longrightarrow}\, Z \thicksim N(0,1).$$

An approximate 95$\%$ C.I for $\theta$ is given by 
$$\widehat{\theta_n} \pm 1.96\, \frac{1}{\sqrt{J(\widehat{\theta_n})}}$$
$$\frac{n}{T} \pm 1.96\, \sqrt{\frac{1}{\dfrac{n}{\widehat{\theta_n}^2}}}$$
$$\frac{n}{T}\pm 1.96 \frac{\widehat{\theta_n}}{\sqrt{n}}$$
$$\frac{n}{T} \pm 1.96 \, \frac{\left(\frac{n}{T}\right)}{\sqrt{n}}$$
$$\frac{n}{T} \pm 1.96\, \frac{\sqrt{n}}{T}$$

\item[(c)] $-\log X_i \thicksim ExP\left(\frac{1}{\theta}\right)$
$$T   = - \sum^n_{i = 1} \log X_i \thicksim GAM \left(n, \frac{1}{\theta}\right)$$
then $Y = 2\theta T\thicksim \chi^2_{2n}$.\\
Now
% DIAGRAM 1


\begin{figure}[H]
\centering
\begin{tikzpicture}[=>stealth]
\draw[-Stealth](0,0)--(7.5,0)node[right]{$y$};
\draw[-Stealth](0,0)--(0,3)node[above]{$g(y)$};
\draw[smooth,domain=0:14,scale=0.5,variable=\x]plot(\x,{8*\x*\x*exp(-1/2*\x)/4});
\draw(5,0)--(5,0.68);
\draw(0.7,0)--(0.7,0.95);
\path[pattern=north west lines, smooth, domain=10:14,scale=0.5,variable=\x]plot(\x,{8*\x*\x*exp(-1/2*\x)/4})--(14,0.18)--(14,0)--(10,0)--(10,0.68)--cycle;
\path[pattern=north east lines, smooth, domain=0:1.4,scale=0.5,variable=\x]plot(\x,{8*\x*\x*exp(-1/2*\x)/4})--(1.4,0.95)--(1.4,0)--(0,0)--cycle;
\draw[-Stealth](6.5,0.8)--(6,0.8)--(6,0.2);
\draw[-Stealth](0.3,0.8)--(0.5,0.8)--(0.5,0.3);
\draw(6.5,0.8)node[right]{0.025}
(5,0)node[below]{$b$}
(0.7,0)node[below]{$a$}
(0.3,0.8)node[left]{$0.025$};
\end{tikzpicture}
\end{figure}










Thus an exact $95\%$ equal tailed C.I is given as follows:
$$P(a < Y < b) = 0.95$$
$$P\left(a < 2\theta T < b\right) = 0.95$$
$$P\left(\frac{1}{2T} < \theta < \frac{b}{2T}\right) = 0.95$$
Therefore $\left[\frac{a}{2T}\, , \, \frac{b}{2T}\right]$ is an exact $95\%$ C.I for $\theta$ where
$$a = \chi^2_{0.025,2n}$$
$$b = \chi^2_{0.975, 2n}$$
i.e
$$P\left(\chi^2_{2n} \leq a\right) = 0.025$$
$$P\left(\chi^2_{2n} \leq b \right) = 0.975.$$
\end{enumerate}
\end{proof}






\begin{exa}
Suppose $X_1, \, \cdots\, , \, X_n$ is a random sample from the $POI(\theta)$ distribution. Find approximate $95\%$ C.I's for $\theta$ and $\tau(\theta) = e^{-\theta}$ based on the asymptotic distribution of $\widehat{\theta_n}$.
\end{exa}

\begin{proof}[\textbf{\textit{Solution.}}]
From example 2.1.2 $$\widehat{\theta}_n = \overline{X}_n$$
$$J(\theta) = \frac{n}{\theta}$$
an approximate $95\%$ C.I for $\theta$ is given by
$$\widehat{\theta}_n \, \pm\, 1.96\, \sqrt{\dfrac{1}{J(\widehat{\theta}_n)}}$$
$$\overline{X}_n \, \pm\, 1.96\, \sqrt{\dfrac{\widehat{\theta}_n}{n}}$$
$$\overline{X}_n \, \pm \, 1.96\, \sqrt{\dfrac{\overline{x}_n}{n}}$$
By Theorem 2.5.3
$$\sqrt{n}\, \left(\tau(\widehat{\theta}_n) - \tau (\theta)\right) \, \underset{D}{\longrightarrow}\, N\left(0, \frac{\left(\tau'(\theta)\right)^2}{J_1(\theta)}\right)$$
i.e
$$T\left(\widehat{\theta}_n\right)  - \tau(\theta) \, \underset{D}{\longrightarrow}\, N\left(0\, , \, \frac{\left(\tau'(\theta)\right)^2}{J(\theta)}\right)$$

$$\tau(\theta) = e^{-\theta}\, \implies\, \tau'(\theta) = - e^{-\theta}$$
an approximate $95\%$ C.I for $\tau(\theta) = e^{-\theta}$ is given by
$$\tau\left(\widehat{\theta}_n\right)\pm 1.96\, \sqrt{\dfrac{\left(\tau'(\widehat{\theta}_n)\right)^2}{J(\widehat{\theta}_n)}}$$

$$e^{-\overline{x}_n} \, \pm 1.96\, \sqrt{\dfrac{\left(e^{-\widehat{\theta}_n}\right)^2}{n/\widehat{\theta}_n}}$$

$$e^{-\overline{x}} \, \pm\, 1.96\, e^{-\overline{x}}\, \sqrt{\dfrac{\overline{x}_n}{n}}.$$
\end{proof}




\begin{exa}
Let $X_1, \, \cdots\, , \, X_n$ be a random sample from the $EXP(1,\theta)$ distribution. Find the $ML$ estimator of $\theta$. Show that the $ML$ estimator $\widehat{\theta}_n$ is a consistent  estimator of $\theta$. Find the distribution of $n\left(\widehat{\theta}_n - \theta\right)$ and a $95\%$ C.I for $\theta$.
\end{exa}

\begin{proof}[\textbf{\textit{Solution.}}]
\begin{align*}
L(\theta) &  = \prod^n_{i = 1} e^{-(x_i - \theta)}\hspace{0.3cm} , \hspace{0.6cm} x_1, \, \cdots\, , \, x_n > \theta\\
          & = e^{-\sum^n_{i  = 1}(x_i - \theta)}\hspace{0.3cm} , \hspace{0.6cm} x_{(1)} > \theta\\
          & = e^{n\theta - \sum^n_{i = 1} x_i}\hspace{0.3cm} , \hspace{0.6cm} x_{(1)} > \theta
\end{align*}
% DIAGRAM 2
\begin{figure}[H]
\centering
\begin{tikzpicture}
\draw[-Stealth](0,0)--(4.5,0)node[right]{$\theta$};
\draw[-Stealth](0,0)--(0,3)node[left]{$\L(\theta)$};
\draw(0,0)parabola(3.5,1.7)
(3.5,1.7)node[draw, shape =  circle, scale = 0.01cm, fill = black!]{}
(3.5,-0.08)--(3.5,0.08)
(1.75,-0.08)--(1.75,0.08)
(3.5,-0.08)node[below]{$x_{(1)}$};
\end{tikzpicture}
\end{figure}








Therefore $\, \widehat{\theta}_n  = X_{(1)}$.
$$\widehat{\theta}_n \underset{P}{\longrightarrow} \theta\hspace{0.6cm} \text{if}\hspace{0.6cm} \lim_{n \rightarrow\infty} P\left(\left|\widehat{\theta}_n - \theta\right| > \varepsilon \right) = 0.$$
\begin{align*}
P\left(\left|\widehat{\theta}_n - \theta\right| > \varepsilon\right) & = P\left(\widehat{\theta}_n -\theta > \varepsilon\right) + P\left(\widehat{\theta}_n - \theta < -\varepsilon\right)\\
& = P(X_{(1)} - \theta > \varepsilon) + P(X_{(1)} - < -\varepsilon)\\
& = P(X_{(1)} > \varepsilon + \theta) + P(X_{(1)} < -\varepsilon + \theta).
\end{align*}

Now, let $Y = X_{(1)}$  $$f_{X_{(1)}}(y) = n\left[1 - F_X(y)\right]^{n  - 1}\, f_X(y)$$

$$F(y)  = \int^y_{\theta} e^{-(x - \theta)}\, dx = -e^{-(x - \theta)}\Big|_{\theta}^y = 1 + e^{-(y - \theta)}.$$

$$f(y) = n  \left(e^{-(y - \theta)}\right)^{n - 1}\, e^{-(y - \theta)} = n\, e^{-n(y - \theta)}\hspace{0.3cm} , \hspace{0.5cm} y > \theta.$$

\begin{align*}
P\left(\left|X_{(1)} - \theta\right| > \varepsilon \right) & = \int^{\infty}_{\theta + \varepsilon} n\, e^{-n(y -\theta)}\, dy \, +\, \underbrace{\int_{-\infty}^{\theta - \varepsilon} n\, e^{-n(y - \theta)}\, dy}_0\\
& =\frac{n\, e^{-n(y - \theta)}}{-n}\Big|_{\theta + \varepsilon}^{\infty}\\
& = - \left(0 - e^{-n(\theta + \varepsilon - \theta)}\right)\\
& = e^{-n\varepsilon}\, \longrightarrow \, 0.
\end{align*}
i.e $\lim_{n\rightarrow \infty} P\left(\left|X_{(1)} - \theta\right| > \varepsilon \right) = 0$. Therefore $\, X_{(1)} \underset{P}{\longrightarrow} \theta$ (it is consistent).\\

Let $\, T = n\left(\widehat{\theta}_n - \theta\right)$
\begin{align*}
F(t) & = P\left(T\leq t\right) = P\left(n\left(\widehat{\theta}_n - \theta\right) \leq t\right)\\
     & = P\left(X_{(1)} - \theta \leq \frac{t}{n}\right)\\
     & = P\left(X_{(1)} \leq \theta + \frac{t}{n}\right)\\
     & = \int^{\theta + \frac{t}{n}}_{\theta} n\, e^{-n(y  - \theta)}\, dy\\
     & = -e^{-n(y - \theta)}\Big|^{\theta + \frac{t}{n}}_{\theta}\\
     & =  1 - e^{-t}.
\end{align*}
We find the pdf by differentiating
$$f(t) = e^{-t}\hspace{0.5cm} , \hspace{0.4cm} t> 0.\,\,\,\, \text{Therefore}\hspace{0.5cm} T\thicksim EXP(1).$$

A $95\%$ C.I for $\theta$ is calculated as follows:
$$P(a < n(X_{(1)} - \theta) < b) = 0.95$$
$$P\left(\frac{a}{n} < X_{(1)} - \theta < \frac{b}{n}\right) = 0.95$$
$$P\left(-\frac{a}{n} > \theta - X_{(1)} > -\frac{b}{n}\right) = 0.95$$
$$P\left(-\frac{b}{n}  < \theta - X_{(1)} < - \frac{a}{n}\right) = 0.95$$
$$P\left(X_{(1)} - \frac{b}{n} < \theta < X_{(1)} - \frac{a}{n}\right) = 0.95$$
$$\left[X_{(1)} - \frac{b}{n}\, , \, X_{(1)} - \frac{a}{n}\right].$$
where
% DIAGRAM 3
\begin{figure}[H]
\centering
\begin{tikzpicture}
\draw[-Stealth](0,0)--(5,0)node[below]{$t$};
\draw[-Stealth](0,0)--(0,3)node[above]{$f(t)$};
\draw[scale = 1, variable = \x, domain = 0.4:4.5, smooth,thick]plot({\x},{1/\x});
\draw[-Stealth](-0.5,1.2)--(0.2,1.2);
\draw[-Stealth](4,0.8)--(4,0.1);
\draw(0.5,0)--(0.5,2)
(0.5,0)node[below]{$a$}
(3.5,0)node[below]{$b$}
(3.5,0)--(3.5,0.285)
(-0.5,1.2)node[left]{$0.0.25$}
(4,0.8)node[above]{$0.025$};
\path[pattern = north east lines, smooth, domain = 0.4:0.5, variable = \x]plot({\x},{1/\x})--(0.5,2)--(0.5,0)--(0,0)--(0,2.5)--cycle;
\path[pattern= north west lines, smooth, domain = 3.5:4.5, variable = \x]plot({\x},{1/\x})--(4.5,0.22)--(4.5,0)--(3.5,0)--(3.5,0.285)--cycle;
\end{tikzpicture}
\end{figure}









$$F(a) = 1 - e^{-a} = 0.025 \, \implies e^{-a} = 0.975$$
$$a = -\log 0.975$$


$$F(b) = 1 - e^{-b} = 0.975 \, \implies e^{-b} = 0.025$$
$$b  = -\log 0.025.$$

$$\left[X_{(1)} + \frac{\log 0.025}{n} \, , \, X_{(1)} + \frac{\log 0.975}{n}\right].$$


\end{proof}








\begin{problem}
Suppose $X_1, \, \cdots\, , \, X_n$ is a random sample from the $UNIF(0,\theta)$ distribution. Show that the $ML$ estimator $\widehat{\theta}_n$ is a consistent estimator of $\theta$ and find the asymptotic distribution of $n\left(\widehat{\theta}_n - \theta\right)$.
\end{problem}


\begin{problem}
The lifetimes of a certain type of electronic equipment follows the $EXP(\theta)$ distribution. Ten components where tested independently with the observed lifetimes, to the nearest days, given by 
\begin{table}[H]
\centering
\begin{tabular}{cccccccccc}
70 & 11 & 66 & 5 & 20 & 4 & 35 & 40 & 29 & 8\\
\end{tabular}
\end{table}
Find 
\begin{enumerate}
\item[(a)] $ML$ estimator of $\theta$ and verify its a maximum.
\item[(b)] Fisher information and evaluate an approximate $95\%$ C.I for $\theta$.
\item[(c)] The exact $95\%$ C.I for $\theta$.
\end{enumerate}
\end{problem}
\vspace{4pt}



\begin{the multiparameter case}
If $\theta = (\theta_1, \, \cdots\, , \, \theta_k)^t$, the likelihood equation consists of $k$ equations in the $K$ unknown parameters. Under similar regularity conditions, the components of $\widehat{\theta}_n$ each converge in probability to the corresponding component of $\theta$.\\
Similarly $$\sqrt{n}\, \left(\widehat{\theta}_n - \theta\right)\, \underset{D}{\longrightarrow}\, Y\, \thicksim\, MVN\left(0\, , \, \left(J_1(\theta)\right)^{-1}\right)$$ where $J_1(\theta)$ is a Fisher information matrix for a sample of size one.\\
It follows that $$\sqrt{n}\, \left(\tau\left(\widehat{\theta}_n\right) - \tau(\theta)\right) \, \underset{D}{\longrightarrow}\, w\thicksim MVN\left(0\, , \, D(\theta)^t\, (J_1(\theta))^{-1}\, D(\theta)\right)$$
where $$D(\theta) = \left(\dfrac{\partial\tau}{\partial\theta_1}\, , \, \cdots\,, \, \frac{\partial\tau}{\partial\theta_k}\right)^t.$$
Approximate C.I's for a single parameter, say $\theta_i$, can be constructed based on the approximate normality of the following:
$$\frac{\left(\widehat{\theta}_n\right)_i - (\theta)_i}{\sqrt{\left[J\left(\widehat{\theta}_n\right)\right]^{-1}_{ii}}}\hspace{1cm} \text{or}\hspace{1cm}\frac{\left(\widehat{\theta}_n\right)_i - (\theta)_i}{\left[I\left(\widehat{\theta}_n\right)\right]^{-1}_{ii}}$$ where $(\theta)_i$ is the $i^{\text{th}}$ component of the vector $\theta$ and $\left[J(\theta)\right]_{ii}$ is the $(i,i)^{\text{th}}$ of the matrix $\left[J(\theta)\right]^{-1}$.
\end{the multiparameter case}





\begin{exa}
Suppose $X_1, \, \cdots\, , \, X_n$ is a random sample from the $GAM(\alpha\, , \, \beta)$ distribution and $\theta = (\alpha\, , \, \beta)^t$. Find an approximate $95\%$ C.I for $\beta$ and $\tau(\theta) = E(X) = \alpha\beta$.
\end{exa}



\begin{proof}[\textbf{\textit{Solution.}}] From example 2.2.8
$$J(\theta) = \begin{pmatrix}
n\, \Psi'(\alpha) & \dfrac{n}{\beta}\\\\
\dfrac{n}{\beta} & \dfrac{n\alpha}{\beta^2}\\
\end{pmatrix}$$

$$\left|J(\theta)\right| = \frac{n^2\, \alpha\, \Psi'(\alpha}{\beta^2} - \frac{n^2}{\beta^2} = \frac{n^2\, \left(\alpha\, \Psi'(\alpha)  - 1\right)}{\beta^2}.$$

$$\left[J(\theta)\right]^{-1} = \frac{\beta^2}{n^2\, \left(\alpha\, \Psi'(\alpha)  - 1\right)}\begin{pmatrix}
\dfrac{n\alpha}{\beta^2} & -\dfrac{n}{\beta}\\\\
-\dfrac{n}{\beta} & n\, \Psi'(\alpha)\\
\end{pmatrix} = \frac{\beta^2}{n\, \left(\alpha\, \Psi'(\alpha)  - 1\right)}\begin{pmatrix}
\dfrac{\alpha}{\beta^2} & -\dfrac{1}{\beta}\\\\
-\dfrac{1}{\beta} &  \Psi'(\alpha)\\
\end{pmatrix}.$$

a $95\%$ C.I for $\beta$ is given by 
$$\widehat{\beta}_n \, \pm\, 1.96\, \sqrt{var\left(\widehat{\beta}_n\right)}$$

$$\widehat{\beta}_n \, \pm\, 1.96\, \sqrt{\left[J\left(\widehat{\theta}_n\right)\right]_{22}}$$

$$\widehat{\beta}_n\, \pm\, 1.96\, \sqrt{\dfrac{\widehat{\beta}_n^2\, \Psi'\left(\widehat{\alpha}_n\right)}{n\,\left(\widehat{\alpha}_n\, \Psi\left(\widehat{\alpha}_n\right)\right)^{-1}}}.$$

$$\tau\left(\widehat{\theta}_n\right) - \tau(\theta) \, \underset{D}{\longrightarrow}\, W\thicksim MVN\left(0\, , \, D(\theta)^t\, \left[J(\theta)\right]^{-1}\, D(\theta)\right)$$

$$\tau(\theta) = \alpha\beta$$

$$D(\theta) = \begin{pmatrix}
\dfrac{\partial\tau}{\partial \alpha}\\\\ \dfrac{\partial \tau}{\partial\beta}\\
\end{pmatrix}  = \begin{pmatrix}
\beta\\  \alpha\\
\end{pmatrix}.$$
\begin{align*}
D(\theta)^t\, \left[J(\theta)\right]^{-1}\, D(\theta) & = \begin{pmatrix}
\beta & \alpha\\
\end{pmatrix} \frac{\beta^2}{n\, \left(\alpha\, \Psi'(\alpha)  - 1\right)}\begin{pmatrix}
\dfrac{\alpha}{\beta^2} & -\dfrac{1}{\beta}\\\\
-\dfrac{1}{\beta} &  \Psi'(\alpha)\\
\end{pmatrix}\, \begin{pmatrix}
\beta\\ \alpha\\
\end{pmatrix}\\\\
& = \frac{\beta^2}{n\, \left(\alpha\, \Psi'(\alpha)  - 1\right)} \begin{pmatrix}
\dfrac{\alpha}{\beta} - \dfrac{\alpha}{\beta} && -1 + \alpha \, \Psi'(\alpha)\\
\end{pmatrix}\, \begin{pmatrix}
\beta\\
\alpha\\
\end{pmatrix}.\\\\
& = \frac{\beta^2}{n\, \left(\alpha\, \Psi'(\alpha)  - 1\right)}\, \begin{pmatrix}
0 + \alpha\left(\alpha \, \Psi'(\alpha) - 1\right)\\
\end{pmatrix}\\\\
& = \frac{\alpha\, \beta^2}{n}
\end{align*}
a $95\%$ C.I for $\tau(\theta) = \alpha\beta$ is given by $$\widehat{\alpha}_n\widehat{\beta}_n\, \pm \, 1.96\, \sqrt{\dfrac{\widehat{\alpha}_n\, \widehat{\beta}_n^2}{n}}.$$
\end{proof}








\begin{problem}
For problem 2.2.5 find an approximate $95\%$ C.I for $$\ \tau (\theta) = cov(X_1, X_2) = -n\theta_1\theta_2.$$
\end{problem}




\begin{problem}
In problem 2.2.9, find an approximate $95\%$ C.I for $\beta$ and $E(X)$.
\end{problem}





\begin{problem}
Suppose $X_1, \, \cdots\, , \, X_n$ is a random sample from  the $BETA(\alpha\, ,\, \beta)$ distribution. Explain how you would find the $ML$ estimate of $\theta = (\alpha, \beta)^t$, approximate $95\%$ C.I's for $\alpha$ and $$\tau(\theta) = E(X) = \frac{\alpha}{\alpha + \beta}.$$
\end{problem}










\newpage
\chapter{OTHER ESTIMATION CRITERIA}
\section{Best Linear Unbiased Estimators}
An estimator is a best linear unbiased estimator (BLUE) if it is Linear and unbiased estimator with smallest variance.

Assume that $Y = (Y_1, \, \cdots\, , \, Y_n)^t$ has expectation $E(Y) = XB$ where $X$ is some $n\times k$ matrix having rank $k$ and $B = (\beta_1, \, \beta_2, \, \cdots\, , \, \beta_k)^t$ is a vector of unknown parameters. Suppose the convariance of $Y$ is $\sigma^2 B$ where $B$ is a non-singular matrix and $\sigma^2$ is scalar. We want to find the BLUE of a linear combination of $B$, say $\theta = a^t B$ where $a$ is a $k$-dimensional vector.



\begin{gauss}
Let $E(Y) = XB$ and $var(Y) = \sigma^2 B$, then the best linear unbiased estimator (BLUE) of $\theta = a^t B$ is
$$\hat{\theta}  = a^t \left(X^t B^{-1}X\right)^{-1}X^tB^{-1}Y.$$
\end{gauss}


\begin{exa}
Suppose $Y_1, \, \cdots\, , \, Y_n$ are independent unbiased estimators of $\theta$ with unknown variance $var(Y_i) = \sigma^2, \, \, i = 1, \, 2, \, \cdots\, , \, n$. Find the BLUE of $\theta$.
\end{exa}

\begin{proof}[\textbf{Solution.}] $Y = XB + \varepsilon$
$$Y_i = \theta + \varepsilon_i \, , \, \hspace{0.3cm} E(Y_i) = \theta\, , \, \hspace{0.3cm} E(\varepsilon_i) = 0\, , \, \hspace{0.2cm} var(Y_i) = var(\varepsilon_i) = \sigma^2.$$

$$\begin{pmatrix}
Y_1\\
\vdots\\
Y_n\\
\end{pmatrix}  = \begin{pmatrix}
1\\
\vdots\\
1\\
\end{pmatrix}\, (\theta) + \begin{pmatrix}
\varepsilon_1\\
\vdots\\
\varepsilon_n\\
\end{pmatrix}$$

$$var(Y) = \sigma^2\, I = \begin{pmatrix}
\sigma^2 & \cdots & 0\\
\vdots &  & \vdots\\
0 & \cdots & \sigma^2\\
\end{pmatrix} = \sigma^2\, \begin{pmatrix}
1 & \cdots & 0\\
\vdots & \vdots & \vdots\\
0 & \cdots & 1\\
\end{pmatrix}.$$

$$\hat{\theta} = a^t\left(X^tB^{-1}X\right)^{-1}X^tB^{-1}Y$$

$$X^tB^{-1} = (1\,\, \cdots\, \, 1)\frac{1}{\sigma^2}I = \frac{1}{\sigma^2}\, (1\,\, \cdots\,\, 1).$$

$$X^tB^{-1}X = \frac{1}{\sigma^2}\, (1 \, \cdots\,  \, 1)\begin{pmatrix}
1\\
\vdots\\
1\\
\end{pmatrix} = \frac{n}{\sigma^2}.$$

$$\left(X^tB^{-1}X\right)^{-1} = \frac{\sigma^2}{n}.$$

$$\left(X^tB^{-1}X\right)^{-1} = \frac{\sigma^2}{n}\cdot\frac{1}{\sigma^2}\, (1\,  \cdots\, 1) = \frac{1}{n}\,(1\, \cdots\, 1).$$
\begin{align*}
\left(X^tB^{-1}X\right)^{-1}X^tB^{-1}Y  & = \frac{1}{n}(1\, \, \cdots\, \, 1)\begin{pmatrix}
Y_1\\
\vdots\\
Y_n\\
\end{pmatrix}
 = \frac{1}{n}\sum^n_{i = 1} Y_i
 = \overline{Y}.
\end{align*}
\end{proof}







\begin{problem}
Suppose $Y_1, \, \cdots\, , \, Y_n$ are independent unbiased estimators of $\theta$ with known variances $var(Y_i) = \sigma_i^2, \, i = 1, \, 2, \, \cdots\, , \, n$. Find the BLUE of $\theta$.
\end{problem}


\begin{problem}
Suppose $Y_{ij}, \, i = 1, \, 2, \ j  = 1, \, 2, \, \cdots\, , n$ are independent random variables with $E(Y_{ij}) = \mu  + \alpha_i$ and $var(Y_{ij}) = \sigma^2$ where $\alpha_1 + \alpha_2 = 0$. Find the BLUE of $\alpha_1$.
\end{problem}






\section{Equivariant Estimators}
A model $\{f_{\theta}(x):\, \theta \in (-\infty,\infty)\}$ such that $f_{\theta}(x) = f_0(x - \theta)$ with $f_0$ known is called a location invariant family and $\theta$ is called a location parameter.

\begin{exa}
Show that the following  families are location invariant:
\begin{multicols}{3}
\begin{enumerate}
\item[(a)] $N(\theta,1)$
\item[(b)] $CAU(1,\theta)$
\item[(c)] $EXP(1,\theta)$
\end{enumerate}
\end{multicols}
\end{exa}


\begin{proof}[\textbf{Solution.}]
\begin{enumerate}
\item[(a)]  $f_{\theta}(x) = \dfrac{1}{\sqrt{2\pi}}e^{-\frac{1}{2}(x - \theta)^2}\, = f_{0}(x - \theta)\, , \hspace{0.2cm} \theta \in \mathbb{R}$ where
$$f_0(z) = \frac{1}{\sqrt{2\pi}}e^{-\frac{1}{2}z^2}.$$
\item[(b)] $f_{\theta}(x) = \dfrac{1}{\pi\left[1 + (x - \theta)^2\right]} = f_0(x - \theta)\,, \hspace{0.3cm} \theta \in \mathbb{R}$ where
$$f_0(z) = \frac{1}{\pi\left[1 + z^2\right]}.$$
\item[(c)] $f_{\theta}(x) = e^{-(x  - \theta)} = f_0(x - \theta)\, , \hspace{0.3cm} x > \theta\, ,\,\, x - \theta  > 0$ where
$$f_0(z) = e^{-z}\, , \hspace{0.3cm} z>0.$$
\end{enumerate}
\end{proof}


\begin{defn}
An estimator $\tilde{\theta}(X_1, \, \cdots\, , \, X_n)$ is location equivariant if $$\tilde{\theta}(x_1 + a, \, \cdots\, , \, x_n + a) = \tilde{\theta}(x_1, \, \cdots\, , \, x_n) + a$$ for any constant $a$ and all values of $(x_1, \, \cdots\, , \, x_n)$.
\end{defn}


\begin{exa}
Suppose $X_1, \, \cdots\, , \, X_n$ is a random sample from a $N(\theta,1)$ distribution. Show that the $UMVUE$ is a location equivariant estimator.
\end{exa}


\begin{proof}[\textbf{Solution.}]
$\tilde{\theta}(X_1, \, \cdots\, , \, X_n) = \overline{X}$  is the $UMVUE$ of $\theta$ 
\begin{align*}
\tilde{\theta}(x_1 + a, \, \cdots\, , \, x_n + a) & = \frac{1}{n}\sum^n_{i = 1} (X_i + a)\\
& = \frac{1}{n}\left(\sum^n_{i = 1} X_i + a\right)\\
& = \frac{1}{n}\sum^n_{i = 1} X_i + a\\
& =\tilde{\theta}(x_1, \, \cdots\, , \, x_n) + a.
\end{align*}
$\therefore\, \overline{X}$ is a location equivariant estimator.
\end{proof}



\begin{defn}
A model $\{f_{\theta}(x):\, \theta > 0\}$ such that $f_{\theta}(x)  = \frac{1}{\theta} \, f_1\left(\frac{x}{\theta}\right)$ with $f_1$ known is called a scale invariant family and $\theta$ is called a scale parameter.
\end{defn}


\begin{defn}
An estimator $\tilde{\theta}^k = \tilde{\theta}^k(x_1, \, \cdots\,, \, x_n)$ is scale invariant if $$\tilde{\theta}^k(cx_1, \, \cdots\,, \, cx_n) = c^k\, \tilde{\theta}^k(x_1, \, \cdots\,, \, x_n)$$ for all values of $(x_1, \, \cdots\, , \, x_n)$ and $c >0$.
\end{defn}




\begin{exa}
Let $X_1, \, \cdots\, , \, X_n$ be a random sample from the $N(0, \sigma^2)$ distribution. Show that
\begin{enumerate}
\item[(a)] $f_{\theta}(x)$ is a scale invariant family
\item[(b)] the $UMVUE$ of $\theta^2$ is a scale equivariant estimator.
\end{enumerate}
\end{exa}





\begin{proof}[\textbf{Solution.}]
\begin{enumerate}
\item[(a).] $\begin{aligned}[t]
f_{\theta}(x) = \frac{1}{\sqrt{2\pi \theta^2}}e^{-\frac{1}{2}x^2} = \frac{1}{\theta\, \sqrt{2\pi}}e^{-\frac{1}{2}\left(\frac{x}{\theta}\right)^2} = \frac{1}{\theta}\, f_1\left(\frac{x}{\theta}\right).
\end{aligned}$ 

where $f_1(z) = \dfrac{1}{\sqrt{2\pi}}e^{-\frac{1}{2}z^2}$.\\

$\therefore\,\, f_{\theta}(x)\,$ is a scale invariant family.



\item[(b).] $f_{\theta}(x)$ is REF distribution with complete sufficient statistic $T = \sum^n_{i  = 1} X_i^2.$ Thus
$\dfrac{1}{n}\sum_{i = 1}^n X_i^2\,$ is the $UMVUE$ of $\theta^2$.
$$\tilde{\theta}^2(x_1, \, \cdots\,, \, x_n) = \frac{1}{n}\sum^n_{i = 1} x_i^2.$$
\begin{align*}
\tilde{\theta}^2(cx_1, \, \cdots\,, \, cx_n) & = \frac{1}{n}\sum^n_{i = 1} \left(cx_i\right)^2
  = \frac{1}{n}\sum^n_{i = 1} c^2\, x_i^2\\
& = c^2\, \frac{1}{n}\sum^n_{i = 1} x^2_i\\
& = c^2\, \tilde{\theta}^2(x_1, \, \cdots\,, \, x_n).
\end{align*}
Therefore the $UMVUE$ of $\theta^2$ is a scale equivariant estimator.
\end{enumerate}
\end{proof}





\begin{problem}
Show that the $UNIF(\theta - 1, \theta + 1)$ distribution is a location invariant family.
\end{problem}


\begin{problem}
Let $X_1, \, \cdots\, , \, X_n$ be a random sample the $UNIF(0,\theta)$ distribution. Show that then
\begin{enumerate}
\item[(a)] $UNIF(0,\theta)$ densities is a scale invariant family.
\item[(b)] $UMVUE$ of $\theta$ is a scale equivariant estimator.
\end{enumerate}
\end{problem}



\section{Estimating Equations}
To find the M.L estimator, we usually solve the likelihood equation $$\sum^n_{i = 1} S_i(\theta,X_i) = 0.$$
A function such as the score function which is a function of both the observations and the parameter is called an estimating function for example, if we know $var(X_i) = \theta$ independently identically distributed $X_i$, then we can use the estimating function 
$$\Psi(\theta, X) = \sum^n_{i = 1} \left(X_i - \overline{X}\right)^2 - (n - 1)\theta$$
to estimate the parameter $\theta$ without any knowledge of the distribution, density or mean of $X$. The estimating function is the set equal to 0 and solved for $\theta$.


\begin{defn}
An estimating function $\Psi(\theta,X)$ is unbiased if $E\left[\Psi(\theta,X)\right] = 0$ for all $\theta$.
\end{defn}



\begin{thm}
Suppose $\widehat{\theta}$ is a root of the estimating equation $\Psi(\theta,X) = 0$, then $\widehat{\theta}$ is approximately normal with mean $\theta$ and variance $\left[J(\Psi,\theta)\right]^{-1}$ where  $$J(\Psi, \theta) = \frac{\left[E\left(\frac{\partial\Psi}{\partial \theta}\right)\right]^2}{var\left[\psi(\theta,X)\right]}.$$
\end{thm}



\textbf{Note:} $J(\Psi, \theta)$ is called the Godambe information of the estimating function.



\begin{exa}
Suppose $X = (X_1, \, \cdots\, , \, X_n)$ is a random sample from a distribution with $\, E\left[\log X_i\right] = e^{\theta}\, $ and $\, var\left[\log X_i \right] = e^{2\theta}\, , \, \, \, i = 1, \, 2, \, \cdots\, , \, n$.\\
Consider an estimating function $$\Psi(\theta, X) = \sum^n_{i = 1}\left(\log (X_i) - e^{\theta}\right)$$
\begin{enumerate}
\item[(a)] Show that $\Psi(\theta, X)$ is an unbiased estimating function.
\item[(b)] Find the estimator $\widehat{\theta}$ which satisfies $\Phi\left(\widehat{\theta}, X\right) = 0$.
\item[(c)] Construct an approximate $95\%$ C.I for $\theta$.
\end{enumerate}
\end{exa}

\begin{proof}[\textbf{Solution.}]
\begin{enumerate}
\item[(a)] 
$\begin{aligned}[t]
E\left[\Psi(\theta, X)\right] & = E\left[\sum^n_{i = 1}\left(\log X_i - e^{\theta}\right)\right]
 = \sum^n_{i = 1} \left(E\left(\log X_i\right) - e^{\theta}\right)\\
&  = \sum^n_{i = 1} \left(e^{\theta} - e^{\theta}\right)\\
& = \sum^n_{i = 1} 0\\
&  = 0.
\end{aligned}$\\
Therefore $\, \Psi(\theta,X)$ is an unbiased estimating function.

\item[(b)] $\Psi(\theta,X) = 0$
$$\sum^n_{i = 1} \left(\log X_i  - e^{\theta} \right) = 0$$
$$\sum^n_{i  = 1} \log X_i - n e^{\theta} = 0$$
$$\sum^n_{i = 1} \log X_i  = ne^{\theta}$$
$$\widehat{\theta} = \log \left[\frac{1}{n}\sum^n_{i = 1} \log X_i\right].$$

\item[(c)] $\widehat{\theta} \thicksim N\left(\theta, \left[J(\Psi, \theta)\right]^{-1}\right)$ a $95\%$ C.I for $\theta$  is given by 
$$\widehat{\theta} \pm 1.96 \sqrt{\dfrac{1}{J(\Psi,\widehat{\theta}})}$$

$$J(\Psi,\theta) = \frac{E\left[E\left(\frac{\partial \Psi}{\partial \theta}\right)\right]^2}{var\left[\Psi(\theta,X)\right]}$$
$$\Psi(\theta, X) = \sum^n_{i = 1} \left(\log X_i - e^{\theta}\right)$$
$$\frac{\partial \Psi}{\partial \theta}  = - \sum^n_{i = 1} e^{\theta} = -n e^{\theta}$$
$$E\left(\frac{\partial  \Psi}{\partial \theta}\right) = -n e^{\theta}.$$
$$\left[E\left(\frac{\partial \Psi}{\partial \theta}\right)\right]^2 = n^2e^{2\theta}$$

\begin{align*}
var\left[\Psi(\theta,X)\right] &  = var\left[\sum^n_{i = 1} \left(\log X_i - e^{\theta}\right)\right]
 = \sum^n_{i = 1} var\left(\log X_i\right)
 = \sum^n_{i = 1} e^{2\theta}
 = ne^{2\theta}.
\end{align*}

$$J(\Psi, \theta) = \frac{n^2\, e^{2\theta}}{n\,e^{2\theta}} = n.$$
$$J(\Psi,\widehat{\theta})  = n$$
$$\therefore\,\, \widehat{\theta} \, \pm\, 1.96\sqrt{\dfrac{1}{n}}\,.$$
\end{enumerate}
\end{proof}

\begin{exa}
Suppose $X_1, \, X_2, \, \cdots\, , \, X_n$ is a random sample from the Bernoulli$(\theta)$ distribution. Suppose that $\varepsilon_1, \, \cdots\, , \, \varepsilon_n$ are independent $N(0,\sigma^2)$ random variables independent of the $X_i's$. Define the estimating function $$\Psi\left[\theta, (X,Y)\right] = \sum^n_{i = 1}\left(Y_i - \theta \, X_i\right)$$
\begin{enumerate}
\item[(a)] Show that this is an unbiased estimating function of $\theta$.
\item[(b)] Find the estimator $\widehat{\theta}$ which satisfies $\Psi\left[\theta, (X,Y)\right] = 0$.
\item[(c)] Construct an approximate $95\%$ C.I for $\theta$.
\end{enumerate}
\end{exa}





\section{Bayes Estimation}
There are two major schools of thought in statistical inference, the frequentist and Bayesian School.  Frequentist consider a parameter as an unknown constant underlying a probability (density) function of the data.\\
Bayesians consider the parameter $\theta$ to be random with a prior distribution $\pi(\theta)$. The purpose of the data then is to adjust this distribution for $\theta$, to result in a posterior distribution of the parameter.



\begin{posterior}
Suppose a parameter is chosed at random according to the prior distribution $\pi(\theta)$, then the posterior distribution of the parameter $\theta$ given the data $x = (x_1, \, \cdots\, , \, x_n)$ of the parameter is the conditional distribution of $\theta$ given the data $x = (x_1, \, \cdots\, , \, x_n)$ $$\pi(\theta /x) = C\, \pi(\theta)\, \prod^n_{i = 1} f_{\theta}(x_i) = C\, \pi(\theta)\, L(\theta)$$ where $$C = \frac{1}{\int^{\infty}_{-\infty} \pi(\theta)\, L(\theta)\, d\theta}$$ and $L(\theta)$ is the likelihood function.
\end{posterior}

\begin{exa}
Suppose a coin is tossed $n$ times with probability of heads $\theta$. Assume that the prior probability of heads follows a $BETA(10,10)$ distribution. If the $n$ tosses result in $x$ heads, find the posterior density function of $\theta$.
\end{exa}

\begin{proof}[\textbf{Solution.}]
$f_{\theta}(x) = \binom{n}{x}\,\theta^x\,(1 - \theta)^{n - x}$
$$L(\theta) = \binom{n}{x}\, \theta^x\, (1 - \theta)^{n - x}$$
$$\theta \thicksim BETA(10,10)$$
$$\pi(\theta) = \frac{1}{B(10,10)} \, \theta^{10 - 1}\, (1 - \theta)^{10 - 1}   = C_1\, \theta^9 \, (1 - \theta)^{9}\, , \hspace{0.3cm} 0 < \theta < 1.$$
\begin{align*}
\pi(\theta / x) & = C_2\, \pi(\theta)\, L(\theta)\\
&  = C_3 \, \theta^9\, (1 - \theta)^{9}\, \theta^x\, (1 - \theta)^{n - x}\\
& = C_3\, \theta^{x  + 9}\, (1 - \theta)^{n + 9 - x}.
\end{align*}
$$C_3 = \frac{1}{\int_0^1\theta^{x + 9}\, (1 - \theta)^{n + 9 - x}\, d\theta} = \frac{1}{B(x + 10, n + 10 - x)}.$$

$$\pi(\theta/x)  = \frac{\theta^{x + 9} (1 - \theta)^{n + 9 - x}}{B(x + 10, n + 10 -x)}$$
$$\theta / x \, \thicksim\, BETA(10 + x, 10 + n - x)$$
\end{proof}

\textbf{Recall:}
$$B(\alpha,\beta) = \frac{\Gamma(\alpha)\, \Gamma(\beta)}{\Gamma(\alpha + \beta)} = \int^1_0x^{\alpha - 1}\, (1 - x)^{\beta  - 1}\, dx.$$


\begin{conjugate}
If a prior distribution has the property that the posterior distribution is in the same family of distributions as the prior then the prior is called a conjugate prior.
\end{conjugate}

\begin{exa}
Let $x_1, \,  X_2, \, \cdots\, , \, x_n$ be a random sample from the distribution with pdf $$f_{\theta}(x) = \theta\, x^{\theta - 1}\, , \hspace{0.3cm} 0 < x < 1, \, \hspace{0.2cm} \theta  > 1.$$
Show that the prior distribution  $\theta \thicksim GAM(a + 1\, , \, 1/b)$ is a conjugate prior.
\end{exa}

\begin{proof}[\textbf{Solution.}]
\begin{align*}
\pi(\theta) & = \frac{1}{\left(\frac{1}{b}\right)^{a + 1}\, \Gamma(a + 1)} \, \theta^a\, e^{-\theta/(\frac{1}{b})}\hspace{0.2cm}, \hspace{0.3cm} \theta > 0\\\\
& = K\, \theta^a\, e^{-b\theta}\, \, , \, \hspace{0.3cm} \theta > 0.
\end{align*}

\begin{align*}
L(\theta) &  = \prod^n_{i = 1} \theta\, x_i^{\theta} = \theta^n\, \left(\prod^n_{i = 1} x_i \right)^{\theta  - 1}\\
& = \theta^n\, \left(\prod^n_{i  = 1} x_i\right)^{\theta - 1}\, \frac{1}{\prod^n_{i = 1} x_i}\\
& = K_1\, \theta^n\, \left(\prod^n_{i = 1} x_i \right)^{\theta}.
\end{align*}


\begin{align*}
\pi(\theta /x) & = C\, \pi(\theta)\, L(\theta)\\
& = C_1\, \theta^a\, e^{-b\theta}\, \theta^n\, \left(\prod^n_{i = 1} x_i\right)^{\theta}\\
& = C_1\, \theta^{a + n}\, e^{-b\theta}\, e^{\log\left(\prod^n_{i = 1} x_i\right)^{\theta}}\\
& = C_1\, \theta^{a + n}\, e^{-b\theta}\, e^{\theta\, \log \prod^n_{i = 1} x_i}\\
& = C_1\, \theta^{a + n} \, e^{\theta\left(b - \log \prod^n_{i = 1} x_i\right)}\\
& = C_1\, \theta^{a + n}\, e^{-\theta/\frac{1}{b - \sum^n_{i = 1} \log x_i}}.
\end{align*}
$$\therefore\,\, \theta/x\, \thicksim \, GAM\left(a  + n + 1\, , \, \frac{1}{b - \sum^n_{i = 1} \log x_i}\right).$$

$\therefore\, \theta \, \thicksim\, GAM(a + 1, 1/b)$ is a conjugate prior since $\theta /x$ is also gamma.
\end{proof}



\begin{bayes}
One method of obtaining a point estimator of $\theta$ is to use the posterior distribution and a suitable loss function.
\end{bayes}





\begin{thm}
The Bayes estimator of $\theta$ for a squared error loss function with respect to the prior $\pi(\theta)$ given data $X$ is the mean of the posterior distribution given by $$\tilde{\theta} = \tilde{\theta}(X) = \int_{-\infty}^{\infty} \theta\, \pi\left(\theta/x\right)\, d\theta.$$
This estimator minimizes
$$E\left[(\tilde{\theta} - \theta)^2| X\right] = \int_{-\infty}^{\infty} \left\{\int_{-\infty}^{\infty} (\title{\theta} -\theta)^2\, f_{\theta}(x)\, dx\right\}\, \pi(\theta)\, d\theta$$
\end{thm}


\begin{exa}
Let  $X_1, \, X_2, \, \cdots\, , \, X_n$ be a random sample from the distribution with pdf $$f_{\theta}(x) = \theta\, x^{\theta  - 1}\,  , \, \hspace{0.2cm} 0 < x < 1\,, \, \hspace{0.2cm} \theta > 1.$$
Using the prior distribution $\theta \thicksim GAM\left(a + 1, 1/b\right)$, find the Bayes estimator of $\theta$ for the squared error loss function.
\end{exa}


\begin{proof}[\textbf{Solution.}]
From Example 3.4.4
$$\pi(\theta / x)  = C_1\, \theta^{a + n}\, e^{-\theta/\left(\frac{1}{b - \sum^n_{i = 1} \log X_i}\right)}$$
$$\theta /X \, \thicksim \, GAM\left(a+n + 1\, , \,\frac{1}{b - \sum^n_{i = 1} \log X_i}\right)$$
$$\tilde{\theta} = E\left[\theta/X\right] =\frac{a + n + 1}{b - \sum^n_{i = 1}\log X_i}.$$
\end{proof}



\begin{problem}
Let  $X_1, \, X_2, \, \cdots\, , \, X_n$ be a random sample from the $POI(\theta)$ distribution.
\begin{enumerate}
\item[(a)] Show that the prior distribution $$\theta \, \thicksim\, GAM\left(b + 1, \frac{1}{a}\right)$$ is a conjugate prior.
\item[(b)] using the conjugate prior in (a) find the Bayes estimator of $\theta$ for a squared error loss function.
\end{enumerate}
\end{problem}



\begin{problem}
Let  $X_1, \, X_2, \, \cdots\, , \, X_n$ be a random sample from the $X\thicksim N(0, \sigma^2)$ distribution where $\sigma^2$ is known. Assuming the prior distribution is $\theta \, \thicksim\, N\left(\frac{b}{a}, \frac{\sigma^2}{a}\right),$ fined the Bayes estimator using a squared error loss function.
\end{problem}



\begin{problem}
Let  $X_1, \, X_2, \, \cdots\, , \, X_n$ be a random sample from the $POI(\theta)$ distribution. Assuming the prior distribution of $\theta$ is $\theta \thicksim BETA(1/2,1)$, find the Bayes estimator of $\theta$ using the squared error loss function.
\end{problem}








\newpage
\chapter{HYPOTHESIS TESTS}
\section{}
A hypothesis test involves a separation of the parameter space $\Omega$ into two distinct regions, $\Omega_0$ and $\Omega - \Omega_0$. This partition of the parameter space corresponds to testing the null hypothesis that the parameter is in $\Omega_0$.
$$\text{i.e}\hspace{0.3cm} H_0:\, \theta \in \Omega_0.$$
The null hypothesis is usually the status quo and it is only rejected if there is reasonably strong evidence against it.\\
The alternative hypothesis determines what departures from the null hypothesis are anticipated.
$$\text{i.e}\hspace{0.4cm} H_1:\, \theta\in \Omega - \Omega_0$$
Since the true value of the parameter is unknown, we base our decision on the observed value of $X$ by dividing the sample space into the critical region or rejection region $R$ and its compliment $R^C$ called the acceptance region.\\
We declare that $H_0$ is false (in favor of the alternative) if we observe $x\in R$.


\begin{defn}
The power function of a test with a critical region $R$ is the function $$\beta(\theta) = P_{\theta}\left(X\in R\right)$$ or the probability that the null hypothesis is rejected as a function of the parameter.
\end{defn}

\textbf{Note:} In order to minimize the two types of errors, the power function $\beta(\theta)$ is expected to be small for $\theta \in \Omega_0$ but large otherwise.


\begin{defn}
A test has level of significance $\alpha$ if $\beta(\theta) \leq \alpha$ for all $\theta \in \Omega_0$.
\end{defn}


\textbf{Note:} The level of significance is simply an upper bound on the probability of a type I error. There is no assurance that the upper bound is tight, that is equality is achieved. The lowest such upper bound is often called the size of the test.

\begin{defn}
The size of a test is equal to $\,\,\sup_{\theta \in \Omega_0}\beta (\theta)$.
\end{defn}



\begin{exa}
Suppose a coin is tossed 1000 times. Let $x$ be the number of heads. Suppose we want to test whether the coin is fair.
\end{exa}


\begin{proof}[\textbf{Solution.}]
$X \, \thicksim\, BIN(1000,\theta)$
$$\Omega = \{\theta |\, 0 < \theta < 1\}$$
$H_0:\, \theta = \frac{1}{2}\,$ (simple hypothesis) or $\Omega_0  = \{\frac{1}{2}\}$\\
$H_1:\, \theta \neq \frac{1}{2}\,$ or $\, \theta \in \Omega - \Omega_0\,$ (composite hypothesis).\\
Suppose the critical region is $R = \{x|\, x \leq 469 $ or $ x \geq 531\} = \{x|\, |x - 500|> 30\}.$\\
Sample space $ = R\cup R^C\, = \{x|\, x = 0, \, 1, \, \cdots\, , \, 1000\}.$\\

If we observe $x\in R$ we reject $H_0$ in favour of $H_1$.
$$\text{Type I error}\, = \, \text{reject}\, \, H_0\,\, \text{if}\, \, H_0\, \, \text{is true}$$
\begin{align*}
P\left(\text{Typer I error}\right) & = P\left(X \in R|\, H_0\, \, \text{is true}\right)\\
& = P\left(\left|X - 500\right| > 0 \, |\, \theta = 1/2\right)\\
& \approx 2P\left(X < 469.5\, |\, \theta = 1/2\right)\\
& = 2P\left(\frac{X - n\theta}{\sqrt{n\theta (1 - \theta)}} \, \leq \, \frac{469.5 - 500}{\sqrt{250}}\right)\\
& = 2\, P\left(Z < -1.93\right)\\
& = 0.054.
\end{align*}
\begin{align*}
\beta(\theta) &  = P_{\theta}\left(X\in R\right)
 = P_{\theta}\left(|X - 530|\, > 30\right)\\
& = P_{\theta}\left(X\leq 469\right) + P_{\theta}\left(X\geq 531\right)\\
& \approx P\left(Z < \frac{469.5 - 1000\, \theta}{\sqrt{1000\, \theta (1 - \theta)}}\right)\, + \, P\left(Z > \frac{530.5 - 1000\, \theta}{\sqrt{1000\, \theta (1 - \theta)}}\right)
\end{align*}
% diagram 1


\begin{figure}[H]
\centering
\begin{tikzpicture}
\draw[-Stealth](-3.3,-1)--(-3.3,2.9)node[above]{$\beta(\theta)$};
\draw[-Stealth](-3.6,-0.5)--(4,-0.5)node[right]{$\theta$};
%\draw[smooth, domain  = 0:3, variable = \x]plot(\x, {sqrt(\x)});
\draw(3,2)parabola(0,0)
(-3,2)parabola(0,0)
(-3.3,0)node[]{-}
(-3.3,2)node[]{-}
(-3.3,0)node[left, scale  = 0.9]{0.054}
(-3.3,2)node[left, scale = 0.9]{1}
(0,-0.58)--(0,-0.42)
(0,-0.58)node[below, scale = 0.9]{0.5}
(2.8,-0.58)--(2.8,-0.42)
(2.8,-0.58)node[below,scale = 0.9]{0.6}
(-2.8,-0.58)--(-2.8,-0.42)
(-2.8,-0.58)node[below,scale = 0.9]{0.4};
\draw[dashed](-3.2,2)--(3.2,2);
\end{tikzpicture}
\end{figure}





For this example, level of significance  = 0.054.\\
Size of the test = 0.054 because $H_0:\, \theta = \frac{1}{2}\,$ is a simple hypothesis.
\end{proof}










\section{Uniformly Most Powerful Tests}
Tests are often constructed by specifying the size of the test, which determines the probability of the type I error and then minimizing the probability that the null hypothesis is accepted when it is false (type II). Equivalently, we try and maximize the power function of the test for $\theta \in \Omega - \Omega_0$.


\begin{defn}
A test with power function $\beta(\theta)$ is a uniformly most powerful test of size $\alpha$, if, for all other tests of the same size $\alpha$ having power function $\beta^*(\theta)$, we  have $\beta (\theta) \geq \beta^*(\theta)$ for all $\theta \in \Omega - \Omega_0$.
\end{defn}



\textbf{Note:}\\ The word uniformly refers to the fact that one function dominates, that is $\beta(\theta) \geq \beta^*(\theta)$ uniformly for all $\theta \in \Omega - \Omega_0$ when the alternative $\Omega - \Omega_0$ consists of a single a point $\{\theta\}$ (i.e simple alternative) then we drop one word ``Uniformly" and refer to a ``most powerful test".



\begin{neyman}
Let $X$ have probability (density) function $f_{\theta}(x), \, \theta \in \Omega$. Consider testing a simple null hypothesis $H_0:\, \theta = \theta_0$ against a simple alternative\\ $H_1:\, \theta = \theta_1$. Suppose the critical region is defined by $$R = \left\{x:\, \frac{f_{\theta_1}(x)}{f_{\theta_0}}\, > \, C\right\}$$ corresponds to a test of size $\alpha$, where $C$ is a constant. Then the test with this critical region is a most powerful test of size $\alpha$ for testing $$H_0:\, \theta  = \theta_0\,\, \text{against}\,\, H_1:\, \theta = \theta_1.$$
\end{neyman}


\begin{proof}
Consider another critical region $R_1$ with the same size. Then
$$P_{\theta_0}\left(X\in R\right) = P_{\theta_0}\left(X\in R_1\right) = \alpha$$

$$\int_{R}f_{\theta_0}(x)\, dx = \int_{R_1}f_{\theta_0}(x)\, dx$$

$$\int_{R\cap R^C_1}f_{\theta_0}(x)\, dx + \int_{R\cap R_1}f_{\theta_0}(x)\, dx = \int_{R_1\cap R^C}f_{\theta_0}(x)\, dx + \int_{R_1\cap R}f_{\theta_0}(x)\, dx$$

$$\implies\hspace{0.3cm} \int_{R\cap R_1^C}f_{\theta_0}(x)\, dx = \int_{R_1 \cap R^C} f_{\theta_0}(x)\, dx\hspace{0.5cm}\cdots\cdots\hspace{0.3cm}(1)$$
for $x \in R\cap R^C_1$ $$\frac{f_{\theta_1}(x)}{f_{\theta_0}(x)}\, > \, C$$

$$\implies\hspace{0.3cm} f_{\theta_1}(x) \, >\, C\, f_{\theta_0}(x)$$

$$\implies\hspace{0.3cm} \int_{R\cap R_1^C}f_{\theta_1}(x)\, dx \, > \, C\int_{R\cap R^C_1}f_{\theta_0}(x)\, dx\hspace{0.5cm}\cdots\cdots\hspace{0.3cm}(2)$$

for $\, x\in R_1\cap R^C$ $$\frac{f_{\theta_1}(x)}{f_{\theta_0}(x)}\, \leq \, C$$
$$\implies\hspace{0.3cm} f_{\theta_1}(x)\, \leq \, C\, f_{\theta_0}(x)$$
$$-f_{\theta_1}(x)\, \geq\, -C \, f_{\theta_0}(x)$$

$$\implies\hspace{0.3cm} - \int_{R_1\cap R^C}f_{\theta_1}(x)\, dx  \geq \, -C\, \int_{R_1\cap R^C}f_{\theta_0}(x)\, dx \hspace{0.5cm}\cdots\cdots\hspace{0.3cm} (3)$$

Using (1), (2) and (3) we have
\begin{align*}
\beta(\theta_1) - \beta_1(\theta_1) & = \int_{R\cap R^C_1} f_{\theta_1}(x)\, dx - \int_{R_1\cap R^C}f_{\theta_1}(x)\, dx\\
& \geq \,C\, \int_{R\cap R_1^C}f_{\theta_0}(x)\, dx - C\int_{R_1\cap R^C}f_{\theta_0}(x)\, dx\\
& = C\left[ \int_{R\cap R_1^C}f_{\theta_0}(x)\, dx - \int_{R_1\cap R^C}f_{\theta_0}(x)\, dx\right]\\
& = 0
\end{align*}
$$\implies\hspace{0.3cm} \beta(\theta_1) - \beta_1(\theta_1)\geq 0$$
Therefore the test with critical region $R$ is the most powerful.
\end{proof}


\begin{exa}
Suppose $(X_1, \, \cdots\, , \, X_n)$ are independent $N(\theta,1)$ random variables. Consider only the parameter Space $\Omega = [0,\infty)$. Suppose we want to test $H_0:\, \theta = 0\,$ against $\, H_1:\, \theta >0$.
\begin{enumerate}
\item[(a)] Choose an arbitrary $\theta_1 > 0$ and obtain the most powerful test of size $\alpha = 0.05$ of $H_0:\, \theta = \theta_0$ against $H_1:\, \theta = \theta_1$.
\item[(b)] Does this depend on the value of $\theta_1$ you choose? Can we conclude that it is uniformly most powerful?
\item[(c)] Sketch the power function of the test.
\end{enumerate}
\end{exa}



\begin{proof}[\textbf{Solution.}]
\begin{enumerate}
\item[(a).] $H_0:\, \theta   = \theta_0 = 0\,$ against $\, H_1:\, \theta = \theta_1\,$ where $\, \theta_1  > 0$.\\

By the Neyman-Pearson Lemma the most powerful test of $H_0:\, \theta  = \theta_0\,$ against $\, H_1:\, \theta = \theta_1\,$ has critical region 
$$R = \left\{x:\, \frac{f_{\theta_1}(x)}{f_{\theta_0}(x)}\, > \, C\right\}.$$
Now
\begin{align*}
f_{\theta}(x) &  = \prod^n_{i = 1} f_{\theta}(x_i)
 =  \prod^n_{i = 1}\frac{1}{\sqrt{2\pi}}\,\exp \left\{-\frac{1}{2}(x_i - \theta)^2\right\}\\
& = \left(2\pi\right)^{-n/2}\, \exp\left\{-\frac{1}{2}\sum^n_{i = 1} (x_i - \theta)^2\right\}\\
& = \left(2\pi\right)^{-n/2}\, \exp\left\{-\frac{1}{2}\sum^n_{i = 1}(x_i^2 - 2\theta x_i + \theta^2)\right\} \\
& = \left(2\pi\right)^{-n/2}\, \exp\left\{-\frac{1}{2}\sum^n_{i = 1} x_i^2 + n\theta\overline{X} - \frac{n\theta^2}{2}\right\}.
\end{align*}
\begin{align*}
\frac{f_{\theta_1}(x)}{f_{\theta_0}(x)} & = \frac{\left(2\pi\right)^{-n/2}\, \exp\left\{-\frac{1}{2}\sum^n_{i = 1}x_i^2 + n\theta_1\overline{x} - \frac{n}{2}\theta_1^2\right\}}{\left(2\pi\right)^{-n/2}\, \exp\left\{-\frac{1}{2}\sum^n_{i = 1}x_i^2 + n\theta_0\overline{x} - \frac{n}{2}\theta_0^2\right\}}\\\\
&  = e^{n\theta_1\overline{x} - \frac{n}{2}\theta_1^2}\, , \hspace{0.3cm} \theta_0 = 0.
\end{align*}
$$\implies\hspace{0.5cm} e^{n\theta_1\overline{X}  - \frac{n}{2}\theta^2_1}\, > C$$

$$\implies\hspace{0.5cm} e^{n\theta_1\overline{X}} \, > \, C\, e^{\frac{n}{2}\theta_1^2}$$

$$\implies\hspace{0.5cm} n\theta_1 \overline{X} \, > \, \log \left[C\, e^{\frac{n}{2}\, \theta_1^2}\right]$$

$$\implies\hspace{0.5cm} \overline{X} \, > \, \frac{1}{n\theta_1}\, \log\left[C\, e^{\frac{n}{2}\theta_1^2}\right]\hspace{0.3cm}\text{since}\hspace{0.2cm} \theta_1 > 0$$
$$\implies\hspace{0.3cm} \overline{X}\, > \, K$$
Determine $K$ for a size of $\alpha = 0.05$ test 
\begin{align*}
0.05 & = P_{\theta_0}\left(\overline{X}\, > \, K\right)\\
& = P\left(\frac{\overline{X} - \theta_0}{\frac{1}{\sqrt{n}}}\, > \, \frac{K -\overbrace{\theta_0}^{0}}{\frac{1}{\sqrt{n}}}\right)\\
& = P\left(Z\, > \, K\sqrt{n}\right)
\end{align*}
% DIAGRAM 

\begin{figure}[H]
\centering
\begin{tikzpicture}
\draw(-3.5,0)--(3.5,0)
(1.5,0)--(1.5,0.75)
(1.5,0)node[below, scale = 0.9]{1.645}
(1.7,0.4)--(2.3,1)
(2.3,1)node[right,scale = 0.9]{0.05};
\draw[dashed](0,0)--(0,2.39);
\path[pattern = north west lines, smooth, domain = 1.5:3.5, variable = \x]plot(\x, {6*1/sqrt(2*pi)*exp(-1/2*\x*\x)})--(3.5,0)--(1.5,0)--(1.5,0.75)--cycle;
\draw[smooth, domain = -3.5:3.5, variable = \x]plot(\x,{6*1/sqrt(2*pi)*exp(-1/2*\x*\x)});
\end{tikzpicture}
\end{figure}
\vspace*{1pt}
$$\hspace{0.5cm} K\sqrt{n} = 1.645\, \implies\,  K =  \frac{1.645}{\sqrt{n}}$$
Therefore the most powerful test has critical region $$R = \left\{x:\, \overline{X}\, > \, \frac{1.645}{\sqrt{n}}\right\}.$$


\item[(b).] The construction of this test used the fact that $\theta_1 > 0$ but it did not depend on $\theta_1$ (a particular value of $\theta_1$). Therefore we would obtain the same critical region $$R  = \left\{x:\, \overline{x} \, > \, \frac{1.645}{\sqrt{n}}\right\}$$ for each $\theta_1 > 0$.\\
$\therefore\,$ the test with critical region $R$ is the UMP test of $H_0:\, \theta  = 0\,$ against $\, H_1:\, \theta > 0\, $ (composite hypothesis).


\item[(c).]
$\begin{aligned}[t]
\beta(\theta) & = P_{\theta}\left(X\in R\right)
 = P_{\theta}\left(\overline{X} > K\right)
 = P\left(\frac{\overline{X}  - \theta}{\frac{1}{\sqrt{n}}}\, > \, \frac{\frac{1.645}{\sqrt{n}}  - \theta}{\frac{1}{\sqrt{n}}}\right)
 = P\left(Z > 1.645 - \theta \sqrt{n}\right)
\end{aligned}$


%DIAGRAM


\begin{figure}[H]
\centering
\begin{tikzpicture}[=>stealth]
\draw[-Stealth](-0.5,-0.5)--(5.5,-0.5)node[right]{$\theta$};
\draw[-Stealth](0,-1)--(0,3.5)node[above]{$\beta(\theta)$};
\draw(0,0)..controls(2,0.42)and(3,2.64)..(5,3);
\draw(0,0)node[]{-}
(0,3)node[]{-}
(0,0)node[left]{0.05}
(0,3)node[left]{1};
\draw[dashed,thick](0,3)--(5.5,3);
\end{tikzpicture}
\end{figure}
\end{enumerate}
\end{proof}





\begin{problem}
Let $X_1, \, \cdots\,, \, X_n$ be a random sample from the $GAM(1/2,\theta)$ distribution. Find the 
\begin{enumerate}
\item[(a)] Most powerful test of size $\alpha = 0.05$ of $H_0:\, \theta = 2$ against the alternative $H_1:\, \theta = \theta_1$ where $\theta_1 < 2$.
\item[(b)] Uniformly most powerful test of size $\alpha = 0.05$ of $H_0:\, \theta = 2$ against the alternative $H_1:\, \theta < 2$ and sketch the power function of this test.
\item[(c)] Uniformly most powerful test of size $\alpha = 0.05$ of $H_0:\, \theta = 2$ against the alternative $H_1:\, \theta > 2$. On the same graph as (b) sketch the power function of this test.
\end{enumerate}
\end{problem}



\begin{problem}
Let $X_1, \, \cdots\,, \, X_n$ be a random sample from the distribution with pdf $$f_{\theta}(x) = \frac{\theta}{x^{1 + \theta}}\, , \, \, x> 1, \, \, \theta > 0$$ Find the
\begin{enumerate}
\item[(a)] most powerful tests of size $\alpha = 0.05$ of $H_0:\, \theta = 1$ against $H_1:\, \theta = \theta_1$ were $\theta_1  > 1$.
\item[(b)] uniformly most powerful test of size $\alpha = 0.05$ of $H_0:\, \theta = 1$ and sketch the power function.
\item[(c)] uniformly most powerful test of size $\alpha = 0.05$ of $H_0:\, \theta = 1$ against $H_1:\, \theta < 1$. On the same graph as in (b), sketch the power function of this test.
\end{enumerate}
\end{problem}

\begin{hypothesis}
There is a close relationship between hypothesis tests and  confidence interval. For example, let $X_1, \, \cdots\,, \, X_n$ be a random sample from the $N(\theta, 1)$ distribution and we want to test hypothesis $H_0:\, \theta = \theta_0$ against $H_1:\, \theta \neq \theta_0$. The critical region
$$\left\{x:\, \left|\overline{x} - \theta_0\right|\, > \frac{1.96}{\sqrt{n}}\right\}$$
 is of size $\alpha = 0.05$ which has a corresponding acceptance region
 $$\left\{x:\, \left|\overline{x} - \theta_0\right|\, \geq \, \frac{1.96}{\sqrt{n}}\right\}.$$
\end{hypothesis}


Note that the hypothesis $H_0:\, \theta = \theta_0$ would not be rejected if $\|\overline{x} - \theta_0| \leq \frac{1.96}{\sqrt{n}}$ or equivalently 
$$\overline{x} - \frac{1.96}{\sqrt{n}} < \theta_0 < \overline{x} + \frac{1.96}{\sqrt{n}}$$ which is a $95\%$ confidence interval for $\theta$.


\begin{problem}
Let $X_1, \, \cdots\,, \, X_n$ be a random sample from the $GAM(2, \theta)$ distributions. Show that
$$R = \left\{x:\, \sum^5_{i = 1}x_i  < 4.7955\, \theta_0\, \, \text{or}\,\, \sum^5_{i = 1} x_i > 17.085\, \theta_0\right\}$$
is a size $\alpha = 0.05$ critical region for testing $H_0:\, \theta = \theta_0$. Show that this critical region may be used to construct a $95\%$ C.I for $\theta$. 
\end{problem}






\section{Locally Most Powerful tests}
It is not always possible to construct uniformly most powerful tests. In most problems, alternative parameter values that are close to the null hypothesis are hard to detect.  It is therefore important to use a test that has high a power as possible for alternative parameter values that are close to the null hypothesis. Such a test is called locally most powerful.



\begin{defn}
A test of $H_0:\, \theta = \theta_0$ against $H_1:\, \theta > \theta_0$ with power function $\beta(\theta)$ is locally most powerful if, for any other test having the same size and power function $\beta^*(\theta)$, there exists an $\varepsilon > 0$ such that $\beta(\theta)\geq \beta^*(\theta)$ for all $\theta_0 < \theta < \theta_0 + \varepsilon$.
\end{defn}



\begin{thm}
Suppose $X_1, \, \cdots\,, \, X_n$ is a random sample from a regular statistical model $\{f_{\theta}(x);\, \theta \in \Omega\}$. A test with critical region $$R = \left\{x:\, S(\theta_0, x) > C\right\}$$ where $C$ is a constant determined by the size of the test, is a locally most powerful test of $H_0:\, \theta = \theta_0$ against $H_1:\, \theta > \theta_0$.
\end{thm}



\begin{exa}
Suppose $X_1, \, \cdots\,, \, X_n$ is a random sample from a $N(\theta ,1)$ distribution. Show that the locally most powerful test of $H_0:\, \theta = 0$ against $H_1:\, \theta > 0$ is also the uniformly most powerful test.
\end{exa}

%
\begin{proof}[\textbf{Solution.}]
$H_0:\, \theta = 0\,$ against $\, H_1:\, \theta > 0$.\\
The rejection region for the locally most powerful test is $R = \{x:\, S(\theta_0,x) > C\}$.
$$L(\theta) = \prod^n_{i = 1} \frac{1}{\sqrt{2\pi}}\, e^{-\frac{1}{2}(x_i - \theta)^2}   = \left(2\pi\right)^{-n/2}\, e^{-\frac{1}{2}\sum^n_{i = 1}(x_i - \theta)^2}.$$

$$\mathcal{L}(\theta)   = -\frac{1}{2}\sum^n_{i = 1} (x_i - \theta)^2 - \frac{n}{2}\, \log 2\pi .$$
\begin{align*}
S(\theta) & = (-1)\, \sum^n_{i = 1}(x_i - \theta)(-)
 =\sum^n_{i = 1}x_i - n\theta
 = n\overline{x} - n\theta
 = n(\overline{x} - \theta).
\end{align*}
$$S(\theta_0,x)\, > \, C$$
$$n\left(\overline{x} - \theta_0\right)\, > \, C$$
$$n\overline{x}\, > \, C$$
$$\overline{x}\, > \, \frac{C}{n}$$
$$\overline{x}\, > \, K$$
where $K$ is determined by the size of the test.\\
Therefore the critical region corresponding to the LMP test is $R = \{x:\,\overline{x} > K\}$ which is the same as that for the UMP test.
\end{proof}




\begin{problem}
Consider a single observation $X$ from the LOG$(1,\theta)$ distribution. Find the locally most powerful test of $H_0:\, \theta = 0$ against $H_1:\, \theta > 0$.
\end{problem}
%----------------------------------------------

\section{Likelihood ratio tests}

The likelihood ratio test for testing the composite hypothesis $H_0: \theta \in \Omega_0$ against $H_1: \theta \in \Omega - \Omega_0$ has critical region $R = \left\{ x: \, \Lambda (x) =  > C\right\}$ where 
$$\Lambda(x) = \frac{\sup_{\theta\in \Omega} f_{\theta}(x)}{\sup_{\theta\in\Omega_0} f_{\theta}(x)} = \frac{\sup_{\theta\in\Omega} L(\theta)}{\sup_{\theta\in \Omega_0} L(\theta)}$$ and $C$ is determined by the size of the test.\\

\textbf{Note:} that $\Lambda(x) = \frac{L\left(\widehat{\theta}\right)}{L\left(\widehat{\theta}_R\right)}$ where $\widehat{\theta}$ is the usual $ML$ estimate and $\widehat{\theta}_R$ is the $ML$ estimate calculated under the restriction $H_0: \theta\in\Omega_0$.


\begin{exa}
Let $X_1, \, \cdots\, , X_n$ be a random sample from then $N(\mu, \sigma^2)$ distribution where $\mu$ and $\sigma^2$ are unknown. Consider a test of $H_0: \mu = 0, \, 0 < \sigma^2 < \infty\,$ against $H_1:\, \mu \neq 0, \, 0  < \sigma^2 < \infty$.\\
Show that the likelihood ration test of $H_0$ against $H_1$ has critical region $$ R = \left\{ x:\, \frac{n\overline{x}^2}{s^2}\, > C\right\}.$$
Show that under $H_0$ the statistic $\, \frac{n\, \overline{X}^2}{S^2}\,$ has an $F(1, n- 1)$ distribution and thus find a size $\alpha = 0.05$ test for $n = 20$.
\end{exa}

\begin{proof}[\textbf{\textit{Solution.}}]
$H_0: \mu = 0, \, 0 < \sigma^2 < \infty\hspace{0.3cm}$ against $\hspace{0.3cm} H_1: \mu \neq 0, \, 0 < \sigma^2 < \infty\,, \, \, \theta = (\mu ,\sigma^2)$
$\Omega = \left\{(\mu, \sigma^2):\, -\infty < \mu < \infty, \, \sigma^2 > 0\right\}\hspace{0.3cm}$ and $\hspace{0.3cm}\Omega_0 = \left\{(0,\sigma^2):\, \sigma^2>0\right\}$
$$\Lambda(x) = \frac{\left(\widehat{\theta}\right)}{\left(\widehat{\theta}_R\right)}$$

$$L(\theta) = \prod^n_{i = 1} \frac{1}{\sqrt{2\pi\sigma^2}} e^{-\frac{1}{2\sigma^2}(x_i - \theta)^2} = \left(2\pi \sigma^2\right)^{-\frac{n}{2}}e^{-\frac{1}{2\sigma^2}\sum^n_{i = 1} (x_i - \mu)^2}.$$
From example 2.2.4
$$\widehat{\mu} = \overline{x}\,, \hspace{0.3cm} \widehat{\sigma}^2 = \frac{1}{n}\sum^n_{i = 1} \left(X_i - \overline{X}\right)^2$$
replacing we obtain
\begin{align*}
L\left(\widehat{\theta}\right) & = \left(2\pi\widehat{\sigma}^2\right)^{-\frac{n}{2}}\, e^{-\frac{1}{2\widehat{\sigma}^2}\sum^n_{i = 1} \left(x_i - \widehat{\mu}\right)^2}\\
& = \left(2\pi \widehat{\sigma}^2\right)^{-\frac{n}{2}}\,e^{-\frac{1}{2\widehat{\sigma}^2} \sum^n_{i  = 1} (x_i - \overline{x})^2}\\
& =  \left(2\pi \widehat{\sigma}^2\right)^{-\frac{n}{2}}\,e^{-\frac{1}{2\widehat{\sigma}^2}\left(n\widehat{\sigma}^2\right)}\\
& =  \left(2\pi \widehat{\sigma}^2\right)^{-\frac{n}{2}}\, e^{-\frac{n}{2}}
\end{align*}


Under $H_0:\, \mu = 0\, , \hspace{0.3cm} \sigma^2 > 0$
$$L(\theta) = \left(2\pi\sigma^2\right)^{-\frac{n}{2}}\, e^{-\frac{1}{2\sigma^2}\sum^n_{i = 1} x_i^2}.$$

$$\mathcal{L}(\theta) = -\frac{1}{2\sigma^2} \sum^n_{i =1} x_i^2 - \frac{n}{2}\log \left(2\pi \sigma^2\right).$$

$$S(\theta) = \frac{\partial\mathcal{L}}{\partial \sigma^2} = \frac{1}{2(\sigma^2)^2}\sum^n_{i = 1} x_i^2 - \frac{n}{2}\cdot\frac{2\pi}{2\pi \sigma^2}= \frac{1}{2(\sigma^2)^2}\sum^n_{i = 1} x_i^2 - \frac{n}{2\sigma^2}.$$
$S(\theta) = 0$
$$\frac{n}{2\sigma^2} = \frac{1}{2(\sigma^2)^2}\sum^n_{ i = 1}x_i^2 \, \implies\, \widehat{\sigma}_R^2 = \frac{1}{n}\sum^n_{i = 1} x_i^2.$$

$$L\left(\widehat{\theta}_R\right) = \left(2\pi \widehat{\sigma}^2_R\right)^{-\frac{n}{2}}\, e^{-\frac{1}{2\widehat{\sigma}^2_R}\sum^n_{i = 1} x_i^2} =  \left(2\pi \widehat{\sigma}^2_R\right)^{-\frac{n}{2}}\, e^{-\frac{1}{2\widehat{\sigma}^2_R}\cdot n\widehat{\sigma}^2_R} = \left(2\pi \widehat{\sigma}^2_R\right)^{-\frac{n}{2}}\, e^{-\frac{n}{2}} $$

\begin{align*}
\frac{L\left(\widehat{\theta}\right)}{L\left(\widehat{\theta}_R\right)} & = \frac{\left(2\pi \widehat{\sigma}^2\right)^{-\frac{n}{2}}\, e^{-\frac{n}{2}}}{\left(2\pi \widehat{\sigma}^2_R\right)^{-\frac{n}{2}}\, e^{-\frac{n}{2}}} = \frac{\left(\widehat{\sigma}_R^2\right)^{\frac{n}{2}}}{\left(\widehat{\sigma}^2\right)^{\frac{n}{2}}} = \left(\frac{\widehat{\sigma}^2_R}{\widehat{\sigma}^2}\right)^{\frac{n}{2}}\\\\
& = \left(\frac{\frac{1}{n}\sum^n_{i = 1} x_i^2}{\frac{1}{n}\sum^n_{i = 1} (x_i - \overline{x})^2}\right)^{\frac{n}{2}}\\\\
& = \left(\frac{\sum^n_{i = 1} \left((x_i - \overline{x})^2 + \overline{x}(x_i - \overline{x}) + \overline{x}^2\right)}{\sum^n_{i = 1} (x_i - \overline{x})^2}\right)^{\frac{n}{2}}
\end{align*}
\begin{align*}
&  = \left(\frac{\sum^n_{i = 1} (x_i - \overline{x})^2 + n\overline{x}^2}{\sum^n_{i = 1} (x_i - \overline{x})^2}\right)^{\frac{n}{2}}\\\\
& =\left(\frac{(n-1)S^2 + n\overline{x}^2}{(n-1) S^2}\right)^{\frac{n}{2}}   = \left(1 + \frac{n\overline{x}^2}{(n -1)S^2}\right)^{\frac{n}{2}}\\\\
& = \left(1 + \frac{1}{n - 1}\cdot \frac{n\overline{x}^2}{S^2}\right)^{\frac{n}{2}}
\end{align*}

$\left(1 + \frac{1}{n - 1}\cdot\frac{n\overline{x}^2}{S^2}\right)^{\frac{n}{2}} > K\,$ is equivalent to $\frac{n\overline{x}^2}{S^2} > C$ where $C$ is determined by the size of the test
$$\therefore\,\, R = \left\{ x:\, \frac{n\overline{x}^2}{S^2} > C\right\}.$$

Under $H_0:\, \mu = 0$
$$T = \frac{n\, \overline{X}^2}{S^2} = \frac{\left(\dfrac{\overline{X} - 0}{\sigma/\sqrt{n}}\right)^2}{\dfrac{(n - 1) S^2/\sigma^2}{n - 1}}$$

$$\left(\frac{\overline{X} - 0}{\frac{\sigma}{\sqrt{n}}}\right)^2 = Z^2 \, \thicksim\, \chi^2_{(1)}\hspace{0.5cm}\text{and}\hspace{0.5cm}\frac{(n - 1)S^2}{\sigma^2} \thicksim \chi^2_{n - 1}$$
$$T = \frac{\dfrac{\chi^2_{(1)}}{1}}{\dfrac{\chi^2_{n - 1}}{n - 1}}\, \thicksim\, F(1, n-1).$$

$\alpha = P_{H_0}(X\in R)$
$$0.05 = P(T>C) = P(F_{1,19} > C)$$
Therefore $\hspace{0.2cm} C = 4.35$.
\end{proof}




\begin{problem}
Suppose $X\thicksim GAM(2,\beta_1)$ and $Y\thicksim GAM(2,\beta_2)$ independently. Show that the likelihood ratio statistic for testing the hypothesis $H_0: \beta_1 = \beta_2\,$ against the alternative $\, H_1:\beta_1\neq \beta_2$ is a function of the statistic 
$$T = \frac{X}{X + Y}.$$
Find the distribution of $T$ under $H_0$. Find the critical region for a size $\alpha = 0.01$ test.\\

\textit{Hint:} If $X\thicksim GAM(a,1)$ and $Y\thicksim GAM(b,1)$ are independent random variable then $$\frac{X}{X + Y} \thicksim BETA(a,b).$$
\end{problem}



\begin{problem}
Suppose $X_1, \, \cdots\, , X_n$ are independent $EXP(\lambda)$ random variables and independently $Y_1, \, \cdots\, , \, Y_n$  are independent $EXP(\mu)$ random variables.\\
Show that the likelihood ratio statistic for testing the hypothesis $H_0: \lambda = \mu$ against the alternative $H_1: \lambda \neq \mu$ is  a function of 
$$T = \frac{\sum^n_{i = 1} X_i}{\sum^n_{i =1} X_i + \sum^n_{i = 1} Y_i}.$$
\end{problem}




\begin{problem}
Suppose $X_1, \, \cdots\, , X_n$  is a random sample from a regular statistical model $\{f_{\theta}(x); \theta \in \Omega\}$ with $\Omega$ being an open set in a $k$-dimensional Euclidean space. Consider a subset of $\Omega$ defined by $$\Omega_0 = \{\theta(\eta):\, \eta \in \, \text{ open subset of  a}\, q-\text{dimensional Euclidean space}\}.$$
Then the likelihood ratio statistic defined by 
$$\Lambda_n(x) = \frac{\sup_{\theta\in \Omega} \prod^n_{i = 1} f_{\theta}(x_i)}{\sup_{\theta\in \Omega_0} \prod^n_{i =1} f_{\theta}(x_i)} = \frac{\sup_{\theta\in \Omega}L(\theta)}{\sup_{\theta\in \Omega_0} L(\theta)}$$
is such that, under $H_0: \theta \in \Omega_0$,
$$2\, \log \Lambda_n(x) \, \underset{D}{\longrightarrow} \, W\thicksim \chi^2_{(k - q)}.$$
\end{problem}


\textbf{Note:}\\
The number of degrees of freedom is the difference between the number of parameters that need to be estimated in the general model and the number of parameters to be estimated under the restrictions imposed by $H_0$.




\begin{exa}
Suppose $X_1, \, \cdots\, , X_n$  are independent $POI(\lambda)$ random variables and independently $Y_1, \, \cdots\, , \, Y_n$  are independent $POI(\mu)$ random variables. Find the likelihood ratio test of $H_0:\lambda = \mu$ against $H_1: \lambda \neq \mu$.
\end{exa}

\begin{proof}[\textbf{\textit{Solution.}}]
$H_0: \lambda = \mu\hspace{0.3cm}$ against $\hspace{0.3cm} H_1: \lambda \neq \mu$.
$$\Omega  = \{(\lambda , \mu):\, \lambda, \, \mu  > 0\} \hspace{0.5cm}\text{and}\hspace{0.5cm} \Omega_0 = \{(\lambda,\lambda):\, \lambda > 0\}.$$

$$L(\theta) = \prod^n_{i = 1} f_{\lambda}(x_i)\, \prod^n_{i = 1} f_{\mu}(y_i)= \prod^n_{i = 1} \frac{e^{-\lambda}\,\lambda^{x_i}}{x_i!}\,\prod^n_{i = 1} \frac{e^{-\mu}\,\mu^{y_i}}{y_i!}  = \frac{e^{-n\lambda}\, \lambda^{\sum^n_{i=1} x_i}}{\prod^n_{i = 1} x_i!}\,\frac{e^{-n\mu}\, \mu^{\sum^n_{i=1} y_i}}{\prod^n_{i = 1} y_i!}.$$

$$\mathcal{L}(\theta) = -n\lambda + \sum^n_{i = 1} x_i \, \log \lambda - \log \prod^n_{i = 1} x_i! - n\mu + \sum^n_{i = 1} y_i \log \mu - \log \prod^n_{i = 1} y_i!.$$

$$S(\theta) = \begin{pmatrix}
 -n  + \frac{\sum^n_{i =1} x_i}{\lambda}\\\\ -n + \frac{\sum^n_{i = 1}y_i}{\mu}\\
\end{pmatrix}$$
$S(\theta) = 0$
$$\implies \hspace{0.3cm} n = \frac{\sum^n_{i =1} x_i}{\lambda}\hspace{0.5cm} , \hspace{0.5cm} n = \frac{\sum^n_{i = 1} y_i}{\mu}$$
$$\widehat{\lambda} = \frac{1}{n}\sum^n_{i = 1} x_i = \overline{x}\hspace{0.5cm} , \hspace{0.5cm} \widehat{\mu} = \frac{1}{n}\sum^n_{i = 1} y_i = \overline{y}.$$

\begin{align*}
L\left(\widehat{\theta}\right)  & = \frac{e^{-n \widehat{\lambda}}\, \widehat{\lambda}^{\sum^n_{i = 1}x_i}}{\prod^n_{i = 1} x_i!}\,\cdot\, \frac{e^{-n \widehat{\mu}}\, \widehat{\mu}^{\sum^n_{i = 1}y_i}}{\prod^n_{i = 1} y_i!} =\frac{e^{-n \overline{x}}\, \overline{x}^{\sum^n_{i = 1}x_i}}{\prod^n_{i = 1} x_i!}\,\cdot\, \frac{e^{-n \overline{y}}\, \overline{y}^{\sum^n_{i = 1}y_i}}{\prod^n_{i = 1} y_i!}\\\\
&  = \frac{e^{-n(\overline{x} + \overline{y})}\, \overline{x}^{n\overline{x}}\, \overline{y}^{n\overline{y}}}{\prod^n_{i = 1} x_i!\ \prod^n_{i = 1} y_i!}\\\\
& =  \frac{e^{-n(\overline{x} + \overline{y})}\, \left(\overline{x}^{\overline{x}}\right)^n\, \left(\overline{y}^{\overline{y}}\right)^n}{\prod^n_{i = 1} x_i!\ \prod^n_{i = 1} y_i!}.
\end{align*}

Under $H_0$
$$L(\theta) = \frac{e^{-n\lambda}\, \lambda^{\sum^n_{i = 1} x_i}}{\prod^n_{i = 1}x_i!}\, \cdot\, \frac{e^{-n\lambda}\, \lambda^{\sum^n_{i = 1} y_i}}{\prod^n_{i = 1}y_i!} = \frac{e^{-2n\lambda}\, \lambda^{\sum^n_{i = 1} x_i + \sum^n_{i = 1} y_i}}{\prod^n_{i = 1} x_i!\, \prod^n_{i = 1} y_i!}.$$

$$\mathcal{L}(\theta) = -2n\lambda + \left(\sum^n_{i = 1} x_i + \sum^n_{i = 1} y_i \right) \log \lambda - \log \prod^n_{i = 1} x_i!  - \log \prod^n_{i = 1} y_i!.$$

$$S(\theta)  = \frac{\partial \mathcal{L}}{\partial\lambda} = -2n + \frac{\sum^n_{i = 1} x_i + \sum^n_{i = 1} y_i}{\lambda}.$$
$S(\theta) = 0$
$$2n\lambda = \sum^n_{i = 1}  + \sum^n_{i = 1} y_i$$

$$\widehat{\lambda}_R = \frac{\sum^n_{i = 1} x_i + \sum^n_{i = 1} y_i}{2n} = \frac{\overline{x} + \overline{y}}{2}.$$
\begin{align*}
L\left(\widehat{\theta}_R\right) & = \frac{e^{-2n\widehat{\lambda}_R}\, \widehat{\lambda}_R^{\sum^n_{i = 1} x_i + \sum^n_{i = 1} y_i}}{\prod^n_{i = 1} x_i!\, \prod^n_{i = 1}y_i!}\\\\
& = \frac{e^{-2n\left(\frac{\overline{x} + \overline{y}}{2}\right)}\, \, \, \left(\dfrac{\overline{x} + \overline{y}}{2}\right)^{n\overline{x} + n\overline{y}}}{\prod^n_{i = 1} x_i!\, \prod^n_{i = 1}y_i!} = \frac{e^{-n\left(\overline{x} + \overline{y}\right)}\, \, \, \left(\dfrac{\overline{x} + \overline{y}}{2}\right)^{n(\overline{x} + \overline{y})}}{\prod^n_{i = 1} x_i!\, \prod^n_{i = 1}y_i!}
\end{align*}

\begin{align*}
\Lambda(x,y) & = \frac{L\left(\widehat{\theta}\right)}{L\left(\widehat{\theta}_R\right)} = \frac{\dfrac{e^{-n(\overline{x} + \overline{y})}\, \left(\overline{x}^{\overline{x}}\right)^n \, \left(\overline{y}^{\overline{y}}\right)^n}{\prod^n_{i = 1} x_i!\, \prod^n_{i = 1} y_i!}}{\dfrac{e^{-n\left(\overline{x} + \overline{y}\right)}\, \, \, \left(\dfrac{\overline{x} + \overline{y}}{2}\right)^{n(\overline{x} + \overline{y})}}{\prod^n_{i = 1} x_i!\, \prod^n_{i = 1}y_i!}}
=\frac{\left(\overline{x}^{\overline{x}}\right)^n\, \, \left(\overline{y}^{\overline{y}}\right)^n}{\left[\left(\dfrac{\overline{x} + \overline{y}}{2}\right)^{\overline{x} + \overline{y}}\right]^n}
 = \left(\frac{\overline{x}^{\overline{x}}\, \, \overline{y}^{\overline{y}}}{\left(\dfrac{\overline{x} + \overline{y}}{2}\right)^{\overline{x} + \overline{y}}}\right)^n.
\end{align*}




$$\therefore\,\,\, R = \left\{(x,y):\, \frac{\overline{x}^{\overline{x}}\, \overline{y}^{\overline{y}}}{\left(\dfrac{\overline{x} + \overline{y}}{2}\right)^{\overline{x}  + \overline{y}}}\, > C\right\}$$ where $C$ is determined by the size of the test.\\

Under $H_0$, 
$$2\, \log \Lambda(X, Y) \, \underset{D}{\longrightarrow}\, W\, \thicksim\, \chi^2_{(2-1)} = \chi^2_{(1)}.$$
\end{proof}








%-----------------------------------------------
\section{Score and Wald Tests}

\begin{score test}
Score tests can be viewed as a more general class of tests of $H_0:\, \theta = \theta_0$ against $H_1:\, \theta\in \Omega - \{\theta_0\}$.\\
If the usual regularity conditions hold then under $H_0$ we have $$S(\theta_0, X)\, \left[J(\theta_0)\right]^{-\frac{1}{2}}\, \underset{D}{\longrightarrow}\, Z\thicksim(0,1)$$ and  thus 
$$R(\theta_0,X) = \left[S(\theta_0,X)\right]^2\, \left[J(\theta_0)\right]^{-1}\, \underset{D}{\longrightarrow}\, W\thicksim \chi^2_{(1)}.$$ For a vector $\theta = (\theta_1, \, \cdots\, , \theta_k)^t$ we have 
$$R(\theta_0, X) = \left[S(\theta_0, X)\right]^t\, \left[J(\theta_0)\right]^{-1}\, \left[S(\theta_0,X)\right]\,\underset{P}{\longrightarrow}\ W\thicksim  \chi^2_{(k)}.$$
The test based on $R(\theta_0,X)$ is called a (Rao) Score test. It has critical region\\ $R =\{x:\, R(\theta_0, x)\, > c\}$ where $c$ is determined by the size of the test i.e $\alpha = P(W>c)$ where $W\thicksim \chi^2_{(k)}$. The test based on $R(\theta_0,X)$ is asymtoticaly equivalent to the likelihood ratio test.
\end{score test}



\begin{wald test}
Suppose that $\widehat{\theta}$ is the ML estimator of $\theta$ over all $\theta\in \Omega$ and we wish to test $H_0:\, \theta = \theta_0$ against $H_1:\, \theta \in \Omega - \{\theta_0\}$. If the usual regularity conditions hold then under $H_0$
$$w(\theta_0,X) - \left(\widehat{\theta} - \theta_0\right)^t \, J(\theta)\, \left(\widehat{\theta} - \theta_0\right) \, \underset{D}{\longrightarrow}\, w\thicksim \chi^2_{(k)}.$$
A test based on the test statistic $w\left(\theta_0, X\right)$ is called a Wald test. It has critical region $$R = \{x:\, w\left(\theta_0, x\right) > c\}$$ where $c$ is determined by the size of test.
\end{wald test}

\textbf{Note:} $J(\theta_0)$ may be replaced by $I\left(\widehat{\theta}\right)$ to give an asymptotically equivalent test statistic.

\begin{exa}
Suppose $X\thicksim BIN(n,\theta)$. Find the Score test and the Wald test for testing $H_0:\, \theta  = \theta_0$ against $H_1:\, \theta\neq \theta_0$.
\end{exa}


\begin{problem}
Suppose $(X_1, \, \cdots\, , X_n)$ is a random sample from the $POI(\theta)$ distribution. Find the Score test and Wald test for testing $H_0:\, \theta = \theta_0$ against $H_1:\, \theta \neq \theta_0$.
\end{problem}

\begin{proof}[\textbf{Solution.}]
$H_0:\, \theta = \theta_0\,$ against $\, H_1:\, \theta \neq \theta_0$
$$L(\theta) = \binom{n}{x}\, \theta^x\, (1  - \theta)^{n - x}.$$

$$\mathcal{L}(\theta) = \log\binom{n}{x}  + x\, \log \theta + (n - x)\, \log(1 - \theta).$$

$$S(\theta) = \frac{x}{\theta} + (n - x)\, \frac{(1)}{1 - \theta} \,  = \, \frac{x}{\theta} - \frac{n - x}{1 - \theta}.$$

$S(\theta) = 0$
$$\frac{x}{\theta} = \frac{n - x}{1 - \theta}$$
$$x - \theta x = n\theta - \theta x$$
$$\widehat{\theta} = \frac{x}{n} = \overline{x}.$$

$$I(\theta) = -\left[\frac{-x}{\theta^2} - \frac{(n - x)\, (-1)\, (-1)}{(1 - \theta)^2}\right]\, = \, \frac{x}{\theta^2} + \frac{n - x}{(1 - \theta)^2}.$$
\begin{align*}
J(\theta) & = E\left[I(\theta,X)\right] = E\left[\frac{x}{\theta^2} + \frac{n - x}{(1 - \theta)^2}\right]\\
& = \frac{E(X)}{\theta^2} + \frac{n - E(X)}{(1 - \theta)^2}\\
& =\frac{n\theta}{\theta^2} + \frac{n - n\theta}{(1 - \theta)^2} \\
& = \frac{n}{\theta} + \frac{n}{1 -\theta}\\
& = \frac{n}{\theta(1  - \theta)}.
\end{align*}

\textbf{Score test:}
\begin{align*}
R(\theta_0,x) &  = \frac{\left(S(\theta_0,x)\right)^2}{J(\theta_0)} = \frac{\left(\frac{x}{\theta_0} - \frac{n - x}{1 - \theta_0}\right)^2}{\frac{n}{\theta_0(1 - \theta_0)}}\\
& = \left(\frac{x}{\theta_0} - \frac{n - x}{1 - \theta_0}\right)^2\, \frac{\theta_0\,(1 - \theta_0)}{n}\\
& = \left(\frac{x - n\theta_0}{\theta_0(1 - \theta_0)}\right)^2\, \frac{\theta_0(1 - \theta_0)}{n}\\
&  = \frac{(x - n\theta_0)^2}{n\theta_0(1 - \theta_0)}.
\end{align*}
with rejection region $R = \{x:\, R(\theta_0,x)>c\}$ where $c$ is determined from $\alpha = P\left(\chi^2_{(1)} > c\right)$.\\

\textbf{Wald test:}
\begin{align*}
w\left(\theta_0,x\right) & = \left(\widehat{\theta} -\theta_0\right)^t\, J(\theta_0)\, \left(\widehat{\theta} - \theta_0\right)\\
&  = \left(\frac{x}{n} - \theta_0\right)\, \frac{n}{\theta_0(1 - \theta_0)}\, \left(\frac{x}{n} - \theta_0\right)\\
& = \left(\frac{x}{n} - \theta_0\right)^2\, \frac{n}{\theta_0\, (1  - \theta_0)}\\
& = \left(\frac{x  - n\theta_0}{n}\right)^2\, \frac{n}{\theta_0\, (1 -\theta_0)}\\
& = \frac{\left(x - n\theta_0\right)^2}{n\theta_0(1 - \theta_0)}
\end{align*}
Therefore $\, R = \{x:\, w(\theta_0, x) > c \}\,$ where $c$ is determined from $\alpha = P\left(\chi^2_{(1)} > c\right)\,$\\ i.e $w(\theta_0, x) = R(\theta_0, x)$.
\end{proof}


\begin{problem}
Let $X_1, \, \cdots\, , X_n$ be a random sample from the distribution with pdf $$f_{\theta}(x) = \frac{\theta}{x^{1 + \theta}}\, , \hspace{0.3cm} x\geq 1$$
Find the Score test and the Wald test for testing $H_0:\, \theta = \theta_0$ against $H_1:\, \theta \neq \theta_0.$
\end{problem}








\newpage
\begin{thebibliography}{}
\bibitem{}{}
Casella, G. and R. L. Berger, 2002. Statistical Inference. $2^{\text{nd}}$ ed. Pacific Groove: Duxbury.
\bibitem{}{}
Garthwaite, P. H., I. T. Jolliffe and B. Jones, 2002. Statistical Inference. $2^{\text{nd}}$ ed. Oxford: Oxford University Press
\bibitem{}{}
Rao, C. R., 1995. Linear Statistical Inference and its Applications. New York: Wiley and Sons.
\bibitem{}{}
Silvey, S. D., 1988. Statistical Inference, London: Chapman and Hall.
\bibitem{}{}
Knight, K., 1999. Mathematical Statistics, New York: Chapman and Hall.
\bibitem{}{}
Bickel, P. J, and K. A. Doksum, 2001. Mathematical Statistics. $2^{\text{nd}}$ ed. Upper Saddle River: Prentice Hall.
\end{thebibliography}


\end{document}
